% Generated mechanically from frozen input p04/ko/constructions.tex.
% Do not edit this generated file; edit the bound locale source instead.

% OK, start here.
%

\setcounter{chapter}{26}
\chapter{스킴의 구성}



\phantomsection
\label{constructions-section-phantom}


\section{서론}
\label{constructions-section-introduction}

\noindent
이 장에서는 주어진 스킴들로부터 새로운 스킴을 구성하는 방법들을
도입한다. 기본 참고문헌은 \cite{EGA}이다.

\section{상대적 붙이기}
\label{constructions-section-relative-glueing}

\noindent
다음 보조정리는 스킴 $X$를 $S$ 위에 구성하려 하며, $X$를 $S$의 아핀
열린집합들에 제한하는 방법을 이미 알고 있는 경우에 쓰인다. 실제 결과는 완전히
일반적이며 (국소) 환 달린 공간의 맥락에서도 성립하지만, 증명은 스킴의 언어로
서술한다.

\begin{lemma}
\label{constructions-lemma-relative-glueing}
$S$를 스킴이라 하자.
$\mathcal{B}$를 $S$의 위상의 기저라 하자.
다음 데이터가 주어졌다고 하자.
\begin{enumerate}
\item 모든 $U \in \mathcal{B}$에 대해 스킴 $f_U : X_U \to U$, 즉
$U$ 위의 스킴이 주어진다.
\item $U, V \in \mathcal{B}$이고 $V \subset U$일 때, 사상
$\rho^U_V : X_V \to X_U$, 즉 $U$ 위의 사상이 주어진다.
\end{enumerate}
다음을 가정하자.
\begin{enumerate}
\item[(a)] 각 $\rho^U_V$가 동형 $X_V \to f_U^{-1}(V)$, 즉 $V$ 위의
스킴들의 동형을 유도하고,
\item[(b)] $W, V, U \in \mathcal{B}$이고 $W \subset V \subset U$일 때
$\rho^U_W = \rho^U_V \circ \rho ^V_W$가 성립한다.
\end{enumerate}
그러면 스킴의 사상 $f : X \to S$와 동형 $i_U : f^{-1}(U) \to X_U$가
각 $U \in \mathcal{B}$ 위에서 존재하며, $V, U \in \mathcal{B}$이고
$V \subset U$일 때 합성
$$
\xymatrix{
X_V \ar[r]^{i_V^{-1}} &
f^{-1}(V) \ar[rr]^{inclusion} & &
f^{-1}(U) \ar[r]^{i_U} &
X_U
}
$$
은 사상 $\rho^U_V$이다. 또한 $X$는 $S$ 위의 유일한 동형을 제외하고
유일하다.
\end{lemma}

\begin{proof}
이를 증명하기 위해 Schemes, Lemma \ref{schemes-lemma-glue-functors}를 사용한다.
먼저 스킴의 범주에서 집합의 범주로 가는 반변 함자 $F$를 정의한다.
즉, 스킴 $T$에 대해 다음과 같이 둔다.
$$
F(T) =
\left\{
\begin{matrix}
(g, \{h_U\}_{U \in \mathcal{B}}),
\ g : T \to S, \ h_U : g^{-1}(U) \to X_U, \\
f_U \circ h_U = g|_{g^{-1}(U)},
\ h_U|_{g^{-1}(V)} = \rho^U_V \circ h_V
\ \forall\ V, U \in \mathcal{B}, V \subset U
\end{matrix}
\right\}.
$$
제한 사상 $F(T) \to F(T')$은 사상 $T' \to T$가 주어졌을 때 단순히
합성으로 얻는다. 임의의 $W \in \mathcal{B}$에 대해, 부분함자
$F_W \subset F$를 계 $(g, \{h_U\})$들 중 $g(T) \subset W$를 만족하는
것들로 이루어지도록 정의한다.

\medskip\noindent
먼저 $F$가 자리스키 위상에 대한 층 조건을 만족함을 보인다. $T$가 스킴이고
$T = \bigcup V_i$가 열린 덮개이며, $\xi_i \in F(V_i)$가
$\xi_i|_{V_i \cap V_j} = \xi_j|_{V_i \cap V_j}$를 만족하는 원소라고 하자.
$\xi_i = (g_i, \{h_{i, U}\})$라 쓰자. 그러면 사상 $g_i$들이 유일한 전역 사상
$g : T \to S$로 붙는다는 것을 곧바로 알 수 있다. 또한
$g^{-1}(U) = \bigcup g_i^{-1}(U)$임이 분명하다. 따라서 사상
$h_{i, U} : g_i^{-1}(U) \to X_U$들은 유일한 사상
$h_U : g^{-1}(U) \to X_U$로 붙는다. 계 $(g, \{h_U\})$가 $F(T)$의
원소임은 쉽게 확인된다. 그러므로 $F$는 자리스키 위상에 대한 층 조건을 만족한다.

\medskip\noindent
다음으로 각 $F_W$, $W \in \mathcal{B}$가 표현가능임을 확인한다. 구체적으로
함자들의 변환
$$
F_W \longrightarrow \Mor(-, X_W), \ (g, \{h_U\}) \longmapsto h_W
$$
이 동형이라고 주장한다. 이를 보기 위해 $T$가 스킴이고
$\alpha : T \to X_W$가 사상이라고 하자. $g = f_W \circ \alpha$로 둔다.
임의의 $U \in \mathcal{B}$가 $U \subset W$를 만족하면,
$h_U : g^{-1}(U) \to X_U$를 합성
$(\rho^W_U)^{-1} \circ \alpha|_{g^{-1}(U)}$로 정의할 수 있다. 이는 상
$\alpha(g^{-1}(U))$가 $f_W^{-1}(U)$에 포함되고 보조정리의 조건 (a)가
성립하기 때문이다. 관계 $f_U \circ h_U = g|_{g^{-1}(U)}$가 그러한
$U$에 대해 성립함이 분명하다. 또한
$V \in \mathcal{B}$가 $V \subset U \subset W$도 만족하면, 보조정리의
성질 (b)에 의해 $\rho^U_V \circ h_V = h_U|_{g^{-1}(V)}$이다.
이제 $h_U$를 임의의 원소 $U \in \mathcal{B}$에 대해서도 정의해야 한다.
$\mathcal{B}$가 $S$의 위상의 기저이므로 열린 덮개
$U \cap W = \bigcup U_i$를 찾을 수 있으며, 여기서
$U_i \in \mathcal{B}$이다. $g$의 상은 $W$ 안에 있으므로
$g^{-1}(U) = g^{-1}(U \cap W) = \bigcup g^{-1}(U_i)$이다.
사상 $h_i = \rho^U_{U_i} \circ h_{U_i} : g^{-1}(U_i) \to X_U$들을
생각하자. 보조정리의 조건 (b)를 사용하면
$h_i|_{g^{-1}(U_i) \cap g^{-1}(U_j)} = h_j|_{g^{-1}(U_i) \cap g^{-1}(U_j)}$
임을 간단히 증명할 수 있다. 따라서 이 사상들은 붙어서 원하는 사상
$h_U : g^{-1}(U) \to X_U$를 준다. 계 $(g, \{h_U\})$가 $F_W(T)$의
원소이고 위의 표시된 화살표 아래에서 $\alpha$로 가는지 확인하는 (쉬운) 과정은
생략한다.

\medskip\noindent
다음으로 각 $F_W \subset F$가 열린 몰입으로 표현가능임을 확인한다.
이는 정의로부터 분명하다.

\medskip\noindent
마지막으로 모음 $(F_W)_{W \in \mathcal{B}}$가 $F$를 덮는지 확인해야 한다.
이는 구성과 $\mathcal{B}$가 $S$의 위상의 기저라는 사실로부터 분명하다.

\medskip\noindent
$X$를 함자 $F$를 표현하는 스킴이라 하자.
$(f, \{i_U\}) \in F(X)$를 ``보편족''이라 하자. 각 $F_W$가 (위에 표시된
함자들의 사상을 통해) $X_W$로 표현되므로, 원하는 대로
$i_W : f^{-1}(W) \to X_W$가 동형임을 알 수 있다. 보조정리가 증명되었다.
\end{proof}

\begin{lemma}
\label{constructions-lemma-relative-glueing-sheaves}
$S$를 스킴이라 하자.
$\mathcal{B}$를 $S$의 위상의 기저라 하자.
다음 데이터가 주어졌다고 하자.
\begin{enumerate}
\item 모든 $U \in \mathcal{B}$에 대해 스킴 $f_U : X_U \to U$, 즉
$U$ 위의 스킴이 주어진다.
\item 모든 $U \in \mathcal{B}$에 대해 준연접 층 $\mathcal{F}_U$, 즉
$X_U$ 위의 층이 주어진다.
\item 모든 쌍 $U, V \in \mathcal{B}$가 $V \subset U$를 만족할 때 사상
$\rho^U_V : X_V \to X_U$가 주어진다.
\item 모든 쌍 $U, V \in \mathcal{B}$가 $V \subset U$를 만족할 때 사상
$\theta^U_V : (\rho^U_V)^*\mathcal{F}_U \to \mathcal{F}_V$가 주어진다.
\end{enumerate}
다음을 가정하자.
\begin{enumerate}
\item[(a)] 각 $\rho^U_V$가 동형 $X_V \to f_U^{-1}(V)$, 즉 $V$ 위의
스킴들의 동형을 유도하고,
\item[(b)] 각 $\theta^U_V$가 동형이며,
\item[(c)] $W, V, U \in \mathcal{B}$이고 $W \subset V \subset U$이면
$\rho^U_W = \rho^U_V \circ \rho ^V_W$가 성립하고,
\item[(d)] $W, V, U \in \mathcal{B}$이고 $W \subset V \subset U$이면
$\theta^U_W = \theta^V_W \circ (\rho^V_W)^*\theta^U_V$가 성립한다.
\end{enumerate}
그러면 스킴의 사상 $f : X \to S$, 준연접 층 $\mathcal{F}$, 즉 $X$ 위의
층, 그리고 동형 $i_U : f^{-1}(U) \to X_U$와
$\theta_U : i_U^*\mathcal{F}_U \to \mathcal{F}|_{f^{-1}(U)}$가 각
$U \in \mathcal{B}$ 위에서 존재하여, $V, U \in \mathcal{B}$이고
$V \subset U$일 때 합성
$$
\xymatrix{
X_V \ar[r]^{i_V^{-1}} &
f^{-1}(V) \ar[rr]^{inclusion} & &
f^{-1}(U) \ar[r]^{i_U} &
X_U
}
$$
은 사상 $\rho^U_V$이고, 합성
\begin{equation}
\label{constructions-equation-glue}
(\rho^U_V)^*\mathcal{F}_U
=
(i_V^{-1})^*((i_U^*\mathcal{F}_U)|_{f^{-1}(V)})
\xrightarrow{\theta_U|_{f^{-1}(V)}}
(i_V^{-1})^*(\mathcal{F}|_{f^{-1}(V)})
\xrightarrow{\theta_V^{-1}}
\mathcal{F}_V
\end{equation}
은 $\theta^U_V$와 같다. 또한 $(X, \mathcal{F})$는 $S$ 위의 유일한
동형을 제외하고 유일하다.
\end{lemma}

\begin{proof}
Lemma \ref{constructions-lemma-relative-glueing}에 의해 스킴 $X$, 즉 $S$ 위의 스킴과
동형 $i_U$들을 얻는다. $\mathcal{F}'_U = i_U^*\mathcal{F}_U$로 두되,
$U \in \mathcal{B}$라 하자. 이는 준연접 $\mathcal{O}_{f^{-1}(U)}$-가군이다.
사상들
$$
\mathcal{F}'_U|_{f^{-1}(V)} =
i_U^*\mathcal{F}_U|_{f^{-1}(V)} =
i_V^*(\rho^U_V)^*\mathcal{F}_U \xrightarrow{i_V^*\theta^U_V}
i_V^*\mathcal{F}_V = \mathcal{F}'_V
$$
은 동형
$(\theta')^U_V : \mathcal{F}'_U|_{f^{-1}(V)} \to \mathcal{F}'_V$를
$V \subset U$인 $\mathcal{B}$의 원소들마다 정의한다. 조건 (d)는 세 원소
$W \subset V \subset U$가 $\mathcal{B}$에 속할 때 이들이 호환된다는 것을
정확히 말한다. 이에 따라
$$
\varphi_{12} :
\mathcal{F}'_{U_1}|_{f^{-1}(U_1 \cap U_2)}
\longrightarrow
\mathcal{F}'_{U_2}|_{f^{-1}(U_1 \cap U_2)}
$$
라는 잘 정의된 동형을 임의의 $U_1, U_2 \in \mathcal{B}$에 대해 얻을 수
있다. 이를 위해 교집합 $U_1 \cap U_2 = \bigcup V_j$를 원소
$V_j$가 $\mathcal{B}$에 속하도록 덮고 다음과 같이 취한다.
$$
\varphi_{12}|_{V_j} =
\left((\theta')^{U_2}_{V_j}\right)^{-1}
\circ
(\theta')^{U_1}_{V_j}.
$$
이 사상들이 실제로 하나의 $\varphi_{12}$로 붙는다는 확인과 층을 위한
붙이기 데이터의 쌍대순환 조건(Sheaves, Section
\ref{sheaves-section-glueing-sheaves}에서와 같은 조건)의 확인은 생략한다.
Sheaves, Lemma \ref{sheaves-lemma-glue-sheaves}에 의해 $\mathcal{F}$를
$X$ 위에서 얻는다. (\ref{constructions-equation-glue})의 확인은 생략한다.
\end{proof}

\begin{remark}
\label{constructions-remark-relative-glueing-functorial}
Lemmas \ref{constructions-lemma-relative-glueing}와
\ref{constructions-lemma-relative-glueing-sheaves}에서 설명한 구성에는 함자성이 있다.
구체적으로 Lemma \ref{constructions-lemma-relative-glueing}에서와 같은 두 데이터 모음
$(f_U : X_U \to U, \rho^U_V)$와
$(g_U : Y_U \to U, \sigma^U_V)$가 주어졌다고 하자.
모든 $U \in \mathcal{B}$에 대해 사상 $h_U : X_U \to Y_U$, 즉 $U$ 위의
사상이 주어지고, 제한 $\rho^U_V$ 및 $\sigma^U_V$와 호환된다고 하자.
함자성이란 이들로부터 스킴의 사상 $h : X \to Y$, 즉 $S$ 위의 사상이
생기며, 이를 제한하면 사상 $h_U$들을 되찾는다는 뜻이다. 여기서
$f : X \to S$는 데이터 $(f_U : X_U \to U, \rho^U_V)$에서 얻고,
$g : Y \to S$는 데이터 $(g_U : Y_U \to U, \sigma^U_V)$에서 얻는다.

\medskip\noindent
마찬가지로 Lemma \ref{constructions-lemma-relative-glueing-sheaves}에서와 같은 두 데이터 모음
$(f_U : X_U \to U, \mathcal{F}_U, \rho^U_V, \theta^U_V)$와
$(g_U : Y_U \to U, \mathcal{G}_U, \sigma^U_V, \eta^U_V)$가
주어졌다고 하자. 모든 $U \in \mathcal{B}$에 대해 사상
$h_U : X_U \to Y_U$, 즉 $U$ 위의 사상이 주어지고 제한 $\rho^U_V$ 및
$\sigma^U_V$와 호환되며, 사상
$\tau_U : h_U^*\mathcal{G}_U \to \mathcal{F}_U$가 사상 $\theta^U_V$ 및
$\eta^U_V$와 호환되도록 주어졌다고 하자. 함자성이란 이들로부터 스킴의 사상
$h : X \to Y$, 즉 $S$ 위의 사상이 생겨 이를 제한하면 사상 $h_U$들을
되찾고, 사상 $h^*\mathcal{G} \to \mathcal{F}$가 생겨 이를 제한하면 사상
$h_U$들을 되찾는다는 뜻이다. 여기서
$(f : X \to S, \mathcal{F})$는 데이터
$(f_U : X_U \to U, \mathcal{F}_U, \rho^U_V, \theta^U_V)$에서 얻고,
$(g : Y \to S, \mathcal{G})$는 데이터
$(g_U : Y_U \to U, \mathcal{G}_U, \sigma^U_V, \eta^U_V)$에서 얻는다.

\medskip\noindent
확인 과정과 상대적 붙이기 데이터와 상대적 대상 사이의
``범주의 동치''를 적절히 서술하는 일은 생략한다.
\end{remark}

\section{붙이기를 통한 상대 스펙트럼}
\label{constructions-section-spec-via-glueing}

\begin{situation}
\label{constructions-situation-relative-spec}
여기서 $S$는 스킴이고 $\mathcal{A}$는 준연접 $\mathcal{O}_S$-대수이다.
이는 $\mathcal{A}$가 $\mathcal{O}_S$-대수의 층이며 $\mathcal{O}_S$-가군으로서
준연접이라는 뜻이다.
\end{situation}

\noindent
이 절에서는 스킴의 사상
$$
\underline{\Spec}_S(\mathcal{A}) \longrightarrow S
$$
을 스펙트럼 $\Spec(\Gamma(U, \mathcal{A}))$들을 붙여 구성하는 방법을
개괄한다. 여기서 $U$는 $S$의 아핀 열린집합들을 달린다. 먼저 아핀
열린집합에서 취한 $\mathcal{A}$의
값들의 스펙트럼이 Lemma \ref{constructions-lemma-relative-glueing}에서와 같은 적절한 스킴
모음을 이룬다는 것을 보인다.

\begin{lemma}
\label{constructions-lemma-spec-inclusion}
Situation \ref{constructions-situation-relative-spec}에서 생각하자.
$U \subset U' \subset S$가 아핀 열린집합들이라고 하자.
$A = \mathcal{A}(U)$ 및 $A' = \mathcal{A}(U')$라 하자.
환의 사상 $A' \to A$는 사상 $\Spec(A) \to \Spec(A')$를 유도하고, 도식
$$
\xymatrix{
\Spec(A) \ar[r] \ar[d] &
\Spec(A') \ar[d] \\
U \ar[r] &
U'
}
$$
은 당김 정사각형이다.
\end{lemma}

\begin{proof}
$R = \mathcal{O}_S(U)$ 및 $R' = \mathcal{O}_S(U')$라 하자.
사상 $R \otimes_{R'} A' \to A$가 동형임에 유의하자. 이는
$\mathcal{A}$가 준연접이기 때문이다(예를 들어 Schemes, Lemma
\ref{schemes-lemma-widetilde-pullback}을 보라). 결과는 아핀 스킴의 올곱에 대한
Schemes, Lemma \ref{schemes-lemma-fibre-product-affine-schemes}의 기술로부터
따른다.
\end{proof}

\noindent
특히 이 보조정리의 사상 $\Spec(A) \to \Spec(A')$는 열린 몰입이다.

\begin{lemma}
\label{constructions-lemma-transitive-spec}
Situation \ref{constructions-situation-relative-spec}에서 생각하자.
$U \subset U' \subset U'' \subset S$가 아핀 열린집합들이라고 하자.
$A = \mathcal{A}(U)$, $A' = \mathcal{A}(U')$ 및
$A'' = \mathcal{A}(U'')$라 하자. Lemma \ref{constructions-lemma-spec-inclusion}의 사상
$\Spec(A) \to \Spec(A')$와
$\Spec(A') \to \Spec(A'')$의 합성은 Lemma
\ref{constructions-lemma-spec-inclusion}의 사상 $\Spec(A) \to \Spec(A'')$를 준다.
\end{lemma}

\begin{proof}
이는 사상 $A'' \to A$가 $A'' \to A'$와 $A' \to A$의 합성이기 때문에
따른다($\mathcal{A}$가 층이기 때문이다).
\end{proof}

\begin{lemma}
\label{constructions-lemma-glue-relative-spec}
Situation \ref{constructions-situation-relative-spec}에서 생각하자.
다음 성질을 지닌 스킴의 사상
$$
\pi : \underline{\Spec}_S(\mathcal{A}) \longrightarrow S
$$
가 존재한다.
\begin{enumerate}
\item 모든 아핀 열린집합 $U \subset S$에 대해 동형
$i_U : \pi^{-1}(U) \to \Spec(\mathcal{A}(U))$가 존재하며, 이는 $U$ 위의
동형이고,
\item 아핀 열린집합 $U \subset U' \subset S$에 대해 합성
$$
\xymatrix{
\Spec(\mathcal{A}(U)) \ar[r]^{i_U^{-1}} &
\pi^{-1}(U) \ar[rr]^{inclusion} & &
\pi^{-1}(U') \ar[r]^{i_{U'}} &
\Spec(\mathcal{A}(U'))
}
$$
은 위 Lemma \ref{constructions-lemma-spec-inclusion}의 열린 몰입이다.
\end{enumerate}
또한 $\underline{\Spec}_S(\mathcal{A})$는 $S$ 위의 유일한 동형을 제외하고
유일하다.
\end{lemma}

\begin{proof}
Lemmas \ref{constructions-lemma-relative-glueing},
\ref{constructions-lemma-spec-inclusion}, 그리고
\ref{constructions-lemma-transitive-spec}에서 곧바로 따른다.
유일성은 Lemma \ref{constructions-lemma-relative-glueing}의 마지막 문장에 서술되어 있다.
\end{proof}

\section{함자로서의 상대 스펙트럼}
\label{constructions-section-spec}

\noindent
Situation \ref{constructions-situation-relative-spec}의 상황, 즉 $S$가 스킴이고
$\mathcal{A}$가 $\mathcal{O}_S$-대수의 준연접 층인 상황에서 생각한다.

\medskip\noindent
임의의 $f : T \to S$에 대해 당김 $f^*\mathcal{A}$는
$\mathcal{O}_T$-대수의 준연접 층이다. 이제 쌍 $(f : T \to S, \varphi)$를
생각할 텐데, 여기서 $f$는 스킴의 사상이고
$\varphi : f^*\mathcal{A} \to \mathcal{O}_T$는 $\mathcal{O}_T$-대수의
사상이다. 이는 $f^{-1}\mathcal{O}_S$-대수 준동형
$\varphi : f^{-1}\mathcal{A} \to \mathcal{O}_T$를 주는 것과 같음에 유의하자.
Sheaves, Lemma \ref{sheaves-lemma-adjointness-tensor-restrict}을 보라. 또한 이는
$\mathcal{O}_S$-대수 사상 $\varphi : \mathcal{A} \to f_*\mathcal{O}_T$를
주는 것과 같다. Sheaves, Lemma
\ref{sheaves-lemma-adjoint-push-pull-modules}를 보라. 앞으로 별도의 언급 없이
$\varphi$를 생각하는 이 세 방식을 모두 사용한다.

\medskip\noindent
그러한 쌍 $(f : T \to S, \varphi)$와 사상 $a : T' \to T$가 주어지면
두 번째 쌍 $(f' = f \circ a, \varphi' = a^*\varphi)$를 얻으며, 이를
$(f, \varphi)$의 당김이라고 부른다. $\varphi' = a^*\varphi$를 기술하는 한
방법은 합성 $\mathcal{A} \to f_*\mathcal{O}_T \to f'_*\mathcal{O}_{T'}$으로
보는 것이다. 여기서 두 번째 사상은 $f_*a^\sharp$이고
$a^\sharp : \mathcal{O}_T \to a_*\mathcal{O}_{T'}$이다.
이로써 함자를 정의하였다.
\begin{eqnarray}
\label{constructions-equation-spec}
F : \Sch^{opp} & \longrightarrow & \textit{Sets} \\
T & \longmapsto & F(T) = \{\text{pairs }(f, \varphi) \text{ as above}\}
\nonumber
\end{eqnarray}

\begin{lemma}
\label{constructions-lemma-spec-base-change}
Situation \ref{constructions-situation-relative-spec}에서 생각하자.
함자를 $F$라 하자. 이는 위에서 $(S, \mathcal{A})$에 결부한 함자이다.
$g : S' \to S$를 스킴의 사상이라 하자.
$\mathcal{A}' = g^*\mathcal{A}$로 두자. 함자 $F'$를 위에서
$(S', \mathcal{A}')$에 결부한 함자라 하자. 그러면 함자들의 표준 동형
$$
F' \cong h_{S'} \times_{h_S} F
$$
이 존재한다.
\end{lemma}

\begin{proof}
쌍 $(f' : T \to S', \varphi' : (f')^*\mathcal{A}' \to \mathcal{O}_T)$를
주는 것은 쌍 $(f, \varphi : f^*\mathcal{A} \to \mathcal{O}_T)$와 함께
$f$의 분해 $f = g \circ f'$를 주는 것과 같다. 실제로 이 표기에서
$(f')^* \mathcal{A}' = (f')^*g^*\mathcal{A} = f^*\mathcal{A}$이다.
따라서 보조정리가 성립한다.
\end{proof}

\begin{lemma}
\label{constructions-lemma-spec-affine}
Situation \ref{constructions-situation-relative-spec}에서 생각하자.
함자 $F$를 위에서 $(S, \mathcal{A})$에 결부한 함자라 하자. $S$가 아핀이면
$F$는 아핀 스킴 $\Spec(\Gamma(S, \mathcal{A}))$으로 표현된다.
\end{lemma}

\begin{proof}
$S = \Spec(R)$ 및 $A = \Gamma(S, \mathcal{A})$라 쓰자.
그러면 $A$는 $R$-대수이고 $\mathcal{A} = \widetilde A$이다.
환의 사상 $R \to A$는 표준 사상
$$
f_{univ} : \Spec(A)
\longrightarrow
S = \Spec(R).
$$
을 준다. Schemes, Lemma \ref{schemes-lemma-widetilde-pullback}에 의해
$f_{univ}^*\mathcal{A} =  \widetilde{A \otimes_R A}$이다. 따라서 표준 사상
$$
\varphi_{univ} :
f_{univ}^*\mathcal{A} = \widetilde{A \otimes_R A}
\longrightarrow
\widetilde A = \mathcal{O}_{\Spec(A)}
$$
이 존재하며, 이는 $A$-가군 사상 $A \otimes_R A \to A$,
$a \otimes a' \mapsto aa'$에서 온다. 이 경우 쌍
$(f_{univ}, \varphi_{univ})$가 $F$를 표현한다고 주장한다. 다시 말해 임의의
스킴 $T$에 대해 사상
$$
\Mor(T, \Spec(A)) \longrightarrow \{\text{pairs } (f, \varphi)\},\quad
a \longmapsto (f_{univ} \circ a, a^*\varphi_{univ})
$$
이 전단사라고 주장한다.

\medskip\noindent
역사상을 구성하자. 임의의 쌍 $(f : T \to S, \varphi)$에서 유도되는 환의 사상
$$
\xymatrix{
A = \Gamma(S, \mathcal{A}) \ar[r]^{f^*} &
\Gamma(T, f^*\mathcal{A}) \ar[r]^{\varphi} &
\Gamma(T, \mathcal{O}_T)
}
$$
을 얻는다. 이는 Schemes, Lemma \ref{schemes-lemma-morphism-into-affine}에 의해
스킴의 사상 $T \to \Spec(A)$를 유도한다.

\medskip\noindent
이 사상이 위에 표시된 사상의 역이라는 확인은 생략한다.
\end{proof}

\begin{lemma}
\label{constructions-lemma-spec}
Situation \ref{constructions-situation-relative-spec}에서 생각하자.
함자 $F$는 스킴으로 표현가능이다.
\end{lemma}

\begin{proof}
Schemes, Lemma \ref{schemes-lemma-glue-functors}를 사용할 것이다.

\medskip\noindent
먼저 $F$가 자리스키 위상에 대한 층 조건을 만족함을 확인한다. 구체적으로
$T$가 스킴이고 $T = \bigcup_{i \in I} U_i$가 열린 덮개이며,
$(f_i, \varphi_i) \in F(U_i)$가
$(f_i, \varphi_i)|_{U_i \cap U_j} = (f_j, \varphi_j)|_{U_i \cap U_j}$를
만족한다고 하자. 그러면 사상 $f_i : U_i \to S$들은 스킴의 사상
$f : T \to S$로 붙으며 $f|_{U_i} = f_i$를 만족한다. Schemes, Section
\ref{schemes-section-glueing-schemes}를 보라. 따라서
$f_i^*\mathcal{A} = f^*\mathcal{A}|_{U_i}$이고, 가정에 의해 사상
$\varphi_i$들은 $U_i \cap U_j$ 위에서 일치한다. 그러므로 Sheaves, Section
\ref{sheaves-section-glueing-sheaves}에 의해 이들은 $\mathcal{O}_T$-대수의 사상
$f^*\mathcal{A} \to \mathcal{O}_T$로 붙는다. 이로써 $F$가 자리스키 위상에
대한 층 조건을 만족함을 증명하였다.

\medskip\noindent
$S = \bigcup_{i \in I} U_i$를 아핀 열린 덮개라 하자.
$F_i \subset F$를 쌍 $(f : T \to S, \varphi)$들 중
$f(T) \subset U_i$를 만족하는 것들로 이루어진 부분함자라 하자.

\medskip\noindent
각 $F_i$가 표현가능임을 보여야 한다. Lemma \ref{constructions-lemma-spec-base-change}에 의해
$F_i$는 $U_i$에 준연접 $\mathcal{O}_{U_i}$-대수
$\mathcal{A}|_{U_i}$를 갖추어 결부한 함자와 동일시되므로 이는 성립한다.
따라서 결과는 Lemma \ref{constructions-lemma-spec-affine}에서 따른다.

\medskip\noindent
다음으로 $F_i \subset F$가 열린 몰입으로 표현가능임을 보인다.
$(f : T \to S, \varphi) \in F(T)$라 하자. $V_i = f^{-1}(U_i)$를 생각하자.
$F_i$의 정의로부터 $a : T' \to T$가 주어졌을 때
$a^*(f, \varphi) \in F_i(T')$일 필요충분조건은 $a(T') \subset V_i$임을
알 수 있다. 이것이 보여야 했던 내용이다.

\medskip\noindent
마지막으로 모음 $(F_i)_{i \in I}$가 $F$를 덮는다는 것을 보여야 한다.
$(f : T \to S, \varphi) \in F(T)$라 하자. $V_i = f^{-1}(U_i)$를
생각하자. $S = \bigcup_{i \in I} U_i$가 $S$의 열린 덮개이므로
$T = \bigcup_{i \in I} V_i$가 $T$의 열린 덮개임을 알 수 있다. 또한
$(f, \varphi)|_{V_i} \in F_i(V_i)$이다. 이로써 보조정리의 증명이 끝난다.
\end{proof}

\begin{lemma}
\label{constructions-lemma-glueing-gives-functor-spec}
Situation \ref{constructions-situation-relative-spec}에서 생각하자.
Lemma \ref{constructions-lemma-glue-relative-spec}에서 구성한 스킴
$\pi : \underline{\Spec}_S(\mathcal{A}) \to S$와 함자 $F$를 표현하는 스킴은
$S$ 위의 스킴으로서 표준적으로 동형이다.
\end{lemma}

\begin{proof}
$X \to S$를 함자 $F$를 표현하는 스킴이라 하자.
$\mathcal{O}_S$-대수의 층
$\mathcal{R} = \pi_*\mathcal{O}_{\underline{\Spec}_S(\mathcal{A})}$를 생각하자.
$\underline{\Spec}_S(\mathcal{A})$의 구성에 의해 동형
$\mathcal{A}(U) \to \mathcal{R}(U)$가 모든 아핀 열린집합 $U \subset S$에
대해 존재한다.
이는 Lemma \ref{constructions-lemma-glue-relative-spec}의 (1)에서 따른다.
열린집합 $U \subset U' \subset S$에 대해 이 동형들은 제한 사상과
호환된다. 이는 Lemma \ref{constructions-lemma-glue-relative-spec}의 (2)에서 따른다.
따라서 Sheaves, Lemma \ref{sheaves-lemma-restrict-basis-equivalence-modules}에 의해
이 동형들은 $\mathcal{O}_S$-대수의 동형
$\varphi : \mathcal{A} \to \mathcal{R}$에서 온다. 따라서 이는 원소
$(\pi, \varphi) \in F(\underline{\Spec}_S(\mathcal{A}))$를 준다.
$X$가 함자 $F$를 표현하므로 이에 대응하는 스킴의 사상
$can : \underline{\Spec}_S(\mathcal{A}) \to X$, 즉 $S$ 위의 사상을 얻는다.

\medskip\noindent
$U \subset S$를 임의의 아핀 열린집합이라 하자. $F_U \subset F$를 $F$의
부분함자로서, 쌍 $(f, \varphi)$들 중 스킴 $T$ 위에 있고
$f(T) \subset U$를 만족하는 것들에 대응하도록 정의하자. 밑변환 $X_U$가 $F_U$를 표현함은
분명하다. 한편 Lemma \ref{constructions-lemma-spec-affine}에 따르면 $F_U$는
$\Spec(\mathcal{A}(U)) = \pi^{-1}(U)$로 표현된다. 다시 말해
$X_U \cong \pi^{-1}(U)$이다. 이 동일시가 사상 $can$을 $U$로 밑변환하여
얻어진다는 확인은 생략한다.
\end{proof}

\begin{definition}
\label{constructions-definition-relative-spec}
$S$를 스킴이라 하자. $\mathcal{A}$를 $\mathcal{O}_S$-대수의 준연접 층이라
하자. {\it $\mathcal{A}$의 $S$ 위 상대 스펙트럼}, 또는 간단히
{\it $\mathcal{A}$의 $S$ 위 스펙트럼}은 Lemma
\ref{constructions-lemma-glue-relative-spec}에서 구성되어 함자 $F$ (\ref{constructions-equation-spec})를
표현하는 스킴이다. Lemma \ref{constructions-lemma-glueing-gives-functor-spec}를 보라.
이를 $\pi : \underline{\Spec}_S(\mathcal{A}) \to S$로 나타낸다.
``보편족''은 $\mathcal{O}_S$-대수의 사상
$$
\mathcal{A}
\longrightarrow
\pi_*\mathcal{O}_{\underline{\Spec}_S(\mathcal{A})}
$$
이다.
\end{definition}

\noindent
다음 보조정리는 특히 상대 스펙트럼을 만드는 것이 밑변환과 가환함을 말한다.

\begin{lemma}
\label{constructions-lemma-spec-properties}
$S$를 스킴이라 하자. $\mathcal{A}$를 $\mathcal{O}_S$-대수의 준연접 층이라
하자. $\pi : \underline{\Spec}_S(\mathcal{A}) \to S$를 $\mathcal{A}$의
$S$ 위 상대 스펙트럼이라 하자.
\begin{enumerate}
\item 모든 아핀 열린집합 $U \subset S$에 대해 역상 $\pi^{-1}(U)$는 아핀이다.
\item 모든 사상 $g : S' \to S$에 대해
$S' \times_S \underline{\Spec}_S(\mathcal{A}) =
\underline{\Spec}_{S'}(g^*\mathcal{A})$이다.
\item
보편 사상
$$
\mathcal{A}
\longrightarrow
\pi_*\mathcal{O}_{\underline{\Spec}_S(\mathcal{A})}
$$
은 $\mathcal{O}_S$-대수의 동형이다.
\end{enumerate}
\end{lemma}

\begin{proof}
(1)은 상대 스펙트럼을 붙이기로 기술한 데서 온다. Lemma
\ref{constructions-lemma-glue-relative-spec}를 보라. (2)는 Lemma
\ref{constructions-lemma-spec-base-change}에서 곧바로 따른다. (3)은 $S$ 위에서 국소적인
명제이고, $S$가 아핀인 경우에는 예를 들어 Lemma \ref{constructions-lemma-spec-affine}에
의해 분명하므로 성립한다.
\end{proof}

\begin{lemma}
\label{constructions-lemma-canonical-morphism}
$f : X \to S$를 스킴의 준콤팩트 준분리 사상이라 하자. Schemes, Lemma
\ref{schemes-lemma-push-forward-quasi-coherent}에 의해 층
$f_*\mathcal{O}_X$는 $\mathcal{O}_S$-대수의 준연접 층이다. $S$ 위의
스킴의 표준 사상
$$
can : X \longrightarrow \underline{\Spec}_S(f_*\mathcal{O}_X)
$$
이 존재한다. 임의의 아핀 열린집합 $U \subset S$에 대해 제한
$can|_{f^{-1}(U)}$은 표준 사상
$$
f^{-1}(U) \longrightarrow \Spec(\Gamma(f^{-1}(U), \mathcal{O}_X))
$$
과 동일시된다. 이 사상은 Schemes, Lemma
\ref{schemes-lemma-morphism-into-affine}에서 온다.
\end{lemma}

\begin{proof}
이 사상은 $\underline{\Spec}$를 함자 $F$를 표현하는 스킴으로 정의한 데 따라
표준 사상 $\varphi : f^*f_*\mathcal{O}_X \to \mathcal{O}_X$에서 온다.
이 사상은 당김과 밂의 수반성에 의해
$\text{id} : f_*\mathcal{O}_X \to f_*\mathcal{O}_X$에 대응한다.
$f^{-1}(U)$로 제한한 것에 관한 명제는 아핀 위의 상대 스펙트럼에 대한 기술에서
따른다. Lemma \ref{constructions-lemma-spec-affine}를 보라.
\end{proof}

\section{아핀 n-공간}
\label{constructions-section-affine-n-space}

\noindent
상대 스펙트럼의 응용으로서, 밑 스킴 위의 아핀 $n$-공간을 다음과 같이 정의한다.
$S$를 밑 스킴이라 하자. 임의의 정수 $n \geq 0$에 대해
$\mathcal{O}_S$-대수의 준연접 층 $\mathcal{O}_S[T_1, \ldots, T_n]$을 생각할
수 있다. 이는 $\mathcal{O}_S$-가군의 층으로서 다중지표로 지표화된
$\mathcal{O}_S$의 복사본들의 직합일 뿐이므로 준연접이다.

\begin{definition}
\label{constructions-definition-affine-n-space}
$S$를 스킴이라 하고 $n \geq 0$이라 하자. $S$ 위의 스킴
$$
\mathbf{A}^n_S =
\underline{\Spec}_S(\mathcal{O}_S[T_1, \ldots, T_n])
$$
을 {\it 아핀 $n$-공간($S$ 위)}이라고 부른다. $S = \Spec(R)$가 아핀이면
이를 {\it 아핀 $n$-공간($R$ 위)}이라고도 부르고 $\mathbf{A}^n_R$로
나타낸다.
\end{definition}

\noindent
$\mathbf{A}^n_R = \Spec(R[T_1, \ldots, T_n])$임에 유의하자.
스킴의 임의의 사상 $g : S' \to S$에 대해
$g^*\mathcal{O}_S[T_1, \ldots, T_n] = \mathcal{O}_{S'}[T_1, \ldots, T_n]$이고,
따라서 $\mathbf{A}^n_{S'} = S' \times_S \mathbf{A}^n_S$는 밑변환이다.
그러므로 아핀 $n$-공간을 다음 공식으로 달리 정의할 수도 있다.
$$
\mathbf{A}^n_S = S \times_{\Spec(\mathbf{Z})} \mathbf{A}^n_{\mathbf{Z}}.
$$
또한 $S$-스킴 $f : X \to S$에서 $\mathbf{A}^n_S$로 가는 사상은
$\mathcal{O}_S$-대수의 준동형
$\mathcal{O}_S[T_1, \ldots, T_n] \to f_*\mathcal{O}_X$로 주어진다.
이는 분명히 $T_i$들의 상을 주는 것과 같다. 다시 말해 $X$에서
$\mathbf{A}^n_S$로 가는 $S$ 위의 사상을 주는 것은 $n$개의 원소
$h_1, \ldots, h_n \in \Gamma(X, \mathcal{O}_X)$를 주는 것과 같다.

\section{벡터 다발}
\label{constructions-section-vector-bundle}

\noindent
$S$를 스킴이라 하자. $\mathcal{E}$를 $\mathcal{O}_S$-가군의 준연접 층이라
하자. Modules, Lemma \ref{modules-lemma-whole-tensor-algebra-permanence}에 의해
대칭 대수 $\text{Sym}(\mathcal{E})$, 즉 $\mathcal{E}$의 $\mathcal{O}_S$ 위의
대칭 대수는 $\mathcal{O}_S$-대수의 준연접 층이다. 따라서 상대 스펙트럼을
적용하는 것이 타당하다.

\begin{definition}
\label{constructions-definition-vector-bundle}
$S$를 스킴이라 하자. $\mathcal{E}$를 준연접 $\mathcal{O}_S$-가군이라
하자\footnote{독자는 여기에서 $\mathcal{E}$가 유한 국소 자유라는 조건을
기대할 수 있다. 그러나 \cite[II, Definition 1.7.8]{EGA}와 일관되도록 그 조건을
부과하지 않는다.}. {\it $\mathcal{E}$에 결부한 벡터 다발}은
$$
\mathbf{V}(\mathcal{E}) = \underline{\Spec}_S(\text{Sym}(\mathcal{E})).
$$
이다.
\end{definition}

\noindent
$\mathcal{E}$에 결부한 벡터 다발에는 구조가 조금 더 주어진다. 즉, 등급화
$$
\pi_*\mathcal{O}_{\mathbf{V}(\mathcal{E})} =
\bigoplus\nolimits_{n \geq 0} \text{Sym}^n(\mathcal{E}).
$$
가 있으며, 이는 $\pi_*\mathcal{O}_{\mathbf{V}(\mathcal{E})}$를 등급
$\mathcal{O}_S$-대수로 만든다. 반대로 $\mathcal{E}$를 그 차수 $1$ 부분에서
복원할 수 있다. 따라서 추상 벡터 다발을 다음과 같이 정의한다.

\begin{definition}
\label{constructions-definition-abstract-vector-bundle}
$S$를 스킴이라 하자. {\it 벡터 다발 $\pi : V \to S$($S$ 위)}은 아핀
사상으로서, $\pi_*\mathcal{O}_V$에 등급 $\mathcal{O}_S$-대수의 구조
$\pi_*\mathcal{O}_V = \bigoplus\nolimits_{n \geq 0} \mathcal{E}_n$이 주어지고,
$\mathcal{E}_0 = \mathcal{O}_S$이며, 사상
$$
\text{Sym}^n(\mathcal{E}_1) \longrightarrow \mathcal{E}_n
$$
이 모든 $n \geq 0$에 대해 동형인 스킴의 사상이다. {\it $S$ 위의 벡터
다발의 사상}은 유도된 사상
$f : V \to V'$가
$$
f^* : \pi'_*\mathcal{O}_{V'} \longrightarrow \pi_*\mathcal{O}_V
$$
로 주어진 등급화와 호환되도록 하는 사상이다.
\end{definition}

\noindent
$S$ 위 벡터 다발의 한 예는 아핀 $n$-공간 $\mathbf{A}^n_S$, 즉 $S$ 위의
아핀 공간이다. Definition \ref{constructions-definition-affine-n-space}를 보라. 이는
$\mathcal{O}_S[T_1, \ldots, T_n] = \text{Sym}(\mathcal{O}_S^{\oplus n})$이기
때문에 성립한다.

\begin{lemma}
\label{constructions-lemma-category-vector-bundles}
스킴 $S$ 위의 벡터 다발의 범주는 준연접 $\mathcal{O}_S$-가군의 범주와
반동치이다.
\end{lemma}

\begin{proof}
생략한다. 힌트: 한 방향에서는 함자
$\underline{\Spec}_S(\text{Sym}^*_{\mathcal{O}_S}(-))$를 사용하고, 다른
방향에서는 함자 $(\pi : V \to S) \leadsto (\pi_*\mathcal{O}_V)_1$을 사용한다.
여기서 아래첨자는 차수 $1$ 부분을 취한다는 뜻이다.
\end{proof}

\section{원뿔}
\label{constructions-section-cone}

\noindent
대수기하학에서 원뿔은 등급 대수에 대응한다. 우리의 규약에 따르면 등급 환 또는
대수 $A$에는 음이 아닌 정수에 의한 등급화
$A = \bigoplus_{d \geq 0} A_d$가 주어진다. Algebra, Section
\ref{algebra-section-graded}를 보라.

\begin{definition}
\label{constructions-definition-cone}
$S$를 스킴이라 하자. $\mathcal{A}$를 등급 준연접 $\mathcal{O}_S$-대수라
하자. $\mathcal{O}_S \to \mathcal{A}_0$가 동형이라고
가정하자\footnote{흔히 $\mathcal{A}$가 $\mathcal{A}_1$에 의해
$\mathcal{O}_S$ 위에서 생성된다고 가정한다. 우리는 \cite[II, (8.3.1)]{EGA}와
일관되도록 이를 가정하지 않는다.}. {\it $\mathcal{A}$에 결부한 원뿔}, 또는
{\it $\mathcal{A}$에 결부한 아핀 원뿔}은
$$
C(\mathcal{A}) = \underline{\Spec}_S(\mathcal{A}).
$$
이다.
\end{definition}

\noindent
등급 $\mathcal{O}_S$-대수의 층에 결부한 원뿔에는 구조가 조금 더 주어진다.
즉, 등급화
$$
\pi_*\mathcal{O}_{C(\mathcal{A})} =
\bigoplus\nolimits_{n \geq 0} \mathcal{A}_n
$$
를 얻는다. 따라서 추상 원뿔을 다음과 같이 정의할 수 있다.

\begin{definition}
\label{constructions-definition-abstract-cone}
$S$를 스킴이라 하자. {\it 원뿔 $\pi : C \to S$($S$ 위)}은 아핀 사상으로서,
$\pi_*\mathcal{O}_C$에 등급 $\mathcal{O}_S$-대수의 구조
$\pi_*\mathcal{O}_C = \bigoplus\nolimits_{n \geq 0} \mathcal{A}_n$이 주어지고
$\mathcal{A}_0 = \mathcal{O}_S$를 만족하는 스킴의 사상이다.
{\it 원뿔의 사상} $\pi : C \to S$에서 $\pi' : C' \to S$로 가는 것은
사상 $f : C \to C'$로서, 유도된 사상
$$
f^* : \pi'_*\mathcal{O}_{C'} \longrightarrow \pi_*\mathcal{O}_C
$$
이 주어진 등급화와 호환되는 것이다.
\end{definition}

\noindent
모든 벡터 다발은 원뿔의 예이다. 실제로 $S$ 위 벡터 다발의 범주는 $S$ 위
원뿔의 범주의 충만한 부분범주이다.

\section{등급 환의 Proj}
\label{constructions-section-proj}

\noindent
이 절에서는 \cite[II, Section 2]{EGA}를 따라 등급 환의 Proj를 구성한다.

\medskip\noindent
$S$를 등급 환이라 하자. 위상공간 $\text{Proj}(S)$를 생각하자. 이는 $S$에
결부한 공간이다.
Algebra, Section \ref{algebra-section-proj}를 보라. 이 공간에 환의 층
$\mathcal{O}_{\text{Proj}(S)}$를 주어, 얻어진 쌍
$(\text{Proj}(S), \mathcal{O}_{\text{Proj}(S)})$이 스킴이 되게 할 것이다.

\medskip\noindent
$\text{Proj}(S)$에는 열린집합 $D_{+}(f)$들로 이루어진 기저가 있음을
상기하자. 여기서 $f \in S_d$, $d \geq 1$이며, 이 열린집합들을 {\it 표준
열린집합}이라고 부른다. Algebra, Section \ref{algebra-section-proj}를 보라.
이 용어는 명시하지 않더라도 언제나 $f$가 양의 차수를 갖는 동차원소임을
뜻한다. 또한 두 표준 열린집합의 교집합도 표준 열린집합이며,
$D_{+}(f) \cap D_{+}(g) = D_{+}(fg)$이다. 여기서 $f, g \in S$는 양의
차수를 갖는 동차원소들이다.

\begin{lemma}
\label{constructions-lemma-standard-open}
$S$를 등급 환이라 하자. $f \in S$를 양의 차수를 갖는 동차원소라 하자.
\begin{enumerate}
\item $g\in S$가 양의 차수를 갖는 동차원소이고
$D_{+}(g) \subset D_{+}(f)$이면 다음이 성립한다.
\begin{enumerate}
\item $f$는 $S_g$에서 가역이고,
$f^{\deg(g)}/g^{\deg(f)}$는 $S_{(g)}$에서 가역이다.
\item $g^e = af$이며, 여기서 어떤 $e \geq 1$ 및 동차원소 $a \in S$가
존재한다.
\item 표준 $S$-대수 사상 $S_f \to S_g$가 존재한다.
\item 표준 $S_0$-대수 사상 $S_{(f)} \to S_{(g)}$가 존재하며, 사상
$S_f \to S_g$와 호환된다.
\item 사상 $S_{(f)} \to S_{(g)}$는 동형
$$
(S_{(f)})_{g^{\deg(f)}/f^{\deg(g)}} \cong S_{(g)},
$$
을 유도한다.
\item 이 사상들은 위상공간의 가환도식
$$
\xymatrix{
D_{+}(g) \ar[d] &
\{\mathbf{Z}\text{-graded primes of }S_g\} \ar[l] \ar[r] \ar[d] &
\Spec(S_{(g)}) \ar[d] \\
D_{+}(f) &
\{\mathbf{Z}\text{-graded primes of }S_f\} \ar[l] \ar[r] &
\Spec(S_{(f)})
}
$$
을 유도하며, 수평 사상들은 위상동형이고 수직 사상들은 열린 몰입이다.
\item 서로 호환되는 표준 $S_f$-가군 및 $S_{(f)}$-가군 사상
$M_f \to M_g$ 및 $M_{(f)} \to M_{(g)}$가 임의의 등급 $S$-가군 $M$에 대해
존재한다.
\item 사상 $M_{(f)} \to M_{(g)}$는 동형
$$
(M_{(f)})_{g^{\deg(f)}/f^{\deg(g)}} \cong M_{(g)}.
$$
을 유도한다.
\end{enumerate}
\item $D_{+}(f)$의 임의의 열린 덮개는 꼴
$D_{+}(f) = \bigcup_{i = 1}^n D_{+}(g_i)$의 유한 열린 덮개로 세분된다.
\item $g_1, \ldots, g_n \in S$를 양의 차수를 갖는 동차원소들이라 하자.
$D_{+}(f) \subset \bigcup D_{+}(g_i)$일 필요충분조건은
$g_1^{\deg(f)}/f^{\deg(g_1)}, \ldots, g_n^{\deg(f)}/f^{\deg(g_n)}$가
$S_{(f)}$에서 단위원 아이디얼을 생성하는 것이다.
\end{enumerate}
\end{lemma}

\begin{proof}
$D_{+}(g) = \Spec(S_{(g)})$이며, 그 동일시는 환의 사상
$S \to S_g \leftarrow S_{(g)}$가 줌을 상기하자. Algebra, Lemma
\ref{algebra-lemma-topology-proj}를 보라. 따라서
$f^{\deg(g)}/g^{\deg(f)}$는 $S_{(g)}$의 원소로서 어느 소 아이디얼에도
속하지 않으므로 가역이다. Algebra, Lemma
\ref{algebra-lemma-Zariski-topology}를 보라. 이로써 (a)가 성립한다.
$f$의 $S_g$에서의 역원을 $a/g^d$로 쓰자. $a$를 차수
$d\deg(g) - \deg(f)$인 동차성분으로 바꾸어도 된다. 이는
$g^d - af$가 $g$의 어떤 거듭제곱에 의해 소멸된다는 뜻이므로
$g^e = af$이고, 여기서 어떤 동차원소 $a \in S$의 차수는
$e\deg(g) - \deg(f)$이다. 이로써 (b)를 증명하였다.
(c)의 경우, (a)와 국소화의 보편성에 의해 사상 $S_f \to S_g$가 존재한다.
또는 $b/f^n$을 $a^nb/g^{ne}$로 보내어 정의할 수 있다. 이는 차수 영인
원소들의 부분환에도 사상 $S_{(f)} \to S_{(g)}$를 분명히 유도한다.
마찬가지로 $M_f \to M_g$와 $M_{(f)} \to M_{(g)}$를 $x/f^n$을
$a^nx/g^{ne}$로 보내어 정의할 수 있다. $S_{(g)}$ 및 $M_{(g)}$를 각각
$S_{(f)}$ 및 $M_{(f)}$의 주국소화로 나타낸 명제들은 위 공식에서 분명하다.
위상공간의 가환도식에 있는 사상들은 위에서 주어진 환의 사상들에 대응한다.
수평 화살표들은 Algebra, Lemma \ref{algebra-lemma-topology-proj}에 의해
위상동형이다. 왼쪽 수직 화살표가 열린 부분집합의 포함이므로 수직 화살표들은
열린 몰입이다.

\medskip\noindent
열린집합 $D_{+}(f)$는 $\Spec(S_{(f)})$와 위상동형이므로 준콤팩트이다.
Algebra, Lemma \ref{algebra-lemma-quasi-compact}를 보라. 따라서 표준
열린집합들이 위상의 기저를 이룬다는 사실에서 두 번째 명제가 바로 따른다.

\medskip\noindent
세 번째 명제는 Algebra, Lemma \ref{algebra-lemma-Zariski-topology}에서
바로 따른다.
\end{proof}

\noindent
Sheaves, Section \ref{sheaves-section-bases}에서 기저 위의 층이라는 개념을
정의하고, 이것이 공간 위의 층이라는 개념과 본질적으로 동치임을 보였다.
Sheaves, Lemmas \ref{sheaves-lemma-extend-off-basis} 및
\ref{sheaves-lemma-extend-off-basis-structures}를 보라. 또한 Sheaves, Lemma
\ref{sheaves-lemma-cofinal-systems-coverings-standard-case}에서 각 표준 열린집합에
대해 열린 덮개들의 공종 체계에서 층 조건을 확인하는 것으로 충분함을 보였다.
위 보조정리에 의해 표준 열린집합들로 이루어진 유한 덮개에서 확인하면 충분하다.

\begin{definition}
\label{constructions-definition-standard-covering}
$S$를 등급 환이라 하자. $D_{+}(f) \subset \text{Proj}(S)$가 표준
열린집합이라고 하자. $D_{+}(f)$의 {\it 표준 열린 덮개}란 덮개
$D_{+}(f) = \bigcup_{i = 1}^n D_{+}(g_i)$를 말하며, 여기서
$g_1, \ldots, g_n \in S$는 양의 차수를 갖는 동차원소들이다.
\end{definition}

\noindent
$S$를 등급 환이라 하자. $M$을 등급 $S$-가군이라 하자. 표준 열린집합들의
기저 위에서 앞층 $\widetilde M$을 정의할 것이다.
$U \subset \text{Proj}(S)$가 표준 열린집합이라고 하자.
$f, g \in S$가 양의 차수를 갖는 동차원소들이며
$D_{+}(f) = D_{+}(g)$라고 하자. 위 Lemma \ref{constructions-lemma-standard-open}에 의해
서로 역인 표준 사상 $M_{(f)} \to M_{(g)}$와 $M_{(g)} \to M_{(f)}$가
존재한다. 따라서 임의의 $f$를 $U = D_{+}(f)$를 만족하도록 택하여
$$
\widetilde M(U) = M_{(f)}.
$$
로 정의할 수 있다. $D_{+}(g) \subset D_{+}(f)$이면, 위 Lemma
\ref{constructions-lemma-standard-open}에 의해 표준 사상
$$
\widetilde M(D_{+}(f)) = M_{(f)} \longrightarrow
M_{(g)} = \widetilde M(D_{+}(g)).
$$
이 존재함에 유의하자. 이는 표준 열린집합들의 기저 위에서 아벨군의 앞층을
분명히 정의한다. $M = S$이면 $\widetilde S$는 표준 열린집합들의 기저 위에서
환의 앞층이다. 일반적인 $M$에 대해서는 $\widetilde M$이 표준 열린집합들의
기저 위에서 $\widetilde S$-가군의 앞층임을 알 수 있다.

\medskip\noindent
$\widetilde M$의 점 $x \in \text{Proj}(S)$에서의 줄기를 계산하자.
$x$가 동차 소 아이디얼 $\mathfrak p \subset S$에 대응한다고 하자.
줄기의 정의에 의해
$$
\widetilde M_x
=
\colim_{f\in S_d, d > 0, f\not\in \mathfrak p} M_{(f)}
$$
이다. 여기서 집합 $\{f \in S_d, d > 0, f \not \in \mathfrak p\}$에는 규칙
$f \geq f' \Leftrightarrow D_{+}(f) \subset D_{+}(f')$으로 전순서를 준다.
$f_1, f_2 \in S \setminus \mathfrak p$가 양의 차수를 갖는 동차원소들이면 이
순서에서 $f_1f_2 \geq f_1$이다. Algebra, Section
\ref{algebra-section-proj}에서 $M_{(\mathfrak p)}$를 원소가 분수 $x/f$인
가군으로 정의하였다. 여기서 $x, f$는 동차원소이고
$\deg(x) = \deg(f)$이며 $f \not \in \mathfrak p$이다.
$\mathfrak p \in \text{Proj}(S)$이므로 양의 차수를 갖는 동차원소
$f_0 \in S$ 중 $f_0 \not\in \mathfrak p$인 것이 적어도 하나 존재한다.
따라서 $x/f = f_0x/ff_0$이고, $M_{(\mathfrak p)}$의 원소의 분모는 언제나
양의 차수를 갖는다고 가정해도 됨을 알 수 있다. 이 고찰들로부터
$$
\widetilde M_x = M_{(\mathfrak p)}.
$$
가 쉽게 따른다.

\medskip\noindent
다음으로 표준 열린 덮개에 대한 층 조건을 확인한다.
$D_{+}(f) = \bigcup_{i = 1}^n D_{+}(g_i)$이면 이 덮개에 대한 층 조건은 수열
$$
0 \to M_{(f)} \to \bigoplus M_{(g_i)} \to \bigoplus M_{(g_ig_j)}.
$$
의 완전성과 동치이다. $D_{+}(g_i) = D_{+}(fg_i)$임에 유의하자. 따라서 이
수열을 다음 수열로 다시 쓸 수 있다.
$$
0 \to M_{(f)} \to \bigoplus M_{(fg_i)} \to \bigoplus M_{(fg_ig_j)}.
$$
Lemma \ref{constructions-lemma-standard-open}에 의해
$g_1^{\deg(f)}/f^{\deg(g_1)}, \ldots, g_n^{\deg(f)}/f^{\deg(g_n)}$는
$S_{(f)}$에서 단위원 아이디얼을 생성하고, 가군 $M_{(fg_i)}$ 및
$M_{(fg_ig_j)}$는 $S_{(f)}$-가군 $M_{(f)}$을 이 원소들 및 그 곱에서
주국소화한 것이다. 따라서 가군 $M_{(f)}$, 즉 $S_{(f)}$ 위의 가군과 원소들
$g_1^{\deg(f)}/f^{\deg(g_1)}, \ldots, g_n^{\deg(f)}/f^{\deg(g_n)}$에
Algebra, Lemma \ref{algebra-lemma-cover-module}을 적용할 수 있다. 이로부터
수열이 완전임을 얻는다. 위 고찰들에 의해 $\widetilde M$은 표준
열린집합들의 기저 위의 층이다.

\medskip\noindent
따라서 Sheaves, Section \ref{sheaves-section-bases}의 내용으로부터 유일한
환의 층 $\mathcal{O}_{\text{Proj}(S)}$가 존재하며 표준 열린집합 위에서
$\widetilde S$와 일치한다고 결론 내린다. 위의 줄기 계산과
Algebra, Lemma \ref{algebra-lemma-proj-prime}에 의해 이 환의 층의 모든
줄기는 국소환임에 유의하자.

\medskip\noindent
마찬가지로 임의의 등급 $S$-가군 $M$에 대해 유일한
$\mathcal{O}_{\text{Proj}(S)}$-가군의 층 $\mathcal{F}$가 존재하며, 표준
열린집합 위에서 $\widetilde M$과 일치한다. Sheaves, Lemma
\ref{sheaves-lemma-extend-off-basis-module}을 보라.

\begin{definition}
\label{constructions-definition-structure-sheaf}
$S$를 등급 환이라 하자.
\begin{enumerate}
\item {\it 구조층 $\mathcal{O}_{\text{Proj}(S)}$($S$의 동차 스펙트럼의)}은
유일한 환의 층 $\mathcal{O}_{\text{Proj}(S)}$이며, 표준 열린집합들의 기저
위에서 $\widetilde S$와 일치한다.
\item 국소 환 달린 공간 $(\text{Proj}(S), \mathcal{O}_{\text{Proj}(S)})$을
$S$의 {\it 동차 스펙트럼}이라 부르고 $\text{Proj}(S)$로 나타낸다.
\item $\mathcal{O}_{\text{Proj}(S)}$-가군의 층 중 $\widetilde M$을
$\text{Proj}(S)$의 모든 열린집합으로 확장하는 것을
$\mathcal{O}_{\text{Proj}(S)}$-가군의 층으로서 $M$에 결부한 층이라고
부른다. 이 층도 $\widetilde M$으로 나타낸다.
\end{enumerate}
\end{definition}

\noindent
지금까지 얻은 결과들을 요약한다.

\begin{lemma}
\label{constructions-lemma-proj-sheaves}
$S$를 등급 환이라 하자. $M$을 등급 $S$-가군이라 하자. $\widetilde M$을
$\mathcal{O}_{\text{Proj}(S)}$-가군의 층으로서 $M$에 결부한 층이라 하자.
\begin{enumerate}
\item 양의 차수를 갖는 모든 동차원소 $f \in S$에 대해
$$
\Gamma(D_{+}(f), \mathcal{O}_{\text{Proj}(S)}) = S_{(f)}.
$$
이다.
\item 양의 차수를 갖는 모든 동차원소 $f\in S$에 대해
$\Gamma(D_{+}(f), \widetilde M) = M_{(f)}$이며, 이는 $S_{(f)}$-가군의
동일시이다.
\item $D_{+}(g) \subset D_{+}(f)$일 때마다 $\mathcal{O}_{\text{Proj}(S)}$와
$\widetilde M$ 위의 제한 사상은 Lemma \ref{constructions-lemma-standard-open}의 사상
$S_{(f)} \to S_{(g)}$ 및 $M_{(f)} \to M_{(g)}$이다.
\item $\mathfrak p$를 $S$의 동차 소 아이디얼로서 $S_{+}$를 포함하지 않는 것이라 하고,
$x \in \text{Proj}(S)$를 이에 대응하는 점이라 하자. 그러면
$\mathcal{O}_{\text{Proj}(S), x} = S_{(\mathfrak p)}$이다.
\item $\mathfrak p$를 $S$의 동차 소 아이디얼로서 $S_{+}$를 포함하지 않는 것이라 하고,
$x \in \text{Proj}(S)$를 이에 대응하는 점이라 하자. 그러면
$(\widetilde M)_x = M_{(\mathfrak p)}$이며, 이는
$S_{(\mathfrak p)}$-가군의 동일시이다.
\item
\label{constructions-item-map}
표준 환의 사상
$
S_0 \longrightarrow \Gamma(\text{Proj}(S), \widetilde S)
$
과 표준 $S_0$-가군 사상
$
M_0 \longrightarrow \Gamma(\text{Proj}(S), \widetilde M)
$
이 존재하며, 위의 표준 열린집합 위의 단면 및 줄기에 대한 기술과 호환된다.
\end{enumerate}
또한 이 동일시들은 모두 등급 $S$-가군 $M$에 대해 함자적이다. 특히 함자
$M \mapsto \widetilde M$은 등급 $S$-가군의 범주에서
$\mathcal{O}_{\text{Proj}(S)}$-가군의 범주로 가는 완전 함자이다.
\end{lemma}

\begin{proof}
(1)--(5)는 위 논의에서 분명하다. (6)은 제한 사상과 호환되는 표준 사상
$M_0 \to M_{(f)}$, $x \mapsto x/1$이 존재하기 때문에 성립한다. 여기서 제한
사상은 (3)에 기술되어 있다. 함자 $M \mapsto \widetilde M$의 완전성은 함자
$M \mapsto M_{(\mathfrak p)}$가 완전이라는 사실(Algebra, Lemma
\ref{algebra-lemma-proj-prime}를 보라)과 짧은 완전열의 완전성을 줄기에서 확인할
수 있다는 사실에서 따른다. Modules, Lemma \ref{modules-lemma-abelian}을 보라.
\end{proof}

\begin{remark}
\label{constructions-remark-global-sections-not-isomorphism}
$M_0$에서 $\widetilde M$의 전역 단면으로 가는 사상은 일반적으로 동형과는
거리가 멀다. 자명한 예로 $S = k[x, y, z]$를
$1 = \deg(x) = \deg(y) = \deg(z)$가 되게 취하고(또는 임의 개수의 변수를
취하고)
$M = S/(x^{100}, y^{100}, z^{100})$을 취하자. 그러면
$\widetilde M = 0$이지만 $M_0 = k$임을 쉽게 알 수 있다.
\end{remark}

\begin{lemma}
\label{constructions-lemma-standard-open-proj}
$S$를 등급 환이라 하자. $f \in S$를 양의 차수를 갖는 동차원소라 하자.
$D(g) \subset \Spec(S_{(f)})$가 표준 열린집합이라고 하자. 그러면 양의 차수를
갖는 어떤 동차원소 $h \in S$가 존재하여, Algebra, Lemma
\ref{algebra-lemma-topology-proj}의 위상동형 아래 $D(g)$가
$D_{+}(h) \subset D_{+}(f)$에 대응한다. 실제로 $h$를
$g = h/f^n$을 어떤 $n$에 대해 만족하도록 취할 수 있다.
\end{lemma}

\begin{proof}
$g = h/f^n$으로 쓰되, 여기서 $h$는 양의 차수를 갖는 어떤 동차원소이고
$n \geq 1$이라 하자. $D_{+}(h)$가 $D_{+}(f)$에 포함되지 않으면 $h$를
$hf$로, $n$을 $n + 1$로 바꾼다. 그러면 $h$는 필요한 꼴을 가지며
$D_{+}(h) \subset D_{+}(f)$는 $D(g) \subset \Spec(S_{(f)})$에 대응한다.
\end{proof}

\begin{lemma}
\label{constructions-lemma-proj-scheme}
$S$를 등급 환이라 하자. 국소 환 달린 공간 $\text{Proj}(S)$는 스킴이다.
표준 열린집합 $D_{+}(f)$들은 아핀 열린집합이다. 임의의 등급 $S$-가군 $M$에
대해 층 $\widetilde M$은 준연접 $\mathcal{O}_{\text{Proj}(S)}$-가군의 층이다.
\end{lemma}

\begin{proof}
표준 열린집합 $D_{+}(f) \subset \text{Proj}(S)$를 생각하자. Lemmas
\ref{constructions-lemma-standard-open} 및 \ref{constructions-lemma-proj-sheaves}에 의해
$\Gamma(D_{+}(f), \mathcal{O}_{\text{Proj}(S)}) = S_{(f)}$이고, 위상동형
$\varphi : D_{+}(f) \to \Spec(S_{(f)})$가 존재한다. 임의의 표준 열린집합
$D(g) \subset \Spec(S_{(f)})$에 대해 Lemma
\ref{constructions-lemma-standard-open-proj}에서와 같은 $h \in S_{+}$를 택할 수 있다.
그러면 $\varphi^{-1}(D(g)) = D_{+}(h)$이고, Lemmas
\ref{constructions-lemma-proj-sheaves} 및 \ref{constructions-lemma-standard-open}에 의해
$$
\Gamma(D_{+}(h), \mathcal{O}_{\text{Proj}(S)})
=
S_{(h)}
=
(S_{(f)})_{h^{\deg(f)}/f^{\deg(h)}}
=
(S_{(f)})_g
=
\Gamma(D(g), \mathcal{O}_{\Spec(S_{(f)})}).
$$
임을 알 수 있다. 따라서 $\mathcal{O}_{\text{Proj}(S)}$를 $D_{+}(f)$에
제한한 것은 위상동형 $\varphi$를 통해 Schemes, Section
\ref{schemes-section-affine-schemes}에서 정의한 층
$\mathcal{O}_{\Spec(S_{(f)})}$에 정확히 대응한다. 이에 따라
$D_{+}(f)$는 $\Spec(S_{(f)})$와 $\varphi$를 통해 동형인 아핀 스킴이고,
$\text{Proj}(S)$는 스킴이라고 결론 내린다.

\medskip\noindent
똑같은 방식으로 $\widetilde M$이 준연접
$\mathcal{O}_{\text{Proj}(S)}$-가군의 층임을 보인다. 실제로 위 논증은
$$
\widetilde M|_{D_{+}(f)} \cong \varphi^*\left(\widetilde{M_{(f)}}\right)
$$
을 보여 주고, 이는 $\widetilde M$이 준연접임을 보인다.
\end{proof}

\begin{lemma}
\label{constructions-lemma-proj-separated}
$S$를 등급 환이라 하자. 스킴 $\text{Proj}(S)$는 분리이다.
\end{lemma}

\begin{proof}
표준 사상 $\text{Proj}(S) \to \Spec(\mathbf{Z})$가 분리임을 보여야 한다.
Schemes, Lemma \ref{schemes-lemma-characterize-separated}를 사용할 것이다.
따라서 임의의 두 표준 열린집합 $D_{+}(f)$와 $D_{+}(g)$에 대해
$D_{+}(f) \cap D_{+}(g) = D_{+}(fg)$가 아핀임(이는 분명하다)과 환의 사상
$$
S_{(f)} \otimes_{\mathbf{Z}} S_{(g)} \longrightarrow S_{(fg)}
$$
이 전사임을 보이면 충분하다. 임의의 원소 $s$는 $S_{(fg)}$에서
$s = h/(f^ng^m)$ 꼴이며, 여기서 $h \in S$는 차수
$n\deg(f) + m\deg(g)$인 동차원소이다. $h$에 적당한 단항식 $f^ig^j$를
곱하여 $n = n' \deg(g)$이고 $m = m' \deg(f)$이라고 가정할 수 있다.
그러면 $s$를
$s = h/f^{(n' + m')\deg(g)} \cdot f^{m'\deg(g)}/g^{m'\deg(f)}$로 다시
쓸 수 있다. 따라서 $s$는 실제로 표시된 화살표의 상에 속한다.
\end{proof}

\begin{lemma}
\label{constructions-lemma-proj-quasi-compact}
$S$를 등급 환이라 하자. 스킴 $\text{Proj}(S)$가 준콤팩트일 필요충분조건은
유한 개의 동차원소 $f_1, \ldots, f_n \in S_{+}$가 존재하여
$S_{+} \subset \sqrt{(f_1, \ldots, f_n)}$를 만족하는 것이다. 이 경우
$\text{Proj}(S) = D_+(f_1) \cup \ldots \cup D_+(f_n)$이다.
\end{lemma}

\begin{proof}
그러한 원소 모음이 주어지면 표준 아핀 열린집합 $D_{+}(f_i)$들이
$\text{Proj}(S)$를 덮는다. Algebra, Lemma
\ref{algebra-lemma-topology-proj}를 보라. 반대로 $\text{Proj}(S)$가
준콤팩트이면 유한 개의 표준 열린집합 $D_{+}(f_i)$,
$i = 1, \ldots, n$으로 덮을 수 있고, 위에서 인용한 보조정리에 의해
$S_{+} \subset \sqrt{(f_1, \ldots, f_n)}$임을 알 수 있다.
\end{proof}

\begin{lemma}
\label{constructions-lemma-structure-morphism-proj}
$S$를 등급 환이라 하자. 스킴 $\text{Proj}(S)$에는 아핀 스킴
$\Spec(S_0)$로 가는 표준 사상이 있으며, 이는 Algebra, Definition
\ref{algebra-definition-proj}에서 오는 위상공간의 사상과 일치한다.
\end{lemma}

\begin{proof}
위에서 $\widetilde S$ 및 $\widetilde M$의 구성이 각각 $S_0$-대수의 층 및
$S_0$-가군의 층을 준다는 것을 보았다. 따라서 Schemes, Lemma
\ref{schemes-lemma-morphism-into-affine}에 의해 사상을 얻는다. 이 사상을
$D_{+}(f)$에 제한하면 표준 환의 사상 $S_0 \to S_{(f)}$에서 온다. 사상
$S \to S_f$, $S_{(f)} \to S_f$는 $S_0$-대수 사상이다. Lemma
\ref{constructions-lemma-standard-open}을 보라. 따라서 동차 소 아이디얼
$\mathfrak p \subset S$가 $\mathbf{Z}$-등급 소 아이디얼
$\mathfrak p' \subset S_f$와 (보통의) 소 아이디얼
$\mathfrak p'' \subset S_{(f)}$에 대응한다면, 이들은 각각 $S_0$에서 같은
역상을 갖는다.
\end{proof}

\begin{lemma}
\label{constructions-lemma-proj-valuative-criterion}
$S$를 등급 환이라 하자. $S$가 $S_0$ 위의 대수로서 유한 생성이면 사상
$\text{Proj}(S) \to \Spec(S_0)$는 값매김 판정법의 존재성과 유일성 부분을
만족한다. Schemes, Definition \ref{schemes-definition-valuative-criterion}을
보라.
\end{lemma}

\begin{proof}
유일성 부분은 $\text{Proj}(S)$가 분리라는 사실에서 따른다(Lemma
\ref{constructions-lemma-proj-separated} 및 Schemes, Lemma
\ref{schemes-lemma-separated-implies-valuative}). 동차원소
$x_i \in S_{+}$, $i = 1, \ldots, n$을 택하자. 이들은 $S$를 $S_0$ 위에서
생성하도록 택한다.
$d_i = \deg(x_i)$라 하고 $d = \text{lcm}\{d_i\}$로 두자. Schemes,
Definition \ref{schemes-definition-valuative-criterion}에서와 같은 도식
$$
\xymatrix{
\Spec(K) \ar[r] \ar[d] & \text{Proj}(S) \ar[d] \\
\Spec(A) \ar[r] & \Spec(S_0)
}
$$
이 주어졌다고 하자. $v : K^* \to \Gamma$를 $A$의 값매김이라 나타내자.
Algebra, Definition \ref{algebra-definition-value-group}을 보라. 양의 차수를 갖는
동차원소 $f \in S_{+}$를 $\Spec(K)$의 상이 $D_{+}(f)$에 들어가도록 택할 수
있다. 그러면 환 사상들의 가환도식
$$
\xymatrix{
K & S_{(f)} \ar[l]^{\varphi} \\
A \ar[u] & S_0 \ar[l] \ar[u]
}
$$
을 얻는다. 번호를 바꾸어, $\varphi(x_i^{\deg(f)}/f^{d_i})$가
$i = 1, \ldots, r$에 대해서는 영이 아니고
$i = r + 1, \ldots, n$에 대해서는 영이라고 가정할 수 있다.
열린집합 $D_{+}(x_i)$들이 $\text{Proj}(S)$를 덮으므로 $r \geq 1$임을 알 수
있다. $i_0 \in \{1, \ldots, r\}$을
$\gamma_i = (d/d_i)v(\varphi(x_i^{\deg(f)}/f^{d_i}))$를 $\Gamma$에서 최소로 만드는
지표라 하자. 편의를 위해 $x_0 = x_{i_0}$ 및 $d_0 = d_{i_0}$로 둔다.
환의 사상 $\varphi$는 사상 $\varphi' : S_{(fx_0)} \to K$를 통해 분해되고,
이는 환의 사상 $S_{(x_0)} \to S_{(fx_0)} \to K$를 준다.
대수 $S_{(x_0)}$는 $S_0$ 위에서 원소
$x_1^{e_1} \ldots x_n^{e_n}/x_0^{e_0}$들로 생성된다. 여기서
$\sum e_i d_i = e_0 d_0$이다. $e_i > 0$인 어떤 $i > r$가 존재하면
$\varphi'(x_1^{e_1} \ldots x_n^{e_n}/x_0^{e_0}) = 0$이다.
$e_i = 0$이 $i > r$에 대해 성립하면
\begin{align*}
d \deg(f) v(\varphi'(x_1^{e_1} \ldots x_r^{e_r}/x_0^{e_0}))
& =
d v(\varphi'(x_1^{e_1 \deg(f)} \ldots x_r^{e_r \deg(f)}/x_0^{e_0 \deg(f)})) \\
& =
d \sum e_i v(\varphi'(x_i^{\deg(f)}/f^{d_i}))
- e_0 v(\varphi'(x_0^{\deg(f)}/f^{d_0})) \\
& =
\sum e_i d_i \gamma_i - e_0 d_0 \gamma_0 \\
& \geq
\sum e_i d_i \gamma_0 - e_0 d_0 \gamma_0 = 0
\end{align*}
이다. 이는 $\gamma_0$가 $\gamma_i$들 가운데 최소이기 때문이다. 따라서
$S_{(x_0)}$는 $A$ 안으로 $\varphi'$를 통해 사상된다. 이에 대응하는 스킴의
사상
$\Spec(A) \to \Spec(S_{(x_0)}) = D_{+}(x_0)
\subset \text{Proj}(S)$
은 이 증명의 첫 번째 가환도식에 맞는 사상을 제공한다.
\end{proof}

\noindent
Lemma \ref{constructions-lemma-proj-valuative-criterion}의 증명에서 그 보조정리의 가정 아래
사상 $\text{Proj}(S) \to \Spec(S_0)$가 준콤팩트이기도 함을 보았다. 따라서
Schemes, Proposition
\ref{schemes-proposition-characterize-universally-closed}에 의해 보조정리의
상황에서 $\text{Proj}(S) \to \Spec(S_0)$는 보편 닫힌이다. 다음 여러 예는
등급 환 $S$에 어떤 가정도 두지 않으면 이 결과들이 성립하지 않음을 보여 준다.

\begin{example}
\label{constructions-example-not-existence-valuative-big-proj}
$\mathbf{C}[X_1, X_2, X_3, \ldots]$를 등급 $\mathbf{C}$-대수로 생각하되,
각 $X_i$의 차수를 $1$로 둔다. $X_i$를 $t^{1/i}$로 보내는
환의 사상
$$
\mathbf{C}[X_1, X_2, X_3, \ldots]
\longrightarrow
\mathbf{C}[t^\alpha ; \alpha \in \mathbf{Q}_{\geq 0}]
$$
을 생각하자. 오른쪽 환을 국소화하면 값매김환 $A$가 되며, 국소화하는
아이디얼은 $\mathfrak m = (t^\alpha ; \alpha > 0)$이다. $K$를 $A$의 분수체라 하자. 위 사상은
$\Spec(K) \to \text{Proj}(\mathbf{C}[X_1, X_2, X_3, \ldots])$를 주지만,
이는 $\Spec(A)$ 전체에서 정의된 사상으로 확장되지 않는다. 그 까닭은
$\Spec(A)$의 상이 어느 $D_{+}(X_i)$ 안에 포함되어야 하는데, 그러면
$X_{i + 1}/X_i$가 $A$의 원소로 가야 하지만 실제로는
$t^{1/(i + 1) - 1/i}$로 가므로 그렇지 않기 때문이다.
\end{example}

\begin{example}
\label{constructions-example-not-existence-valuative-small-proj}
$R = \mathbf{C}[t]$라 하고
$$
S = R[X_1, X_2, X_3, \ldots]/(X_i^2 - tX_{i + 1}).
$$
로 두자. 등급화는 $R = S_0$이고 $\deg(X_i) = 2^{i - 1}$이 되게 한다.
$\mathfrak p \in \text{Proj}(S)$이면 $t \not \in \mathfrak p$임에 유의하자.
그렇지 않으면 $\mathfrak p$가 모든 $X_i$를 포함해야 하는데, 동차
스펙트럼의 원소에는 이것이 허용되지 않기 때문이다. 따라서
$D_{+}(X_i) = D_{+}(X_{i + 1})$가 모든 $i$에 대해 성립함을 알 수 있다. 그러므로
$\text{Proj}(S)$는 준콤팩트이고, 실제로 $D_{+}(X_1)$과 같으므로 아핀이다.
$\text{Proj}(S) \to \Spec(R)$의 상이 $D(t)$임을 쉽게 알 수 있다. 따라서
사상 $\text{Proj}(S) \to \Spec(R)$는 닫힌 사상이 아니다. 그러므로
값매김 판정법을 적용할 수 없다. 적용할 수 있다면 이 사상이 닫힌 사상임을
함의하기 때문이다(Schemes, Proposition
\ref{schemes-proposition-characterize-universally-closed}를 보라).
\end{example}

\begin{example}
\label{constructions-example-trivial-proj}
$A$를 환이라 하자. $S = A[T]$를 등급 $A$-대수로서 $T$의 차수가 $1$이게 하자.
그러면 표준 사상 $\text{Proj}(S) \to \Spec(A)$는 동형이다(Lemma
\ref{constructions-lemma-structure-morphism-proj}를 보라).
\end{example}

\begin{example}
\label{constructions-example-open-subset-proj}
$X = \Spec(A)$를 아핀 스킴이라 하고 $U \subset X$를 열린 부분스킴이라 하자.
$A[T]$를 $\deg T = 1$로 두어 등급화하자. $S$를 $A[T]$의 부분환으로서
$A$와 모든 $fT^i$로 생성된 것이라 정의하자. 여기서 $i \ge 0$이고
$f \in A$는 $D(f) \subset U$를 만족한다. $S$는 등급 환이고
$S_0 = A$이며 $\text{Proj}(S) \cong U$이고, 이 동형은 Lemma
\ref{constructions-lemma-structure-morphism-proj}의 표준 사상
$\text{Proj}(S) \to \Spec(A)$를 포함 $U \subset X$와 동일시한다고
주장한다.

\medskip\noindent
$\mathfrak p \in \text{Proj}(S)$가 모든 $fT \in S_1$을 $\mathfrak p$에
포함한다고 하자. 그러면 모든 생성원 $fT^i$ 중 $i \ge 1$인 것이
$\mathfrak p$에 속한다. 실제로
$(fT^i)^2 = (fT)(fT^{2i-1}) \in \mathfrak p$이고 $\mathfrak p$는 근기
아이디얼이기 때문이다. 그러나 그러면 $\mathfrak p \supset S_+$이며, 이는
불가능하다. 따라서 $\text{Proj}(S)$는 표준 아핀 열린 부분집합
$\{D_+(fT)\}_{fT \in S_1}$들로 덮인다.

\medskip\noindent
$fT \in S_1$이면 포함 $S \subset A[T]$가 $S[(fT)^{-1}]$과
$A[T, T^{-1}, f^{-1}]$ 사이의 등급 동형을 유도함에 유의하자. 따라서 표준
열린 부분집합 $D_+(fT) \cong \Spec(S_{(fT)})$는
$\Spec(A[T, T^{-1}, f^{-1}]_0) = \Spec(A[f^{-1}])$와 동형이다. 이 동형이
표준 사상 $\text{Proj}(S) \to \Spec(A)$의 제한임은 분명하다. 또한
$gT \in S_1$이면
$S[(fT)^{-1}, (gT)^{-1}] \cong A[T, T^{-1}, f^{-1}, g^{-1}]$이 등급 환의
동형이므로
$D_+(fT) \cap D_+(gT) \cong \Spec(A[f^{-1}, g^{-1}])$이다. 그러므로
$\text{Proj}(S)$는 열린 부분스킴 $D_+(fT)$들의 합집합이고, 이들은 표준
사상 아래 열린 부분스킴 $D(f) \subset X$와 동형이며, 이 열린 부분스킴들은
$\text{Proj}(S)$ 안에서 $X$ 안에서와 같은 방식으로 교차한다. 따라서 표준
사상은 $\text{Proj}(S)$에서 모든 $D(f) \subset U$의 합집합으로 가는
동형이며, 그 합집합은 $U$이다.
\end{example}

\section{Proj 위의 준연접 층}
\label{constructions-section-quasi-coherent-proj}

\noindent
$S$를 등급 환이라 하자. $M$을 등급 $S$-가군이라 하자. Lemma
\ref{constructions-lemma-proj-sheaves}에서 준연접 가군의 층 $\widetilde{M}$, 즉
$\text{Proj}(S)$ 위의 층과 사상
\begin{equation}
\label{constructions-equation-map-global-sections}
M_0 \longrightarrow \Gamma(\text{Proj}(S), \widetilde{M})
\end{equation}
을 구성하는 방법을 보았다. 이 사상은 차수 $0$ 부분, 즉 $M$의 그 부분에서
$\widetilde{M}$의 전역 단면으로 간다. 차수 $0$ 부분은 $n$번째 꼬임
$M(n)$을 등급 가군 $M$에 대해 취했을 때(Algebra, Section
\ref{algebra-section-graded}를 보라)
$M_n$과 같다. 따라서 사상
\begin{equation}
\label{constructions-equation-map-global-sections-degree-n}
M_n \longrightarrow \Gamma(\text{Proj}(S), \widetilde{M(n)}).
\end{equation}
을 얻을 수 있다. 임의의 준연접 층 $\mathcal{F}$, 즉 $\text{Proj}(S)$ 위의 층에
대해서도 이 연산을 수행하고 싶다. 이를 위해 구조층의 $n$번째 꼬임과
텐서곱을 취할 것이다. Definition \ref{constructions-definition-twist}를 보라. 두 개념을
관련짓기 위해 다음 보조정리를 사용한다.

\begin{lemma}
\label{constructions-lemma-widetilde-tensor}
$S$를 등급 환이라 하자. $(X, \mathcal{O}_X) = (\text{Proj}(S), \mathcal{O}_{\text{Proj}(S)})$를 Lemma
\ref{constructions-lemma-proj-scheme}의 스킴이라 하자. $f \in S_{+}$를 동차원소라 하자.
$x \in X$를 동차 소 아이디얼 $\mathfrak p \subset S$에 대응하는 점이라
하자. $M$, $N$을 등급 $S$-가군이라 하자. 표준
$\mathcal{O}_{\text{Proj}(S)}$-가군 사상
$$
\widetilde M \otimes_{\mathcal{O}_X} \widetilde N
\longrightarrow
\widetilde{M \otimes_S N}
$$
이 존재하며, 이는 표준 사상
$
M_{(f)} \otimes_{S_{(f)}} N_{(f)}
\to
(M \otimes_S N)_{(f)}
$
을 $D_{+}(f)$ 위의 단면에서 유도하고 표준 사상
$
M_{(\mathfrak p)} \otimes_{S_{(\mathfrak p)}} N_{(\mathfrak p)}
\to
(M \otimes_S N)_{(\mathfrak p)}
$
을 $x$에서의 줄기에서 유도한다. 또한 도식
$$
\xymatrix{
M_0 \otimes_{S_0} N_0 \ar[r] \ar[d] &
(M \otimes_S N)_0 \ar[d] \\
\Gamma(X, \widetilde M \otimes_{\mathcal{O}_X} \widetilde N) \ar[r] &
\Gamma(X, \widetilde{M \otimes_S N})
}
$$
은 가환하며, 여기서 수직 사상들은
(\ref{constructions-equation-map-global-sections})로 주어진다.
\end{lemma}

\begin{proof}
표시된 것과 같은 사상을 구성하는 것은 $\mathcal{O}_X$-쌍선형 사상
$$
\widetilde M \times \widetilde N
\longrightarrow
\widetilde{M \otimes_S N}
$$
을 구성하는 것과 같다. Modules, Section
\ref{modules-section-tensor-product}을 보라. 제한 사상들과 호환되도록 열린집합
$D_{+}(f)$ 위의 단면에서 이를 정의하면 충분하다. $D_{+}(f)$ 위에서는
$S_{(f)}$-쌍선형 사상
$M_{(f)} \times N_{(f)} \to (M \otimes_S N)_{(f)}$,
$(x/f^n, y/f^m) \mapsto (x \otimes y)/f^{n + m}$을 사용한다. 세부사항은
생략한다.
\end{proof}

\begin{remark}
\label{constructions-remark-not-isomorphism}
일반적으로 위 Lemma \ref{constructions-lemma-widetilde-tensor}에서 구성한 사상은 동형이
아니다. 다음은 한 예이다. $k$를 체라 하자. $S = k[x, y, z]$라 하고,
$k$는 차수 $0$에 두며 $\deg(x) = 1$, $\deg(y) = 2$,
$\deg(z) = 3$으로 둔다. $M = S(1)$ 및 $N = S(2)$라 하자. 표기법은
Algebra, Section \ref{algebra-section-graded}를 보라. 그러면
$M \otimes_S N = S(3)$이다. 다음에 유의하자.
\begin{eqnarray*}
S_z
& = &
k[x, y, z, 1/z] \\
S_{(z)}
& = &
k[x^3/z, xy/z, y^3/z^2]
\cong
k[u, v, w]/(uw - v^3) \\
M_{(z)} & = & S_{(z)} \cdot x + S_{(z)} \cdot y^2/z \subset S_z \\
N_{(z)} & = & S_{(z)} \cdot y + S_{(z)} \cdot x^2 \subset S_z \\
S(3)_{(z)} & = & S_{(z)} \cdot z \subset S_z
\end{eqnarray*}
극대 아이디얼 $\mathfrak m = (u, v, w) \subset S_{(z)}$를 생각하자.
$M_{(z)}/\mathfrak mM_{(z)}$와 $N_{(z)}/\mathfrak mN_{(z)}$는 모두
차원 $2$이고, 이는 $\kappa(\mathfrak m)$ 위에서의 차원임을 어렵지 않게
알 수 있다. 그러나
$S(3)_{(z)}/\mathfrak mS(3)_{(z)}$의 차원은 $1$이다. 따라서 사상
$M_{(z)} \otimes N_{(z)} \to S(3)_{(z)}$는 동형이 아니다.
\end{remark}

\section{Proj 위의 가역층}
\label{constructions-section-invertible-on-proj}

\noindent
등급 환 위의 등급 가군에 결부한 꼬인 가군 $M(n)$의 구성을 Algebra,
Section \ref{algebra-section-graded}에서 상기하자.

\begin{definition}
\label{constructions-definition-twist}
$S$를 등급 환이라 하자. $X = \text{Proj}(S)$라 하자.
\begin{enumerate}
\item $\mathcal{O}_X(n) = \widetilde{S(n)}$로 정의한다. 이를
{\it $n$번째 구조층의 꼬임($\text{Proj}(S)$의)}이라고 부른다.
\item 임의의 $\mathcal{O}_X$-가군의 층 $\mathcal{F}$에 대해
$\mathcal{F}(n) = \mathcal{F} \otimes_{\mathcal{O}_X} \mathcal{O}_X(n)$로
둔다.
\end{enumerate}
\end{definition}

\noindent
Lemma \ref{constructions-lemma-widetilde-tensor}을 사용하여 몇 가지 표준 사상을 구성한다.
$S(n) \otimes_S S(m) = S(n + m)$이므로 표준 사상
\begin{equation}
\label{constructions-equation-multiply}
\mathcal{O}_X(n) \otimes_{\mathcal{O}_X} \mathcal{O}_X(m)
\longrightarrow
\mathcal{O}_X(n + m).
\end{equation}
들이 존재한다. 일반적으로 이 사상들은 동형이 아니다. Remark
\ref{constructions-remark-not-isomorphism}의 예를 보라. 같은 예는 일반적으로
$\mathcal{O}_X(n)$이 $X$ 위의 가역층이 {\it 아니다}라는 것도 보여 준다.
임의의 $\mathcal{O}_X$-가군 $\mathcal{F}$와 텐서곱을 취하면 사상
\begin{equation}
\label{constructions-equation-multiply-on-sheaf}
\mathcal{O}_X(n) \otimes_{\mathcal{O}_X} \mathcal{F}(m)
\longrightarrow
\mathcal{F}(n + m).
\end{equation}
을 얻는다. 전역 단면 위의 사상
(\ref{constructions-equation-map-global-sections-degree-n})들은 등급 환의 사상
\begin{equation}
\label{constructions-equation-global-sections}
S \longrightarrow \bigoplus\nolimits_{n \geq 0} \Gamma(X, \mathcal{O}_X(n)).
\end{equation}
을 준다. 임의의 $\mathcal{O}_X$-가군 $\mathcal{F}$에 대해서도 사상
(\ref{constructions-equation-multiply-on-sheaf})들은 등급 가군 구조
\begin{equation}
\label{constructions-equation-global-sections-module}
\bigoplus\nolimits_{n \geq 0} \Gamma(X, \mathcal{O}_X(n))
\times
\bigoplus\nolimits_{m \in \mathbf{Z}} \Gamma(X, \mathcal{F}(m))
\longrightarrow
\bigoplus\nolimits_{m \in \mathbf{Z}} \Gamma(X, \mathcal{F}(m))
\end{equation}
을 주고, (\ref{constructions-equation-global-sections})을 통해 $S$-가군 구조도 준다.
더 일반적으로 임의의 등급 $S$-가군 $M$이 주어지면
$M(n) = M \otimes_S S(n)$이다. 따라서 사상
\begin{equation}
\label{constructions-equation-multiply-more-generally}
\widetilde M(n)
=
\widetilde M
\otimes_{\mathcal{O}_X}
\mathcal{O}_X(n)
\longrightarrow
\widetilde{M(n)}.
\end{equation}
을 얻는다. 전역 단면 위에서
(\ref{constructions-equation-map-global-sections-degree-n})은 등급 $S$-가군의 사상
\begin{equation}
\label{constructions-equation-global-sections-more-generally}
M \longrightarrow
\bigoplus\nolimits_{n \in \mathbf{Z}} \Gamma(X, \widetilde{M(n)}).
\end{equation}
을 정의한다. 다음은 정의에서 거의 곧바로 따르는 중요한 사실이다.

\begin{lemma}
\label{constructions-lemma-when-invertible}
$S$를 등급 환이라 하자. $X = \text{Proj}(S)$로 두자. $f \in S$를 차수
$d > 0$인 동차원소라 하자. 층 $\mathcal{O}_X(nd)|_{D_{+}(f)}$은 가역이고,
실제로 모든 $n \in \mathbf{Z}$에 대해 자명하다(Modules,
Definition \ref{modules-definition-invertible}를 보라). $D_{+}(f)$에 제한한
(\ref{constructions-equation-multiply})의 사상
$$
\mathcal{O}_X(nd)|_{D_{+}(f)} \otimes_{\mathcal{O}_{D_{+}(f)}}
\mathcal{O}_X(m)|_{D_{+}(f)}
\longrightarrow
\mathcal{O}_X(nd + m)|_{D_{+}(f)},
$$
$D_+(f)$에 제한한 (\ref{constructions-equation-multiply-on-sheaf})의 사상
$$
\mathcal{O}_X(nd)|_{D_{+}(f)} \otimes_{\mathcal{O}_{D_{+}(f)}}
\mathcal{F}(m)|_{D_{+}(f)}
\longrightarrow
\mathcal{F}(nd + m)|_{D_{+}(f)},
$$
그리고 $D_{+}(f)$에 제한한 (\ref{constructions-equation-multiply-more-generally})의 사상
$$
\widetilde M(nd)|_{D_{+}(f)}
=
\widetilde M|_{D_{+}(f)}
\otimes_{\mathcal{O}_{D_{+}(f)}}
\mathcal{O}_X(nd)|_{D_{+}(f)}
\longrightarrow
\widetilde{M(nd)}|_{D_{+}(f)}
$$
은 모든 $n, m \in \mathbf{Z}$에 대해 동형이다.
\end{lemma}

\begin{proof}
등급 $S$-가군 사상 $S \to S(nd)$ 및 $M \to M(nd)$은
$x \mapsto f^n x$로 주어지며, $f$를 가역화한 뒤 동형이 된다. 첫 번째 사상은
$S_{(f)} \cong S(nd)_{(f)}$임을 보여 주며, 이는 동형
$\mathcal{O}_{D_{+}(f)} \cong \mathcal{O}_X(nd)|_{D_{+}(f)}$을 준다.
두 번째 사상은 사상
$S(nd)_{(f)} \otimes_{S_{(f)}} M_{(f)} \to M(nd)_{(f)}$이 동형임을
보여 준다. 사상 (\ref{constructions-equation-multiply-on-sheaf})의 경우는 사상
(\ref{constructions-equation-multiply})의 경우에서 따른다.
\end{proof}

\begin{lemma}
\label{constructions-lemma-apply-modules}
$S$를 등급 환이라 하자. $M$을 등급 $S$-가군이라 하자.
$X = \text{Proj}(S)$로 두자. $X$가 표준 열린집합 $D_+(f)$들로 덮이고
$f \in S_1$이라고 가정하자. 예를 들어 $S$가 $S_1$로 $S_0$ 위에서
생성되면 그러하다. 그러면 층 $\mathcal{O}_X(n)$들은 가역이고 사상
(\ref{constructions-equation-multiply}), (\ref{constructions-equation-multiply-on-sheaf}), 그리고
(\ref{constructions-equation-multiply-more-generally})은 동형이다. 특히 이 사상들은 동형
$$
\mathcal{O}_X(1)^{\otimes n} \cong
\mathcal{O}_X(n)
\quad
\text{and}
\quad
\widetilde{M} \otimes_{\mathcal{O}_X} \mathcal{O}_X(n) =
\widetilde{M}(n) \cong \widetilde{M(n)}
$$
을 유도한다. 따라서 (\ref{constructions-equation-map-global-sections-degree-n})은
사상
\begin{equation}
\label{constructions-equation-map-global-sections-degree-n-simplified}
M_n \longrightarrow \Gamma(X, \widetilde{M}(n))
\end{equation}
이 되고, (\ref{constructions-equation-global-sections-more-generally})은 사상
\begin{equation}
\label{constructions-equation-global-sections-more-generally-simplified}
M \longrightarrow
\bigoplus\nolimits_{n \in \mathbf{Z}} \Gamma(X, \widetilde{M}(n)).
\end{equation}
이 된다.
\end{lemma}

\begin{proof}
보조정리의 가정 아래 $X$는 열린 부분집합 $D_{+}(f)$들로 덮이며
$f \in S_1$이다. 따라서 이 보조정리는 위 Lemma
\ref{constructions-lemma-when-invertible}의
결과이다.
\end{proof}

\begin{lemma}
\label{constructions-lemma-where-invertible}
$S$를 등급 환이라 하자. $X = \text{Proj}(S)$로 두자. 정수
$d \geq 1$을 고정하자. 다음 $X$의 열린 부분집합들은 서로 같다.
\begin{enumerate}
\item 가장 큰 열린 부분집합 $W = W_d \subset X$로서, 각
$\mathcal{O}_X(dn)|_W$가 가역이고 모든 곱셈 사상
$\mathcal{O}_X(nd)|_W \otimes_{\mathcal{O}_W} \mathcal{O}_X(md)|_W
\to \mathcal{O}_X(nd + md)|_W$
이 동형이 되는 것. (\ref{constructions-equation-multiply})를 보라.
\item 열린 부분집합 $D_{+}(fg)$들의 합집합. 여기서 $f, g \in S$는
동차원소이고 $\deg(f) = \deg(g) + d$이다.
\end{enumerate}
또한 모든 사상
$\widetilde M(nd)|_W = \widetilde M|_W \otimes_{\mathcal{O}_W}
\mathcal{O}_X(nd)|_W \to \widetilde{M(nd)}|_W$
은 동형이다. (\ref{constructions-equation-multiply-more-generally})를 보라.
\end{lemma}

\begin{proof}
$x \in D_{+}(fg)$이고 $\deg(f) = \deg(g) + d$이면 $D_{+}(fg)$ 위에서 층
$\mathcal{O}_X(dn)$은 원소 $(f/g)^n = f^{2n}/(fg)^n$로 생성된다. Lemma
\ref{constructions-lemma-when-invertible}의 증명에서와 같이 논증하면, 이는 $x$가 (1)에서
정의한 열린 부분집합 $W$에 속함을 뜻한다.

\medskip\noindent
반대로 $\mathcal{O}_X(dn)$이 모든 $n$에 대해 어떤 열린 근방 $V$에서
계수 1의 자유 가군이고, 이 근방은 $x \in X$의 근방이며, 모든 곱셈 사상
$\mathcal{O}_X(nd)|_V \otimes_{\mathcal{O}_V} \mathcal{O}_X(md)|_V
\to \mathcal{O}_X(nd + md)|_V$가 동형이라고 하자. 양의 차수를 갖는 동차원소
$h \in S_{+}$를 $x \in D_{+}(h) \subset V$가 되게 택할 수 있다.
구조층의 꼬임의 정의로부터 원소 $s$가 $(S_h)_d$에 존재하여, $s^n$이
$(S_h)_{nd}$의 $S_{(h)}$ 위 기저가 되고 이는 모든 $n \in \mathbf{Z}$에
대해 성립함을 얻는다. $s = f/h^m$으로 쓸 수 있는데, 여기서 어떤
$m \geq 1$ 및 $f \in S_{d + m \deg(h)}$가 존재한다.
$g = h^m$로 두면 $s = f/g$이다. 구성에 의해 $x \in D_{+}(g)$임에
유의하자. 또한 $g^d \in (S_h)_{d\deg(g)}$임에 유의하자. 가정에 의해 이를
$s^{\deg(g)} = f^{\deg(g)}/g^{\deg(g)}$의 배수, 즉
$g^d = a/g^e \cdot f^{\deg(g)}/g^{\deg(g)}$로 쓸 수 있다. 따라서
$g^{d + e + \deg(g)} = a f^{\deg(g)}$이고 $x \in D_{+}(f)$이기도 하다.
그러므로 $x$는 (2)에서 정의한 집합의 원소이다.

\medskip\noindent
생성 단면 $s = f/g$가 아핀 열린집합 $D_{+}(fg)$ 위에 존재하고, 그
거듭제곱들이 가군의 층 $\mathcal{O}_X(nd)$을 자유롭게 생성한다. 이 사실은
곱셈 사상
$\widetilde M(nd)|_W = \widetilde M|_W \otimes_{\mathcal{O}_W}
\mathcal{O}_X(nd)|_W \to \widetilde{M(nd)}|_W$
이 동형임을 쉽게 함의한다. (\ref{constructions-equation-multiply-more-generally})를 보라.
Lemma \ref{constructions-lemma-when-invertible}의 증명과 비교하라.
\end{proof}

\noindent
가역층 $\mathcal{L}$이 국소 환 달린 공간 $X$ 위에 주어지고, 전역 단면
$s$가 $\mathcal{L}$에 주어졌을 때 집합
$X_s = \{x \in X \mid s \not \in \mathfrak m_x\mathcal{L}_x\}$가 열린집합임을
Modules, Lemma \ref{modules-lemma-s-open}에서 상기하자.

\begin{lemma}
\label{constructions-lemma-principal-open}
$S$를 등급 환이라 하자. $X = \text{Proj}(S)$로 두자. 정수 $d \geq 1$을
고정하자. $W = W_d \subset X$를 Lemma \ref{constructions-lemma-where-invertible}에서
정의한 열린 부분스킴이라 하자. $n \geq 1$ 및 $f \in S_{nd}$라 하자.
$s \in \Gamma(W, \mathcal{O}_W(nd))$를 $f$의 상인 단면이라 나타내자. 이 상은
(\ref{constructions-equation-global-sections})을 $W$로 제한하여 얻는다. 그러면
$$
W_s = D_{+}(f) \cap W.
$$
이다.
\end{lemma}

\begin{proof}
$D_{+}(ab) \subset W$를 표준 아핀 열린집합이라 하자. 여기서
$a, b \in S$는 동차원소이고 $\deg(a) = \deg(b) + d$이다.
$D_{+}(ab) \cap D_{+}(f) = D_{+}(abf)$임에 유의하자. 한편 $s$를
$D_{+}(ab)$에 제한한 것은 원소
$f/1 = b^nf/a^n (a/b)^n \in (S_{ab})_{nd}$에 대응한다. Lemma
\ref{constructions-lemma-where-invertible}의 증명에서 $(a/b)^n$이
$\mathcal{O}_W(nd)$의 생성원임을 보았으며, 이는 $D_{+}(ab)$ 위에서의
생성원이다. 따라서
$W_s \cap D_{+}(ab)$는 $b^nf/a^n \in \mathcal{O}_X(D_{+}(ab))$에 결부한
주 열린집합이다. 이로부터 보조정리의 결과가 분명하다.
\end{proof}

\noindent
다음 보조정리는 나중에 풍부한 가역층을 지닌 스킴을 특징짓는 데 사용할 성질들을
서술한다.

\begin{lemma}
\label{constructions-lemma-ample-on-proj}
$S$를 등급 환이라 하자. $X = \text{Proj}(S)$라 하자.
$Y \subset X$를 준콤팩트 열린 부분스킴이라 하자. $\mathcal{O}_Y(n)$을
$\mathcal{O}_X(n)$을 $Y$에 제한한 것으로 나타내자. 다음을 만족하는 정수
$d \geq 1$이 존재한다.
\begin{enumerate}
\item 부분스킴 $Y$는 Lemma \ref{constructions-lemma-where-invertible}에서 정의한 열린집합
$W_d$에 포함된다.
\item 층 $\mathcal{O}_Y(dn)$은 모든 $n \in \mathbf{Z}$에 대해 가역이다.
\item Equation (\ref{constructions-equation-multiply})의 모든 사상
$\mathcal{O}_Y(nd) \otimes_{\mathcal{O}_Y} \mathcal{O}_Y(m)
\longrightarrow
\mathcal{O}_Y(nd + m)$
은 동형이다.
\item 모든 사상
$\widetilde M(nd)|_Y = \widetilde M|_Y \otimes_{\mathcal{O}_Y}
\mathcal{O}_X(nd)|_Y \to \widetilde{M(nd)}|_Y$
은 동형이다. (\ref{constructions-equation-multiply-more-generally})를 보라.
\item $f \in S_{nd}$가 주어졌을 때 $s \in \Gamma(Y, \mathcal{O}_Y(nd))$를
$f$의 상으로 나타내자. 이는 (\ref{constructions-equation-global-sections})을 $Y$로 제한하여
얻는다. 그러면 $D_{+}(f) \cap Y = Y_s$이다.
\item $Y$의 위상의 기저는 열린집합 $Y_s$들의 모음으로 주어지며, 여기서
$s \in \Gamma(Y, \mathcal{O}_Y(nd))$이고 $n \geq 1$이다.
\item $Y$의 위상의 기저는 아핀인 열린집합 $Y_s \subset Y$들로 주어지며,
여기서 $s \in \Gamma(Y, \mathcal{O}_Y(nd))$이고 $n \geq 1$이다.
\end{enumerate}
\end{lemma}

\begin{proof}
$Y$가 준콤팩트이므로 유한 개의 동차원소 $f_i \in S_{+}$,
$i = 1, \ldots, n$이 존재하여 표준 열린집합 $D_{+}(f_i)$들이 $Y$의 열린
덮개를 이룬다. $d_i = \deg(f_i)$라 하고 $d = d_1 \ldots d_n$으로 두자.
$D_{+}(f_i) = D_{+}(f_i^{d/d_i})$임에 유의하자. 따라서 Lemma
\ref{constructions-lemma-where-invertible}의 특징 (2)에 의해, 또는 Lemma
\ref{constructions-lemma-when-invertible}을 이용한 특징 (1)에 의해
$Y \subset W_d$임을 곧바로 알 수 있다. (1)은 Lemma
\ref{constructions-lemma-where-invertible}에 의해 (2), (3), (4)를 함의한다. ((3)은
(4)의 특수한 경우임에 유의하자.) (5)는 Lemma
\ref{constructions-lemma-principal-open}에서 따른다. (6)과 (7)은 열린 부분집합
$D_{+}(f)$들이 $X$의 위상의 기저를 이루고 아핀이기 때문에 따른다.
\end{proof}

\begin{lemma}
\label{constructions-lemma-comparison-proj-quasi-coherent}
$S$를 등급 환이라 하자. $X = \text{Proj}(S)$로 두자. $\mathcal{F}$를
준연접 $\mathcal{O}_X$-가군이라 하자.
$M = \bigoplus_{n \in \mathbf{Z}} \Gamma(X, \mathcal{F}(n))$로 두고, 이를
등급 $S$-가군으로 보되,
(\ref{constructions-equation-global-sections-module}) 및 (\ref{constructions-equation-global-sections})을
사용한다. 그러면 표준 $\mathcal{O}_X$-가군 사상
$$
\widetilde{M} \longrightarrow \mathcal{F}
$$
이 존재한다. 이는 $\mathcal{F}$에 대해 함자적이며, 유도된 사상
$M_0 \to \Gamma(X, \mathcal{F})$는 항등사상이다.
\end{lemma}

\begin{proof}
$f \in S$를 차수 $d > 0$인 동차원소라 하자. Lemma
\ref{constructions-lemma-proj-sheaves}에 의해 $\widetilde{M}|_{D_{+}(f)}$는
$S_{(f)}$-가군 $M_{(f)}$에 대응함을 상기하자. 따라서 표준 사상
$$
M_{(f)} \longrightarrow \Gamma(D_+(f), \mathcal{F}),\quad
m/f^n \longmapsto m|_{D_+(f)} \otimes f|_{D_+(f)}^{-n}
$$
을 정의할 수 있다. 이는 $f|_{D_+(f)}$가 가역층
$\mathcal{O}_X(d)|_{D_+(f)}$의 자명화 단면이므로 의미가 있다. Lemma
\ref{constructions-lemma-when-invertible}과 그 증명을 보라. $\widetilde{M}$이 준연접이므로
Schemes, Lemma \ref{schemes-lemma-compare-constructions}을 통해 표준 사상
$$
\widetilde{M}|_{D_+(f)} \longrightarrow \mathcal{F}|_{D_+(f)}
$$
을 얻는다. 표시된 사상들이 겹침 위에서 붙는다는 것을 증명하면 전역 사상을
얻는다. 그 증명은 생략한다. 마지막 명제의 증명도 생략한다.
\end{proof}






\section{Proj의 함자성}
\label{constructions-section-functoriality-proj}

\noindent
등급 환의 사상 $\psi : A \to B$가 언제나 결부한 사영 동차 스펙트럼의 사상을
주는 것은 아니다. 그 이유는 역상 $\psi^{-1}(\mathfrak q)$이, 동차 소
아이디얼 $\mathfrak q \subset B$의 역상으로서, 무관 소 아이디얼 $A_{+}$를
포함할 수 있기 때문이다. 이는 $\mathfrak q$가 $B_{+}$를 포함하지 않을 때도
가능하다. 올바른 결과는 다음과 같다.

\begin{lemma}
\label{constructions-lemma-morphism-proj}
$A$, $B$를 두 등급 환이라 하자. $X = \text{Proj}(A)$ 및
$Y = \text{Proj}(B)$로 두자. $\psi : A \to B$를 등급 환의 사상이라 하자.
다음과 같이 두자.
$$
U(\psi)
=
\bigcup\nolimits_{f \in A_{+}\ \text{homogeneous}} D_{+}(\psi(f))
\subset Y.
$$
그러면 스킴의 표준 사상
$$
r_\psi :
U(\psi)
\longrightarrow
X
$$
과 $\mathbf{Z}$-등급 $\mathcal{O}_{U(\psi)}$-대수의 사상
$$
\theta = \theta_\psi :
r_\psi^*\left(
\bigoplus\nolimits_{d \in \mathbf{Z}} \mathcal{O}_X(d)
\right)
\longrightarrow
\bigoplus\nolimits_{d \in \mathbf{Z}} \mathcal{O}_{U(\psi)}(d).
$$
이 존재한다. 삼중쌍 $(U(\psi), r_\psi, \theta)$은 다음 성질들로
특징지어진다.
\begin{enumerate}
\item 모든 $d \geq 0$에 대해 도식
$$
\xymatrix{
A_d \ar[d] \ar[rr]_{\psi} & &
B_d \ar[d] \\
\Gamma(X, \mathcal{O}_X(d)) \ar[r]^-\theta &
\Gamma(U(\psi), \mathcal{O}_Y(d)) &
\Gamma(Y, \mathcal{O}_Y(d)) \ar[l]
}
$$
은 가환한다.
\item 임의의 동차원소 $f \in A_{+}$에 대해
$r_\psi^{-1}(D_{+}(f)) = D_{+}(\psi(f))$이고, $r_\psi$를
$D_{+}(\psi(f))$에 제한한 것은 환의 사상
$A_{(f)} \to B_{(\psi(f))}$에 대응하며, 이 사상은 $\psi$가 유도한다.
\end{enumerate}
\end{lemma}

\begin{proof}
조건 (2)가 스킴의 사상과 열린 부분집합 $U(\psi)$를 유일하게 결정함은
분명하다. $f \in A_d$를 택하되 $d \geq 1$이라 하자.
$\mathcal{O}_X(n)|_{D_{+}(f)}$은 $A_{(f)}$-가군 $(A_f)_n$에 대응하고,
$\mathcal{O}_Y(n)|_{D_{+}(\psi(f))}$은 $B_{(\psi(f))}$-가군
$(B_{\psi(f)})_n$에 대응함에 유의하자. 다시 말해 $\theta$를
$D_{+}(\psi(f))$에 제한하면 $\mathbf{Z}$-등급 $B_{(\psi(f))}$-대수의 사상
$$
A_f \otimes_{A_{(f)}} B_{(\psi(f))}
\longrightarrow
B_{\psi(f)}
$$
에 대응한다. 조건 (1)은 $A$의 모든 원소의 상을 결정한다. $f$는 가역원이고
$\psi(f)$로 사상되므로, $1/f^m$은 $1/\psi(f)^m$으로 사상된다. 이로부터
$\theta$가 유일하게 결정되고 다음 규칙으로 주어짐이 쉽게 따른다.
$$
a/f^m \otimes b/\psi(f)^e \longmapsto \psi(a)b/\psi(f)^{m + e}.
$$
존재성을 보이기 위해, 위의 유일성 증명이 사상 $r$과 사상 $\theta$를
꼴 $D_{+}(\psi(f)) \subset U(\psi)$인 모든 표준 열린집합에서
$D_{+}(f)$로 제한했을 때 잘 정의하는 처방도 주었음을 지적하자.
이들을 $r_f$와 $\theta_f$라 부르자. 따라서
$D_{+}(f) \subset D_{+}(g)$가 어떤 동차원소 $f, g \in A_{+}$에 대해
성립하면 $r_g$를
$D_{+}(\psi(f))$에 제한한 것이 $r_f$와 일치함만 확인하면 된다. 이는 위에서
$r$과 $\theta$에 대해 준 공식들로부터 분명하다.
\end{proof}


\begin{lemma}
\label{constructions-lemma-morphism-proj-transitive}
$A$, $B$, $C$를 등급 환이라 하자.
$X = \text{Proj}(A)$, $Y = \text{Proj}(B)$ 및 $Z = \text{Proj}(C)$로 두자.
$\varphi : A \to B$, $\psi : B \to C$를 등급 환의 사상이라 하자.
그러면
$$
U(\psi \circ \varphi) = r_\psi^{-1}(U(\varphi))
\quad
\text{and}
\quad
r_{\psi \circ \varphi}
=
r_\varphi \circ r_\psi|_{U(\psi \circ \varphi)}.
$$
또한 자명한 표기 아래
$$
\theta_\psi \circ r_\psi^*\theta_\varphi
=
\theta_{\psi \circ \varphi}
$$
가 성립한다.
\end{lemma}

\begin{proof}
생략한다.
\end{proof}

\begin{lemma}
\label{constructions-lemma-surjective-graded-rings-map-proj}
위의 Lemma \ref{constructions-lemma-morphism-proj}와 같은 가정과 표기를 사용하자.
$A_d \to B_d$가 모든 $d \gg 0$에 대해 전사라고 가정하자. 그러면
\begin{enumerate}
\item $U(\psi) = Y$이고,
\item $r_\psi : Y \to X$는 닫힌 몰입이며,
\item 사상 $\theta : r_\psi^*\mathcal{O}_X(n) \to \mathcal{O}_Y(n)$은
전사이지만 일반적으로 동형은 아니다($A \to B$가 전사인 경우에도 그러하다).
\end{enumerate}
\end{lemma}

\begin{proof}
(1)은 $U(\psi)$의 정의와 $D_{+}(f) = D_{+}(f^n)$가 임의의
$n > 0$에 대해 성립한다는 사실에서 따른다. 동차원소 $f \in A_{+}$에 대해
$A_{(f)} \to B_{(\psi(f))}$가 전사임을 알 수 있다. 실제로
$B_{(\psi(f))}$의 임의의 원소는 분수 $b/\psi(f)^n$으로 나타낼 수 있고,
이때 $n$을 얼마든지 크게 택할 수 있다(따라서 $b \in B$의 차수는 커질 수밖에
없다). 이로써 (2)를 증명하였다. 같은 논증은 사상
$$
A_f \to B_{\psi(f)}
$$
이 전사임을 보이며, 따라서 $\theta$의 전사성이 따른다.
이 사상이 동형이 아닌 예로, 등급 환 $A = k[x, y]$를 생각하자. 여기서
$k$는 체이고 $\deg(x) = 1$, $\deg(y) = 2$이다.
$I = (x)$로 두면 $B = k[y]$이다. 이 경우
$\mathcal{O}_Y(1) = 0$임에 유의하자. 그러나
$r_\psi^*\mathcal{O}_X(1)$이 영이 아님은 쉽게 알 수 있다.
(이보다 덜 인위적인 예들도 있다.)
\end{proof}

\begin{lemma}
\label{constructions-lemma-eventual-iso-graded-rings-map-proj}
위의 Lemma \ref{constructions-lemma-morphism-proj}와 같은 가정과 표기를 사용하자.
$A_d \to B_d$가 모든 $d \gg 0$에 대해 동형이라고 가정하자. 그러면
\begin{enumerate}
\item $U(\psi) = Y$이고,
\item $r_\psi : Y \to X$는 동형이며,
\item 사상 $\theta : r_\psi^*\mathcal{O}_X(n) \to \mathcal{O}_Y(n)$은
동형이다.
\end{enumerate}
\end{lemma}

\begin{proof}
Lemma \ref{constructions-lemma-surjective-graded-rings-map-proj}에 의해 (1)이 성립한다.
$f \in A_{+}$를 동차원소라 하자. $\psi$에 대한 가정은
$A_f \to B_f$가 동형임을 함의한다(세부사항은 생략한다). 따라서
$r_\psi$와 $\theta$를 $D_{+}(f)$ 위로 제한하면 동형이 됨은 분명하다.
보조정리가 따른다.
\end{proof}


\begin{lemma}
\label{constructions-lemma-surjective-graded-rings-generated-degree-1-map-proj}
위의 Lemma \ref{constructions-lemma-morphism-proj}와 같은 가정과 표기를 사용하자.
$A_d \to B_d$가 $d \gg 0$에 대해 전사이고 $A$가 $A_1$에 의해
$A_0$ 위에서 생성된다고 가정하자. 그러면
\begin{enumerate}
\item $U(\psi) = Y$이고,
\item $r_\psi : Y \to X$는 닫힌 몰입이며,
\item 사상 $\theta : r_\psi^*\mathcal{O}_X(n) \to \mathcal{O}_Y(n)$은
동형이다.
\end{enumerate}
\end{lemma}

\begin{proof}
Lemmas \ref{constructions-lemma-eventual-iso-graded-rings-map-proj} 및
\ref{constructions-lemma-morphism-proj-transitive}에 의해 $B$를 $A \to B$의 상으로
바꾸어도 $X$ 또는 층 $\mathcal{O}_X(n)$은 바뀌지 않는다.
따라서 $A \to B$가 전사라고 가정해도 된다. Lemma
\ref{constructions-lemma-surjective-graded-rings-map-proj}에 의해 (1), (2) 및 (3)의
전사성을 얻는다. Lemma \ref{constructions-lemma-apply-modules}에 의해
$\mathcal{O}_X(n)$과 $\mathcal{O}_Y(n)$은 모두 가역이다. 따라서
$\theta$는 동형이다.
\end{proof}

\begin{lemma}
\label{constructions-lemma-base-change-map-proj}
\begin{reference}
\cite[II, Proposition 2.8.10]{EGA}
\end{reference}
위의 Lemma \ref{constructions-lemma-morphism-proj}와 같은 가정과 표기를 사용하자.
환의 사상 $R \to A_0$과 환의 사상 $R \to R'$이 존재하여
$B = R' \otimes_R A$가 된다고 가정하자. 그러면
\begin{enumerate}
\item $U(\psi) = Y$이고,
\item 도식
$$
\xymatrix{
Y = \text{Proj}(B) \ar[r]_{r_\psi} \ar[d] &
\text{Proj}(A) = X \ar[d] \\
\Spec(R') \ar[r] &
\Spec(R)
}
$$
은 올곱 도식이며,
\item 사상 $\theta : r_\psi^*\mathcal{O}_X(n) \to \mathcal{O}_Y(n)$은
동형이다.
\end{enumerate}
\end{lemma}

\begin{proof}
표준 열린집합 $D_{+}(f)$ 위에서 일어나는 일을 살펴보면 곧바로 따른다.
여기서 $f \in A_{+}$이다.
\end{proof}

\begin{lemma}
\label{constructions-lemma-localization-map-proj}
위의 Lemma \ref{constructions-lemma-morphism-proj}와 같은 가정과 표기를 사용하자.
$g \in A_0$가 존재하여 $\psi$가 동형 $A_g \to B$를 유도한다고
가정하자. 그러면 $U(\psi) = Y$이고 $r_\psi : Y \to X$는 열린 몰입이며,
$Y$를 $D(g) \subset \Spec(A_0)$의 역상과 동형으로 만든다. 또한 사상
$\theta$는 동형이다.
\end{lemma}

\begin{proof}
이는 위의 Lemma \ref{constructions-lemma-base-change-map-proj}의 특수한 경우이다.
\end{proof}


\begin{lemma}
\label{constructions-lemma-d-uple}
$S$를 등급 환이라 하자. $d \geq 1$이라 하자. Algebra, Section
\ref{algebra-section-graded}의 표기를 사용하여 $S' = S^{(d)}$로 두자.
$X = \text{Proj}(S)$ 및 $X' = \text{Proj}(S')$로 두자. 다음을 만족하는
스킴의 표준 동형 $i : X \to X'$이 존재한다.
\begin{enumerate}
\item 임의의 등급 $S$-가군 $M$에 대해 $M' = M^{(d)}$로 두면 표준 동형
$\widetilde{M} \to i^*\widetilde{M'}$이 존재한다.
\item 표준 동형
$\mathcal{O}_{X}(nd) \to i^*\mathcal{O}_{X'}(n)$이 존재한다.
\end{enumerate}
이 동형들은 Lemma \ref{constructions-lemma-widetilde-tensor}의 곱셈 사상들과 호환되며,
따라서 사상
(\ref{constructions-equation-multiply}),
(\ref{constructions-equation-multiply-on-sheaf}),
(\ref{constructions-equation-global-sections}),
(\ref{constructions-equation-global-sections-module}),
(\ref{constructions-equation-multiply-more-generally}) 및
(\ref{constructions-equation-global-sections-more-generally})과도 호환된다
(정확한 진술은 증명을 보라).
\end{lemma}

\begin{proof}
단사인 환 사상 $S' \to S$(우리의 규약 아래 이는 등급 환의 준동형이 아니다)는
사상 $j : \Spec(S) \to \Spec(S')$를 유도한다. 등급 소 아이디얼
$\mathfrak p \subset S$가 주어지면
$\mathfrak p' = j(\mathfrak p) = S' \cap \mathfrak p$는 $S'$의 등급 소
아이디얼임을 알 수 있다. 또한 $f \in S_+$가 동차이고
$f \not \in \mathfrak p$이면 $f^d \in S'_+$이고
$f^d \not \in \mathfrak p'$이다. 역으로, 등급 소 아이디얼
$\mathfrak p' \subset S'$가 주어지고 이것이 어떤 동차원소
$f \in S'_+$를 포함하지 않으면,
$\mathfrak p = \{g \in S \mid g^d \in \mathfrak p'\}$는 $S$의 등급 소
아이디얼이고 $f$를 포함하지 않으며, $j$에 의한 그 상은 $\mathfrak p'$이다.
$\mathfrak p$가 아이디얼임을 보이기 위해 $g, h \in \mathfrak p$이면
이항 공식에 의해 $(g + h)^{2d} \in \mathfrak p'$이고, 따라서
$g + h \in \mathfrak p'$이다. 여기서 $\mathfrak p'$가 소 아이디얼임을
사용하였다. 이와 같이
$j$가 위상동형 $i : X \to X'$을 유도함을 알 수 있다. 더 나아가 동차원소
$f \in S_+$가 주어지면 $S_{(f)} \cong S'_{(f^d)}$이다. 이 동형들은 Lemma
\ref{constructions-lemma-standard-open}의 제한 사상들과 호환되므로, 동형
$i^\sharp : i^{-1}\mathcal{O}_{X'} \to \mathcal{O}_X$가 $X$와 $X'$ 위의
구조층 사이에 존재한다. 따라서 $i$는 스킴의 동형이다.

\medskip\noindent
$M$을 등급 $S$-가군이라 하자. 동차원소 $f \in S_+$가 주어지면
$M_{(f)} \cong M'_{(f^d)}$이므로, 위와 정확히 같은 방식으로 (1)의 동형을
얻는다. (2)의 동형들은 $M = S(nd)$인 경우의 (1)이며, 이때
$M' = S'(n)$이다. $M$과 $N$을 등급 $S$-가군이라 하자. 그러면
$$
M' \otimes_{S'} N' =
(M \otimes_S N)^{(d)} =
(M \otimes_S N)'
$$
이고, 이는 정의로부터 직접 확인할 수 있다. 이제 Lemma
\ref{constructions-lemma-widetilde-tensor}의 곱셈 사상과의 호환성은 도식
$$
\xymatrix{
\widetilde M \otimes_{\mathcal{O}_X} \widetilde N
\ar[d]_{(1) \otimes (1)} \ar[r] &
\widetilde{M \otimes_S N} \ar[d]^{(1)} \\
i^*(\widetilde{M'} \otimes_{\mathcal{O}_{X'}} \widetilde{N'}) \ar[r] &
i^*(\widetilde{M' \otimes_{S'} N'})
}
$$
의 가환성이다. 이는 열린집합 $D_+(f) = D_+(f^d)$ 위에서 사상들의 구성을
살펴보면 알 수 있다. 위쪽 수평 화살표는 사상
$M_{(f)} \times N_{(f)} \to (M \otimes_S N)_{(f)}$으로 주어지고,
아래쪽 수평 화살표는 사상
$M'_{(f^d)} \times N'_{(f^d)} \to (M' \otimes_{S'} N')_{(f^d)}$으로
주어진다. 이 사상들은 식별 $M_{(f)} = M'_{(f^d)}$ 등을 통해 일치하므로
원하는 호환성을 얻는다. 다른 호환성들의 증명은 생략한다.
\end{proof}


\section{Proj로의 사상}
\label{constructions-section-morphisms-proj}

\noindent
$S$를 등급 환이라 하자.
$X = \text{Proj}(S)$를 $S$의 동차 스펙트럼이라 하자.
$d \geq 1$을 정수라 하자. 열린 부분스킴
\begin{equation}
\label{constructions-equation-Ud}
U_d = \bigcup\nolimits_{f  \in S_d} D_{+}(f)
\quad\subset\quad
X = \text{Proj}(S)
\end{equation}
을 생각하자. $d | d' \Rightarrow U_d \subset U_{d'}$이고
$X = \bigcup_d U_d$임에 유의하자. $X$와 $U_d$ 어느 쪽도 준콤팩트일
필요가 없다. Algebra, Lemma \ref{algebra-lemma-topology-proj}를 보라.
$\mathcal{O}_{U_d}(n) = \mathcal{O}_X(n)|_{U_d}$로 쓰자. Lemma
\ref{constructions-lemma-when-invertible}에 의해 $\mathcal{O}_{U_d}(nd)$는
$n \in \mathbf{Z}$에 대해 가역 $\mathcal{O}_{U_d}$-가군이고,
(\ref{constructions-equation-multiply})의 모든 곱셈 사상
$\mathcal{O}_{U_d}(nd) \otimes_{\mathcal{O}_{U_d}} \mathcal{O}_{U_d}(m)
\to \mathcal{O}_{U_d}(nd + m)$은 동형임을 안다. 특히
$\mathcal{O}_{U_d}(nd) \cong \mathcal{O}_{U_d}(d)^{\otimes n}$이다.
전역 단면 위의 등급 환 사상 (\ref{constructions-equation-global-sections})과
$U_d$로의 제한을 합성하면 등급 환의 준동형
\begin{equation}
\label{constructions-equation-psi-d}
\psi^d : S^{(d)} \longrightarrow \Gamma_*(U_d, \mathcal{O}_{U_d}(d)).
\end{equation}
을 얻는다. 표기 $S^{(d)}$에 대해서는 Algebra, Section
\ref{algebra-section-graded}를 보라. 표기 $\Gamma_*$에 대해서는 Modules,
Definition \ref{modules-definition-gamma-star}를 보라. 또한 $U_d$는 열린집합
$D_{+}(f)$, $f \in S_d$로 덮이므로, $\mathcal{O}_{U_d}(d)$는
$\psi^d_1 : S^{(d)}_1 = S_d \to \Gamma(U_d, \mathcal{O}_{U_d}(d))$의
상에 속하는 단면들에 의해 전역 생성됨을 알 수 있다. Modules, Definition
\ref{modules-definition-globally-generated}를 보라.

\medskip\noindent
$Y$를 스킴이라 하고 $\varphi : Y \to X$를 스킴의 사상이라 하자.
상 $\varphi(Y)$가 열린 부분스킴 $U_d$에 포함된다고 가정하자. 이는
$X$의 열린 부분스킴이다.
Modules, Definition \ref{modules-definition-gamma-star} 뒤의 논의에 의해
등급 환의 준동형
$$
\Gamma_*(U_d, \mathcal{O}_{U_d}(d))
\longrightarrow
\Gamma_*(Y, \varphi^*\mathcal{O}_X(d)).
$$
을 얻는다. 이를 $\psi^d$와 합성하면 등급 환의 준동형
\begin{equation}
\label{constructions-equation-psi-phi-d}
\psi_\varphi^d :
S^{(d)}
\longrightarrow
\Gamma_*(Y, \varphi^*\mathcal{O}_X(d))
\end{equation}
을 얻으며, 이는 가역층 $\varphi^*\mathcal{O}_X(d)$가
$(S^{(d)})_1 = S_d \to \Gamma(Y, \varphi^*\mathcal{O}_X(d))$의 상에 속하는
단면들에 의해 전역 생성된다는 성질을 가진다.

\begin{lemma}
\label{constructions-lemma-converse-construction}
$S$를 등급 환이라 하고 $X = \text{Proj}(S)$라 하자. $d \geq 1$ 및 위와 같은
$U_d \subset X$를 택하자. $Y$를 스킴이라 하자. $\mathcal{L}$을 $Y$ 위의
가역층이라 하자. $\psi : S^{(d)} \to \Gamma_*(Y, \mathcal{L})$를 등급 환의
준동형이라 하고, $\mathcal{L}$이
$\psi|_{S_d} : S_d \to \Gamma(Y, \mathcal{L})$의 상에 속하는 단면들에
의해 생성된다고 가정하자. 그러면 사상 $\varphi : Y \to X$로서
$\varphi(Y) \subset U_d$를 만족하는 것과 동형
$\alpha : \varphi^*\mathcal{O}_{U_d}(d) \to \mathcal{L}$이 존재하여,
$\psi_\varphi^d$가 $\psi$와 $\alpha$를 통해 일치한다. 즉,
$$
\xymatrix{
\Gamma_*(Y, \mathcal{L}) &
\Gamma_*(Y, \varphi^*\mathcal{O}_{U_d}(d)) \ar[l]^-\alpha &
\Gamma_*(U_d, \mathcal{O}_{U_d}(d)) \ar[l]^-{\varphi^*} \\
S^{(d)} \ar[u]^\psi & &
S^{(d)} \ar[u]^{\psi^d} \ar[ul]^{\psi^d_\varphi} \ar[ll]_{\text{id}}
}
$$
은 가환한다. 또한 쌍 $(\varphi, \alpha)$는 유일하다.
\end{lemma}

\begin{proof}
$f \in S_d$를 택하자. $s = \psi(f) \in \Gamma(Y, \mathcal{L})$로 나타내자.
$Y_s$를 $s$가 영이 되지 않는 열린집합이라 하자. 이 위에서 $s$에 의한 곱셈은 동형
$\mathcal{O}_{Y_s} \to \mathcal{L}|_{Y_s}$을 유도한다. Modules, Lemma
\ref{modules-lemma-s-open}을 보라. 이 사상의 역을 $x \mapsto x/s$로
나타내고, $\mathcal{L}$의 거듭제곱에 대해서도 같은 표기를 사용하겠다.
이를 이용하여 환의 사상
$\psi_{(f)} : S_{(f)} \to \Gamma(Y_s, \mathcal{O}_Y)$을 정의하되, 분수
$a/f^n$을 $\psi(a)/s^n$으로 보낸다.
Schemes, Lemma \ref{schemes-lemma-morphism-into-affine}에 의해 이는 사상
$\varphi_f : Y_s \to \Spec(S_{(f)}) = D_{+}(f)$에 대응한다. 또한 동형
$\alpha_f : \varphi_f^*\mathcal{O}_{D_{+}(f)}(d) \to \mathcal{L}|_{Y_s}$을
도입한다. 이는 자명화 단면 $f$의 $D_{+}(f)$ 위의 당김을 자명화 단면
$s$로 보내며, 후자는 $Y_s$ 위의 단면이다. 이 선택 아래 보조정리의 도식은 $Y$를 $Y_s$로,
$\varphi$를 $\varphi_f$로, $\alpha$를 $\alpha_f$로 바꾸어도 가환한다.
확인은 생략한다.

\medskip\noindent
$f' \in S_d$가 두 번째 원소라고 하고
$s' = \psi(f') \in \Gamma(Y, \mathcal{L})$로 나타내자. 그러면
$Y_s \cap Y_{s'} = Y_{ss'}$이고, 마찬가지로
$D_{+}(f) \cap D_{+}(f') = D_{+}(ff')$이다. Lemma
\ref{constructions-lemma-ample-on-proj}에서 $D_{+}(f') \cap D_{+}(f)$는 $D_{+}(f)$의
점들 가운데 $\mathcal{O}_X(d)$의 단면 중 $f'$가 정의하는 것이 영이 되지 않는
점들의 집합과 같음을 보았다. 따라서
$\varphi_f^{-1}(D_{+}(f') \cap D_{+}(f)) = Y_s \cap Y_{s'}
= \varphi_{f'}^{-1}(D_{+}(f') \cap D_{+}(f))$이다.
$D_{+}(f) \cap D_{+}(f')$ 위에서 분수 $f/f'$는 구조층의 가역 단면이고
그 역은 $f'/f$이다. 또한
$\psi_{(f')}(f/f') = \psi(f)/s' = s/s'$이고
$\psi_{(f)}(f'/f) = \psi(f')/s = s'/s$임에 유의하자. 다음 도식을
가환하게 만드는 유일한 환의 사상
$S_{(ff')} \to \Gamma(Y_{ss'}, \mathcal{O}_Y)$이 존재한다고 주장한다.
$$
\xymatrix{
\Gamma(Y_s, \mathcal{O}_Y) \ar[r] &
\Gamma(Y_{ss'}, \mathcal{O}_Y) &
\Gamma(Y_{s, '} \mathcal{O}_Y) \ar[l]\\
S_{(f)} \ar[r] \ar[u]^{\psi_{(f)}} &
S_{(ff')} \ar[u] &
S_{(f')} \ar[l] \ar[u]^{\psi_{(f')}}
}
$$
위의 공식들에 의해 ``작동하는'' 규칙
$x/(ff')^n \mapsto \psi(x)/(ss')^n$을 사용할 수 있으므로 이 사상은
존재한다. 유일성은 $\text{Proj}(S)$가 분리되어 있다는 사실에서 따른다.
Lemma \ref{constructions-lemma-proj-separated}와 그 증명을 보라. 이로써 사상
$\varphi_f$와 $\varphi_{f'}$는 $Y_s \cap Y_{s'}$ 위에서 일치함을 알 수
있다. $\alpha_f$와 $\alpha_{f'}$의 제한도 $Y_s \cap Y_{s'}$ 위에서
일치한다. 이는 정칙함수 $s/s'$와 $\psi_{(f')}(f/f')$가 일치하기 때문이다.
따라서 사상 $\psi_f$들은 붙어서 $Y$에서 $U_d \subset X$로 가는 전역
사상을 이루고, 사상 $\alpha_f$들은 붙어서 보조정리의 조건을 만족하는
동형을 이룬다.

\medskip\noindent
이제 쌍 $(\varphi, \alpha)$가 유일함을 보여야 한다. $(\varphi', \alpha')$를
그러한 두 번째 쌍이라 하자. $f \in S_d$라 하자. 보조정리의 도식들이
가환하므로 $D_{+}(f)$의 $\varphi$와 $\varphi'$에 의한 역상은 모두
$Y_{\psi(f)}$와 같다. 열린집합 $D_{+}(f)$들은 $X$의 위상의 기저를 이루고
$X$는 소버 위상공간이므로(Schemes, Lemma
\ref{schemes-lemma-scheme-sober}를 보라), 이는 사상 $\varphi$와
$\varphi'$가 바탕 위상공간 위에서 같다는 뜻이다.
$s = \psi(f)$를 사용하여 가역층 $\mathcal{L}$을 $Y_{\psi(f)}$ 위에서
자명화하자. 도식들의 가환성에 의해
$\alpha^{\otimes n}(\psi^d_{\varphi}(x)) =
\psi(x) = (\alpha')^{\otimes n}(\psi^d_{\varphi'}(x))$가
모든 $x \in S_{nd}$에 대해 성립한다.
$\psi^d_{\varphi}$와 $\psi^d_{\varphi'}$의 구성에 의해
$\psi^d_{\varphi}(x) = \varphi^\sharp(x/f^n) \psi^d_{\varphi}(f^n)$이
$Y_{\psi(f)}$ 위에서 성립하고,
$\psi^d_{\varphi'}$에 대해서도 마찬가지이다. 보조정리 도식들의 가환성으로부터
$\varphi^\sharp(x/f^n) = (\varphi')^\sharp(x/f^n)$를 얻는다. 이로써
$\varphi$와 $\varphi'$는 $Y_{\psi(f)}$에서 아핀 스킴
$D_{+}(f) = \Spec(S_{(f)})$로 가는 같은 사상을 유도함을 증명하였다.
따라서 $\varphi$와 $\varphi'$는 사상으로서 같다. 마지막으로, 보조정리의
도식 가환성이 주어진 $\varphi$에 대해 유일한 $\alpha$를 골라냄을 보여야
한다. 확인은 생략한다.
\end{proof}

\noindent
보조정리 앞의 논의를 계속하자. $S$를 등급 환이라 하자. $Y$를 스킴이라 하자.
다음 조건을 만족하는 {\it 삼중쌍} $(d, \mathcal{L}, \psi)$를 생각하겠다.
\begin{enumerate}
\item $d \geq 1$은 정수이다.
\item $\mathcal{L}$은 가역 $\mathcal{O}_Y$-가군이다.
\item $\psi : S^{(d)} \to \Gamma_*(Y, \mathcal{L})$는 등급 환의
준동형이고, $\mathcal{L}$은 전역 단면 $\psi(f)$들에 의해 생성된다. 여기서
$f \in S_d$이다.
\end{enumerate}
사상 $h : Y' \to Y$와 삼중쌍 $(d, \mathcal{L}, \psi)$가 주어지고 후자가
$Y$ 위에 있으면,
이를 삼중쌍 $(d, h^*\mathcal{L}, h^* \circ \psi)$으로 당길 수 있다.
두 삼중쌍 $(d, \mathcal{L}, \psi)$와 $(d, \mathcal{L}', \psi')$가
같은 정수 $d$를 갖는다고 하자. 동형
$\beta : \mathcal{L} \to \mathcal{L}'$이 존재하여
$\beta \circ \psi = \psi'$가 등급 환의 사상
$S^{(d)} \to \Gamma_*(Y, \mathcal{L}')$으로서 성립하면 이들을
{\it 엄밀히 동치}라고 한다.

\medskip\noindent
각 정수 $d \geq 1$에 대해 위에서 정의한 당김을 사용하여
\begin{eqnarray*}
F_d : \Sch^{opp} & \longrightarrow & \textit{Sets}, \\
Y & \longmapsto &
\{\text{strict equivalence classes of triples }
(d, \mathcal{L}, \psi)
\text{ as above}\}
\end{eqnarray*}
를 정의한다.

\begin{lemma}
\label{constructions-lemma-proj-functor-strict}
$S$를 등급 환이라 하자. $X = \text{Proj}(S)$라 하자. 열린 부분스킴
$U_d \subset X$ (\ref{constructions-equation-Ud})는 함자 $F_d$를 표현하고, 위에서 정의한
삼중쌍 $(d, \mathcal{O}_{U_d}(d), \psi^d)$은 보편족이다(Schemes, Section
\ref{schemes-section-representable}을 보라).
\end{lemma}

\begin{proof}
이는 Lemma \ref{constructions-lemma-converse-construction}을 다시 서술한 것이다.
\end{proof}

\begin{lemma}
\label{constructions-lemma-apply}
$S$가 $S_0$-대수로서 $S_1$의 원소들에 의해 생성되는 등급 환이라 하자.
이 경우 스킴 $X = \text{Proj}(S)$는 스킴 $Y$에 다음 조건을 만족하는 쌍
$(\mathcal{L}, \psi)$들의 엄밀한 동치류의 집합을 대응시키는 함자를 표현한다.
\begin{enumerate}
\item $\mathcal{L}$은 가역 $\mathcal{O}_Y$-가군이다.
\item $\psi : S \to \Gamma_*(Y, \mathcal{L})$는 등급 환의 준동형이고,
$\mathcal{L}$은 전역 단면 $\psi(f)$들에 의해 생성된다. 여기서 $f \in S_1$이다.
\end{enumerate}
여기서 엄밀한 동치는 위와 같이 정의한다.
\end{lemma}

\begin{proof}
보조정리의 가정 아래 $X = U_1$이고, 이는 위의 Lemma
\ref{constructions-lemma-proj-functor-strict}을 다시 서술한 것이다.
\end{proof}

\noindent
$\text{Proj}(S)$에 대응하는 함자에 관한 논의를 일반 등급 환 $S$에 대해 이어가며 이 절을
마치겠다. 독자는 이 절의 나머지를 건너뛰어도 좋다.

\medskip\noindent
임의의 등급 환 $S$를 고정하자. $T$를 스킴이라 하자. 두 삼중쌍
$(d, \mathcal{L}, \psi)$와 $(d', \mathcal{L}', \psi')$가 $T$ 위에 있고,
그 정수 $d$, $d'$는 서로 달라도 된다고 하자. 동형
$\beta : \mathcal{L}^{\otimes d'} \to (\mathcal{L}')^{\otimes d}$가 $T$ 위의
가역층 사이에 존재하여
$\beta \circ \psi|_{S^{(dd')}}$와 $\psi'|_{S^{(dd')}}$가 등급 환의 사상
$S^{(dd')} \to \Gamma_*(Y, (\mathcal{L}')^{\otimes dd'})$으로서 일치하면
두 삼중쌍이 {\it 동치}라고 한다.

\begin{lemma}
\label{constructions-lemma-equivalent}
$S$를 등급 환이라 하자. $X = \text{Proj}(S)$로 두자. $T$를 스킴이라 하자.
$(d, \mathcal{L}, \psi)$와 $(d', \mathcal{L}', \psi')$를 $T$ 위의 두
삼중쌍이라 하자. 다음은 서로 동치이다.
\begin{enumerate}
\item $n = \text{lcm}(d, d')$로 두고 $n = ad = a'd'$로 쓰자. 다음 성질을
갖는 동형 $\beta : \mathcal{L}^{\otimes a} \to (\mathcal{L}')^{\otimes a'}$가
존재한다. 즉, $\beta \circ \psi|_{S^{(n)}}$와 $\psi'|_{S^{(n)}}$는 등급
환의 사상 $S^{(n)} \to \Gamma_*(Y, (\mathcal{L}')^{\otimes n})$으로서
일치한다.
\item 삼중쌍 $(d, \mathcal{L}, \psi)$와 $(d', \mathcal{L}', \psi')$는
동치이다.
\item 어떤 양의 정수 $n = ad = a'd'$에 대해, 다음 성질을 갖는 동형
$\beta : \mathcal{L}^{\otimes a} \to (\mathcal{L}')^{\otimes a'}$가
존재한다. 즉, $\beta \circ \psi|_{S^{(n)}}$와 $\psi'|_{S^{(n)}}$는 등급
환의 사상 $S^{(n)} \to \Gamma_*(Y, (\mathcal{L}')^{\otimes n})$으로서
일치한다.
\item 사상 $\varphi : T \to X$와 $\varphi' : T \to X$가 같으며, 이들은
각각 $(d, \mathcal{L}, \psi)$와 $(d', \mathcal{L}', \psi')$에 결부한
사상이다.
\end{enumerate}
\end{lemma}

\begin{proof}
(1)이 (2)를 함의함은 분명하고, 더 나눌 수 있는 차수와 가역층의 거듭제곱으로
제한하면 (2)가 (3)을 함의한다. 또한 Lemma
\ref{constructions-lemma-converse-construction}의 유일성 명제에 의해 (3)은 (4)를
함의한다. 따라서 (4)가 (1)을 함의함을 증명하면 된다. (4), 다시 말해
$\varphi = \varphi'$라고 가정하자. 이로부터
$\mathcal{L} = \varphi^*\mathcal{O}_X(d)$ 및
$\mathcal{L}' = \varphi^*\mathcal{O}_X(d')$로 쓸 수 있음에 유의하자.
또한 이 식별을 통해 등급 환의 사상 $\psi$와 $\psi'$는 표준 등급 환 사상
$$
S \longrightarrow \bigoplus\nolimits_{n \geq 0} \Gamma(X, \mathcal{O}_X(n))
$$
을 $S^{(d)}$ 및 $S^{(d')}$에 제한한 뒤 $\varphi$로 당긴 것에 대응한다
(다시 Lemma \ref{constructions-lemma-converse-construction}에 의한다). 따라서 $\beta$를
동형
$$
(\varphi^*\mathcal{O}_X(d))^{\otimes a} =
\varphi^*\mathcal{O}_X(n) =
(\varphi^*\mathcal{O}_X(d'))^{\otimes a'}
$$
으로 택하면 된다.
\end{proof}

\noindent
$S$를 등급 환이라 하자. $X = \text{Proj}(S)$라 하자. 열린 부분스킴
$U_d \subset X = \text{Proj}(S)$ (\ref{constructions-equation-Ud}) 위에는 삼중쌍
$(d, \mathcal{O}_{U_d}(d), \psi^d)$이 있다. $d | d'$이면 삼중쌍
$(d, \mathcal{O}_{U_d}(d), \psi^d)$과
$(d', \mathcal{O}_{U_{d'}}(d'), \psi^{d'})$은 열린집합 $U_d$로 제한했을
때 동치임이 분명하다(그리고 이는 $U_{d'}$의 부분집합이다). 이를 Lemma
\ref{constructions-lemma-converse-construction}과 합치면, 사상 $Y \to X$는 대략 $Y$ 위의
삼중쌍의 동치류에 대응한다. 하지만 $Y$가 준콤팩트가 아니면 작동하는 단 하나의
삼중쌍이 존재하지 않을 수도 있으므로 이는 완전히 정확하지는 않다. 따라서 대응하는
함자를 정의할 때 약간 주의해야 한다.

\medskip\noindent
이를 수행하는 한 가지 방법은 다음과 같다. $d' = ad$라고 가정하자. 함자의 변환
$F_d \to F_{d'}$를 생각하자. 이는 삼중쌍 $(d, \mathcal{L}, \psi)$가
$T$ 위에 주어지면 이를 삼중쌍
$(d', \mathcal{L}^{\otimes a}, \psi|_{S^{(d')}})$에 대응시킨다.
Lemma \ref{constructions-lemma-equivalent}의 함의 중 하나는
변환 $F_d \to F_{d'}$가 단사라는 것이다! 준콤팩트 스킴 $T$에 대해
$$
F(T) = \bigcup\nolimits_{d \in \mathbf{N}} F_d(T)
$$
로 정의하되, 전이사상은 위에서 설명한 것을 사용한다. 이는 준콤팩트 스킴의
범주에서 집합의 범주로 가는 반변함자를 분명히 정의한다. 일반 스킴 $T$에 대해서는
$$
F(T)
=
\lim_{V \subset T\text{ quasi-compact open}} F(V).
$$
로 정의한다. 다시 말해 원소 $\xi$가 $F(T)$에 속한다는 것은 호환되는 원소
$\xi_V \in F(V)$들을 택하는 것에 대응한다. 여기서 $V$는 $T$의 준콤팩트
열린집합들을 달린다. 당김사상 $F(T) \to F(T')$의 정의는 스킴의 사상
$T' \to T$에 대해 생략한다. 이로써 함자
\begin{eqnarray*}
F : \Sch^{opp} & \longrightarrow & \textit{Sets}
\end{eqnarray*}
를 정의하였다.

\begin{lemma}
\label{constructions-lemma-proj-functor}
$S$를 등급 환이라 하자. $X = \text{Proj}(S)$라 하자. 위에서 정의한 함자
$F$는 스킴 $X$에 의해 표현가능하다.
\end{lemma}

\begin{proof}
위에서 함자 $F_d$가 열린 부분스킴 $U_d \subset X$에 대응함을 보았다. 또한
위에서 정의한 함자의 변환 $F_d \to F_{d'}$($d | d'$인 경우)은 포함사상
$U_d \to U_{d'}$에 대응한다(위의 논의를 보라). 따라서 $F$가 $X$에 의해
표현됨을 보이려면, 사상 $T \to X$가 준콤팩트 스킴 $T$에 대해 어떤
$U_d$ 안으로 들어가고, 일반 스킴 $T$에 대해
$$
\Mor(T, X)
=
\lim_{V \subset T\text{ quasi-compact open}} \Mor(V, X).
$$
임을 보이면 충분하다. 이 확인들은 생략한다.
\end{proof}

\section{사영공간}
\label{constructions-section-projective-space}

\noindent
사영공간은 대수기하학에서 연구하는 기본적인 대상 가운데 하나이다. 이 절에서는
다항식환의 $\text{Proj}$로서의 구성만 제시한다. 뒤에서 그 아름다운 성질들을
많이 발견할 것이다.

\begin{lemma}
\label{constructions-lemma-projective-space}
$S = \mathbf{Z}[T_0, \ldots, T_n]$이고 $\deg(T_i) = 1$이라 하자. 스킴
$$
\mathbf{P}^n_{\mathbf{Z}} = \text{Proj}(S)
$$
은 스킴 $Y$에 다음 조건을 만족하는 쌍
$(\mathcal{L}, (s_0, \ldots, s_n))$을 대응시키는 함자를 표현한다.
\begin{enumerate}
\item $\mathcal{L}$은 가역 $\mathcal{O}_Y$-가군이다.
\item $s_0, \ldots, s_n$은 $\mathcal{L}$의 전역 단면들이고
$\mathcal{L}$을 생성한다.
\end{enumerate}
여기서 다음 동치로 나눈다.
$(\mathcal{L}, (s_0, \ldots, s_n)) \sim
(\mathcal{N}, (t_0, \ldots, t_n))$ $\Leftrightarrow$ 동형
$\beta : \mathcal{L} \to \mathcal{N}$이 존재하여
$\beta(s_i) = t_i$가 $i = 0, \ldots, n$에 대해 성립한다.
\end{lemma}

\begin{proof}
이는 위의 Lemma \ref{constructions-lemma-apply}의 특수한 경우이다. 실제로 임의의 등급 환
$A$에 대해
\begin{eqnarray*}
\Mor_{graded rings}(\mathbf{Z}[T_0, \ldots, T_n], A)
& = &
A_1 \times \ldots \times A_1 \\
\psi & \mapsto & (\psi(T_0), \ldots, \psi(T_n))
\end{eqnarray*}
이고, 차수 $1$ 부분을 $\Gamma_*(Y, \mathcal{L})$에서 취하면 바로
$\Gamma(Y, \mathcal{L})$이다.
\end{proof}

\begin{definition}
\label{constructions-definition-projective-space}
스킴 $\mathbf{P}^n_{\mathbf{Z}} = \text{Proj}(\mathbf{Z}[T_0, \ldots, T_n])$을
{\it 사영 $n$-공간($\mathbf{Z}$ 위)}이라 한다. 그 기저변환
$\mathbf{P}^n_S$를 스킴 $S$에 대해 취한 것을
{\it 사영 $n$-공간($S$ 위)}이라 한다. $R$이 환이면
$\Spec(R)$로의 기저변환을 $\mathbf{P}^n_R$로 나타내고
{\it 사영 $n$-공간($R$ 위)}이라 한다.
\end{definition}

\noindent
스킴 $Y$가 $S$ 위에 있고, Lemma \ref{constructions-lemma-projective-space}에서와 같은 쌍
$(\mathcal{L}, (s_0, \ldots, s_n))$이 주어졌을 때, 유도된
$\mathbf{P}^n_S$로의 사상을
$$
\varphi_{(\mathcal{L}, (s_0, \ldots, s_n))} :
Y \longrightarrow \mathbf{P}^n_S
$$
로 나타낸다. 이 쌍은 $\mathbf{P}^n_{\mathbf{Z}}$로의 사상을 정의하고 이미
$S$로의 구조사상이 있으므로, 이 둘을 합치면
$\mathbf{P}^n_S = \mathbf{P}^n_{\mathbf{Z}} \times S$로의 사상을 얻는다.
따라서 이 표기는 의미가 있다. Lemma \ref{constructions-lemma-relative-proj-base-change}에
의해 $\mathbf{P}^n_S$를 등급 $\mathcal{O}_S$-대수
$\mathcal{O}_S[T_0, \ldots, T_n]$의 상대 Proj와 식별할 수 있고,
$\varphi_{(\mathcal{L}, (s_0, \ldots, s_n))}$를 다음을 만족하는 유일한
$S$-사상으로 특징지을 수 있다.
$$
\mathcal{L} =
\varphi_{(\mathcal{L}, (s_0, \ldots, s_n))}^*\mathcal{O}_{\mathbf{P}^n_S}(1)
\quad
\text{and}
\quad
s_i = \varphi_{(\mathcal{L}, (s_0, \ldots, s_n))}^*T_i
$$
여기서 $T_i$는 $\mathcal{O}_{\mathbf{P}^n_S}(1)$의 전역 단면으로 본다.
이는 Lemma \ref{constructions-lemma-glue-relative-proj-twists}의 사상 $\psi$를 통한다. 세부사항은
생략한다.

\begin{lemma}
\label{constructions-lemma-standard-covering-projective-space}
사영 $n$-공간($\mathbf{Z}$ 위)은 $n + 1$개의 표준 열린집합으로 덮인다.
$$
\mathbf{P}^n_{\mathbf{Z}} =
\bigcup\nolimits_{i = 0, \ldots, n} D_{+}(T_i)
$$
각 $D_{+}(T_i)$는 $\mathbf{A}^n_{\mathbf{Z}}$와 동형이며, 후자는 아핀
$n$-공간($\mathbf{Z}$ 위)이다.
\end{lemma}

\begin{proof}
이는
$\mathbf{Z}[T_0, \ldots, T_n]_{+} = (T_0, \ldots, T_n)$이고,
$$
\Spec
\left(
\mathbf{Z}
\left[\frac{T_0}{T_i}, \ldots, \frac{T_n}{T_i}
\right]
\right)
\cong
\mathbf{A}^n_{\mathbf{Z}}
$$
가 자명한 방식으로 성립하기 때문이다.
\end{proof}

\begin{lemma}
\label{constructions-lemma-projective-space-separated}
$S$를 스킴이라 하자. 구조사상 $\mathbf{P}^n_S \to S$는
\begin{enumerate}
\item 분리이고,
\item 준콤팩트이며,
\item 값매김 판정법의 존재 부분과 유일성 부분을 만족하고,
\item 보편 닫힌 사상이다.
\end{enumerate}
\end{lemma}

\begin{proof}
이 성질들은 모두 기저변환 아래 안정적이다. 마지막 두 성질에 대해서는 분명하고,
다른 두 성질에 대해서는 Schemes, Lemmas
\ref{schemes-lemma-separated-permanence} 및
\ref{schemes-lemma-quasi-compact-preserved-base-change}를 보라. 따라서 사상
$\mathbf{P}^n_{\mathbf{Z}} \to \Spec(\mathbf{Z})$에 대해 증명하면 충분하다.
분리성은 Lemma \ref{constructions-lemma-proj-separated}이다. 준콤팩트성은 Lemma
\ref{constructions-lemma-standard-covering-projective-space}에서 따른다. 값매김 판정법의
존재성과 유일성은 Lemma \ref{constructions-lemma-proj-valuative-criterion}에서 따른다.
보편 닫힘은 위의 결과와 Schemes, Proposition
\ref{schemes-proposition-characterize-universally-closed}에서 따른다.
\end{proof}

\begin{remark}
\label{constructions-remark-missing-finite-type}
위의 성질 목록에서 빠진 것은 무엇인가? 물론 유한형이라는 성질이다. 여기서 이를
열거하지 않은 까닭은 이 시점에는 아직 유한형의 개념을 정의하지 않았기 때문이다.
(빠진 또 다른 성질은 ``매끄러움''이다. 독자는 틀림없이 그 밖에도 많은 성질을
생각해 낼 수 있을 것이다.)
\end{remark}

\begin{lemma}[세그레 매장]
\label{constructions-lemma-segre-embedding}
$S$를 스킴이라 하자. 닫힌 몰입
$$
\mathbf{P}^n_S \times_S \mathbf{P}^m_S
\longrightarrow
\mathbf{P}^{nm + n + m}_S
$$
이 존재하며, 이를 {\it 세그레 매장}이라 한다.
\end{lemma}

\begin{proof}
$S = \Spec(\mathbf{Z})$일 때 증명하면 충분하다. 따라서 첨자 $S$를 생략하고
절대적인 상황에서 논하겠다.
$\mathbf{P}^n = \text{Proj}(\mathbf{Z}[X_0, \ldots, X_n])$,
$\mathbf{P}^m = \text{Proj}(\mathbf{Z}[Y_0, \ldots, Y_m])$ 및
$\mathbf{P}^{nm + n + m} =
\text{Proj}(\mathbf{Z}[Z_0, \ldots, Z_{nm + n + m}])$로 쓰자.
$\mathbf{P}^{nm + n + m}$로 사상하려면 왼쪽 변 위의 가역층
$\mathcal{L}$과 이를 생성하는 $(n + 1)(m + 1)$개의 단면 $s_i$를
적어야 한다. Lemma \ref{constructions-lemma-projective-space}를 보라. 우리가 택하는
가역층은
$$
\mathcal{L} =
\text{pr}_1^*\mathcal{O}_{\mathbf{P}^n}(1)
\otimes
\text{pr}_2^*\mathcal{O}_{\mathbf{P}^m}(1)
$$
이다. 택하는 단면들은
$$
s_0 = X_0Y_0, \ s_1 = X_1Y_0, \ldots, \ s_n = X_nY_0,
\ s_{n + 1} = X_0Y_1, \ldots, \ s_{nm + n + m} = X_nY_m.
$$
이다. 이 단면들은 $\mathcal{L}$을 생성한다. 실제로 단면 $X_i$들이
$\mathcal{O}_{\mathbf{P}^n}(1)$을 생성하고 단면 $Y_j$들이
$\mathcal{O}_{\mathbf{P}^m}(1)$을 생성한다. 유도된 사상 $\varphi$는
$$
\varphi^{-1}(D_{+}(Z_{i + (n + 1)j})) = D_{+}(X_i) \times D_{+}(Y_j).
$$
라는 성질을 가진다. 따라서 이는 아핀 사상이다. $(i, j) = (0, 0)$인 경우
대응하는 환의 사상은
$$
\mathbf{Z}[Z_1/Z_0, \ldots, Z_{nm + n + m}/Z_0]
\longrightarrow
\mathbf{Z}[X_1/X_0, \ldots, X_n/X_0, Y_1/Y_0, \ldots, Y_n/Y_0]
$$
이며, $Z_i/Z_0$을 원소 $X_i/X_0$으로 보낸다. 여기서 $i \leq n$이다. 또한
$Z_{(n + 1)j}/Z_0$을 원소 $Y_j/Y_0$으로 보낸다. 따라서 전사이다. 다른 아핀
열린 부분집합에 대해서도 같은 논증이 성립한다. 그러므로 사상 $\varphi$는 닫힌
몰입이다(Schemes, Lemma \ref{schemes-lemma-closed-local-target} 및
Example \ref{schemes-example-closed-immersion-affines}를 보라).
\end{proof}

\noindent
다음 두 보조정리는 뒤에서 나올 더 일반적인 결과들의 특수한 경우이지만, 여기서
직접 증명하는 편이 타당할 것이다. 첫 번째 것은 Divisors, Lemma
\ref{divisors-lemma-closed-subscheme-proj}의 특수한 경우이다.

\begin{lemma}
\label{constructions-lemma-closed-in-projective-space}
$R$을 환이라 하자. $Z \subset \mathbf{P}^n_R$를 닫힌 부분스킴이라 하자.
다음과 같이 두자.
$$
I_d = \Ker\left(
R[T_0, \ldots, T_n]_d
\longrightarrow
\Gamma(Z, \mathcal{O}_{\mathbf{P}^n_R}(d)|_Z)\right)
$$
그러면 $I = \bigoplus I_d \subset R[T_0, \ldots, T_n]$은 등급 아이디얼이고
$Z = \text{Proj}(R[T_0, \ldots, T_n]/I)$이다.
\end{lemma}

\begin{proof}
$I$가 등급 아이디얼임은 분명하다.
$Z' = \text{Proj}(R[T_0, \ldots, T_n]/I)$로 두자. Lemma
\ref{constructions-lemma-surjective-graded-rings-generated-degree-1-map-proj}에 의해
$Z'$는 $\mathbf{P}^n_R$의 닫힌 부분스킴이다. 등식 $Z = Z'$를 보이려면
표준 아핀 열린집합 $D_{+}(T_i)$ 위에서 확인하면 충분하다. 동차좌표의 번호를
다시 매겨 $i = 0$이라고 가정해도 된다.
$Z \cap D_{+}(T_0)$ 및 $Z' \cap D_{+}(T_0)$가 각각 아이디얼 $J$ 및 $J'$에
의해 잘리고, 이들은 $R[T_1/T_0, \ldots, T_n/T_0]$의 아이디얼이라고 하자.
그러면 $J'$는 원소 $F/T_0^{\deg(F)}$들로 생성되는 아이디얼이다. 여기서
$F \in I$는 동차원소이다. $F \in I$의 차수가 $d$라고 가정하자. $F$는
$\mathcal{O}_{\mathbf{P}^n_R}(d)$의 단면으로서 $Z$에 제한하면
영이므로 $F/T_0^d$는 $J$의 원소이다. 따라서 $J' \subset J$이다.

\medskip\noindent
역으로 $f \in J$라고 가정하자. $f$의 전체 차수가 $d$이고, 이는
$T_1/T_0, \ldots, T_n/T_0$에 관한 차수라고 하자. 그러면
$f = F/T_0^d$로 쓸 수 있는 어떤 $F \in R[T_0, \ldots, T_n]_d$가 존재한다.
$i \in \{1, \ldots, n\}$를 택하자. 그러면
$Z \cap D_{+}(T_i)$는 어떤 아이디얼
$J_i \subset R[T_0/T_i, \ldots, T_n/T_i]$에 의해 잘린다. 또한
$$
J \cdot
R\left[
\frac{T_1}{T_0}, \ldots, \frac{T_n}{T_0},
\frac{T_0}{T_i}, \ldots, \frac{T_n}{T_i}
\right]
=
J_i \cdot
R\left[
\frac{T_1}{T_0}, \ldots, \frac{T_n}{T_0},
\frac{T_0}{T_i}, \ldots, \frac{T_n}{T_i}
\right]
$$
이다. 왼쪽은 $J$의 국소화로서 원소 $T_i/T_0$에 관한 것이고, 오른쪽은
$J_i$의 국소화로서 원소 $T_0/T_i$에 관한 것이다. 따라서
$T_0^{d_i}F/T_i^{d + d_i}$가 $J_i$의 원소인 충분히 큰 어떤 $d_i$가
존재한다. 이로부터
$T_0^{\max(d_i)}F$가 $I$의 원소임이 따른다. 실제로 그 제한은 각 표준 아핀
열린집합 $D_{+}(T_i)$에서 닫힌 부분스킴 $Z \cap D_{+}(T_i)$ 위에 영이다.
따라서 $f \in J'$이고, 원하는 대로 $J \subset J'$를 얻는다.
\end{proof}

\noindent
다음 보조정리는 더 일반적인 Properties, Lemmas
\ref{properties-lemma-ample-quasi-coherent} 또는
\ref{properties-lemma-proj-quasi-coherent}의 특수한 경우이다.

\begin{lemma}
\label{constructions-lemma-quasi-coherent-projective-space}
$R$을 환이라 하자. $\mathcal{F}$를 $\mathbf{P}^n_R$ 위의 준연접층이라 하자.
$d \geq 0$에 대해
$$
M_d
=
\Gamma(\mathbf{P}^n_R,
\mathcal{F} \otimes_{\mathcal{O}_{\mathbf{P}^n_R}}
\mathcal{O}_{\mathbf{P}^n_R}(d))
=
\Gamma(\mathbf{P}^n_R, \mathcal{F}(d))
$$
로 두자. 그러면 $M = \bigoplus_{d \geq 0} M_d$는 등급
$R[T_0, \ldots, R_n]$-가군이고, 표준 동형
$\mathcal{F} = \widetilde{M}$이 존재한다.
\end{lemma}

\begin{proof}
곱셈 사상
$$
R[T_0, \ldots, R_n]_e \times M_d \longrightarrow M_{d + e}
$$
은 자연스러운 동형
$$
\mathcal{O}_{\mathbf{P}^n_R}(e)
\otimes_{\mathcal{O}_{\mathbf{P}^n_R}}
\mathcal{F}(d)
\longrightarrow
\mathcal{F}(e + d)
$$
에서 온다. Equation (\ref{constructions-equation-global-sections-module})을 보라. 사상
$c : \widetilde{M} \to \mathcal{F}$를 구성하자. 각 표준 아핀
$U_i = D_{+}(T_i)$ 위에서
$\Gamma(U_i, \widetilde{M}) = (M[1/T_i])_0$이고, 여기서 첨자 ${}_0$은
차수 $0$ 부분을 뜻한다. 이 원소는 $m/T_i^d$로 쓸 수 있으며, 여기서
$m \in M_d$이다. $T_i$는 $\mathcal{O}(1)$의 생성자이고 이는 $U_i$ 위에서
$m|_{U_i} = m_i \otimes T_i^d$로 쓸 수 있고, 여기서
$m_i \in \Gamma(U_i, \mathcal{F})$는 유일한 단면이다. 따라서 자연스러운
추측은 $c(m/T_i^d) = m_i$이다. 여기서 생략하는 짧은 논증에 따르면,
이 처방이 잘 정의된 사상 $c : \widetilde{M} \to \mathcal{F}$를 주는지는
$$
(T_i/T_j)^d m_i|_{U_i \cap U_j} = m_j|_{U_i \cap U_j}
$$
가 $M[1/T_iT_j]$에서 성립함을 보이는 것으로 귀착된다. 그러나 겹침 위에서
생성자 $T_i$와 $T_j$는 $\mathcal{O}(1)$에 속하고 가역함수 $T_i/T_j$만큼
다르므로 이는 분명하다.

\medskip\noindent
$c$의 단사성. 아핀 열린집합 $U_i$들 위에서 단사성을 확인해도 된다.
$i \in \{0, \ldots, n\}$라 하고, 원소 $s$가
$s = m/T_i^d \in \Gamma(U_i, \widetilde{M})$를 만족하며
$c(m/T_i^d) = 0$을 만족한다고 하자. 위의 $c$에 대한 설명에 의해 이는
$m_i = 0$이라는 뜻이고, 따라서 $m|_{U_i} = 0$이다. 그러므로
$T_i^em = 0$이 $M$에서 성립하는 어떤 $e$가 존재한다. 따라서 원하는 대로
$s = m/T_i^d = T_i^e/T_i^{e + d} = 0$이다.

\medskip\noindent
$c$의 전사성. 아핀 열린집합 $U_i$들 위에서 전사성을 확인해도 된다. 번호를
다시 매기면 $U_0$ 위에서 확인하는 것으로 충분하다.
$s \in \mathcal{F}(U_0)$라 하자. $\mathcal{F}|_{U_i} = \widetilde{N_i}$로
쓰자. 여기서 어떤 $R[T_0/T_i, \ldots, T_0/T_i]$-가군 $N_i$를 사용하며,
$\mathcal{F}$가 준연접이므로 이렇게 쓸 수 있다. 따라서 $s$는 어떤 원소
$x \in N_0$에 대응한다. 이제
$$
(N_i)_{T_j/T_i} \cong (N_j)_{T_i/T_j}
$$
이고(첨자는 ‘주 아이디얼 국소화’를 뜻한다), 이는 환
$$
R\left[
\frac{T_0}{T_i}, \ldots, \frac{T_n}{T_i},
\frac{T_0}{T_j}, \ldots, \frac{T_n}{T_j}
\right].
$$
위의 가군의 동형이다. 따라서 어떤 큰 정수 $d$에 대해 원소
$s_i \in N_i$, $i = 1, \ldots, n$이 존재하여
$$
s = (T_i/T_0)^d s_i
$$
가 $U_0 \cap U_i$ 위에서 성립한다. 다음으로 $U_i \cap U_j$ 위의 차
$$
t_{ij} = s_i - (T_j/T_i)^d s_j
$$
를 보자. 여기서 $0 < i < j$이다. $s_i$의 선택에 의해
$t_{ij}|_{U_0 \cap U_i \cap U_j} = 0$임을 안다. 따라서 어떤 큰 정수
$e$에 대해 $(T_0/T_i)^et_{ij} = 0$이다.
$s_i' = (T_0/T_i)^es_i$ 및 $s_0' = s$로 두자. 그러면
$$
s_a' = (T_b/T_a)^{e + d} s_b'
$$
가 $U_a \cap U_b$ 위에서 모든 $a, b$에 대해 성립한다. 이는 원소 $s'_a$들이 붙어서 전역 단면
$m \in \Gamma(\mathbf{P}^n_R, \mathcal{F}(e + d))$을 이룬다는 조건과
정확히 같다. 더욱이 구성에 의해 $c(m/T_0^{e + d}) = s$이다. 따라서
$c$는 전사이고 증명이 끝난다.
\end{proof}

\begin{lemma}
\label{constructions-lemma-globally-generated-omega-twist-1}
$X$를 스킴이라 하자. $\mathcal{L}$을 가역층이라 하고,
$s_0, \ldots, s_n$을 $\mathcal{L}$을 생성하는 전역 단면들이라 하자.
$\mathcal{F}$를 유도된 사상
$\mathcal{O}_X^{\oplus n + 1} \to \mathcal{L}$의 핵이라 하자. 그러면
$\mathcal{F} \otimes \mathcal{L}$은 전역 생성된다.
\end{lemma}

\begin{proof}
실제로 $X$가 임의의 국소 환 달린 공간이어도 결과는 참이다. 층
$\mathcal{F}$는 유한 국소 자유 $\mathcal{O}_X$-가군이고 그 계수는 $n$이다. 원소
$$
s_{ij} = (0, \ldots, 0, s_j, 0, \ldots, 0, -s_i, 0, \ldots, 0)
\in \Gamma(X, \mathcal{L}^{\oplus n + 1})
$$
에서 $s_j$는 $i$번째 자리에 있고 $-s_i$는 $j$번째 자리에 있으며, 이 원소는
$\mathcal{L}^{\otimes 2}$에서 영으로 사상된다. 따라서
$s_{ij} \in \Gamma(X, \mathcal{F} \otimes_{\mathcal{O}_X} \mathcal{L})$이다.
국소 계산에 의해 이 단면들이 $\mathcal{F} \otimes \mathcal{L}$을 생성한다.

\medskip\noindent
다른 증명. 사상 $\varphi : X \to \mathbf{P}^n_\mathbf{Z}$을 생각하자. 이는
쌍 $(\mathcal{L}, (s_0, \ldots, s_n))$에 결부한 것이다. $\mathcal{O}(1)$의
당김은 $\mathcal{L}$이고 $T_i$의 당김은 $s_i$이므로,
$\mathbf{P}^n_\mathbf{Z}$의 경우에 보조정리를 증명하면 충분하다. 이 경우 층
$\mathcal{F}$는 등급 $S = \mathbf{Z}[T_0, \ldots, T_n]$-가군 $M$에
대응하며, 이는 짧은 완전열
$$
0 \to M \to S^{\oplus n + 1} \to S(1) \to 0
$$
에 놓인다. 여기서 두 번째 사상은 $T_0, \ldots, T_n$으로 주어진다. 이 경우
위의 명제는 원소
$$
T_{ij} = (0, \ldots, 0, T_j, 0, \ldots, 0, -T_i, 0, \ldots, 0)
\in M(1)_0
$$
들이 등급 가군 $M(1)$을 $S$ 위에서 생성한다는 명제로 옮겨진다. 세부사항은
생략한다.
\end{proof}

\section{가역층과 Proj로의 사상}
\label{constructions-section-invertible-proj}

\noindent
$T$를 스킴이라 하고 $\mathcal{L}$을 $T$ 위의 가역층이라 하자. 단면
$s \in \Gamma(T, \mathcal{L})$에 대해 $T_s$로 $s$가 영이 되지 않는 점들의
열린 부분집합을 나타낸다. Modules, Lemma \ref{modules-lemma-s-open}을 보라.
다음 보조정리는 Lemma \ref{constructions-lemma-apply}의 약간의 일반화로 볼 수 있다. 또한
Lemma \ref{constructions-lemma-morphism-proj}의 일반화이기도 하다.

\begin{lemma}
\label{constructions-lemma-invertible-map-into-proj}
$A$를 등급 환이라 하자. $X = \text{Proj}(A)$로 두자. $T$를 스킴이라 하자.
$\mathcal{L}$을 가역 $\mathcal{O}_T$-가군이라 하자.
$\psi : A \to \Gamma_*(T, \mathcal{L})$를 등급 환의 준동형이라 하자. 다음과
같이 두자.
$$
U(\psi) = \bigcup\nolimits_{f \in A_{+}\text{ homogeneous}} T_{\psi(f)}
$$
사상 $\psi$는 스킴의 표준 사상
$$
r_{\mathcal{L}, \psi} :
U(\psi) \longrightarrow X
$$
과 함께 $\mathbf{Z}$-등급 $\mathcal{O}_T$-대수의 사상
$$
\theta :
r_{\mathcal{L}, \psi}^*\left(
\bigoplus\nolimits_{d \in \mathbf{Z}} \mathcal{O}_X(d)
\right)
\longrightarrow
\bigoplus\nolimits_{d \in \mathbf{Z}} \mathcal{L}^{\otimes d}|_{U(\psi)}.
$$
을 유도한다. 삼중쌍 $(U(\psi), r_{\mathcal{L}, \psi}, \theta)$은 다음
성질들로 특징지어진다.
\begin{enumerate}
\item 동차원소 $f \in A_{+}$에 대해
$r_{\mathcal{L}, \psi}^{-1}(D_{+}(f)) = T_{\psi(f)}$이다.
\item 모든 $d \geq 0$에 대해 도식
$$
\xymatrix{
A_d \ar[d]_{(\ref{constructions-equation-global-sections})} \ar[r]_{\psi} &
\Gamma(T, \mathcal{L}^{\otimes d}) \ar[d]^{restrict} \\
\Gamma(X, \mathcal{O}_X(d)) \ar[r]^{\theta} &
\Gamma(U(\psi), \mathcal{L}^{\otimes d})
}
$$
은 가환한다.
\end{enumerate}
또한 임의의 $d \geq 1$과 열린 부분스킴 $V \subset T$에 대해, $\psi(A_d)$의
단면들이 $\mathcal{L}^{\otimes d}|_V$를 생성한다고 하자. 그러면 사상
$r_{\mathcal{L}, \psi}|_V$는 사상 $\varphi : V \to \text{Proj}(A)$와
일치하고, 사상 $\theta|_V$는 사상
$\alpha : \varphi^*\mathcal{O}_X(d) \to \mathcal{L}^{\otimes d}|_V$와
일치한다. 여기서 $(\varphi, \alpha)$는
$\psi|_{A^{(d)}} : A^{(d)} \to \Gamma_*(V, \mathcal{L}^{\otimes d})$에
결부한 Lemma \ref{constructions-lemma-converse-construction}의 쌍이다.
\end{lemma}

\begin{proof}
(1)과 (2)를 만족하는 두 삼중쌍 $(U, r : U \to X, \theta)$ 및
$(U', r' : U' \to X, \theta')$이 있다고 하자. 성질 (1)에 의해
$U = U' = U(\psi)$이고, 바탕 위상공간의 사상으로서 $r = r'$이다. 이는
열린집합 $D_{+}(f)$들이 $X$의 위상의 기저를 이루고 $X$가 소버
위상공간이기 때문이다(Algebra, Section \ref{algebra-section-proj} 및
Schemes, Lemma \ref{schemes-lemma-scheme-sober}를 보라).
$f \in A_{+}$를 동차원소라 하자.
$\Gamma(D_{+}(f), \bigoplus_{n \in \mathbf{Z}} \mathcal{O}_X(n)) = A_f$
임에 유의하자. 이는 $\mathbf{Z}$-등급 대수의 등식이다. 두 $\mathbf{Z}$-등급 환의 사상
$$
\theta, \theta' :
A_f
\longrightarrow
\Gamma(T_{\psi(f)}, \bigoplus \mathcal{L}^{\otimes n}).
$$
을 생각하자. 왼쪽에서는 $f$에 의한 곱셈이, 오른쪽에서는 $\psi(f)$에 의한
곱셈이 동형임을 안다. 또한 (2)에 의해
$\theta(x/1) = \theta'(x/1) = \psi(x)|_{T_{\psi(f)}}$임을 모든 $x \in A$에
대해 안다. 따라서
원하는 대로 $\theta = \theta'$임을 쉽게 얻는다. 차수 $0$ 부분을 생각하면
$r^\sharp = (r')^\sharp$, 즉 스킴의 사상으로서 $r = r'$임을 얻는다.
이로써 유일성을 증명하였다.

\medskip\noindent
이제 존재성을 보이자. 방금 증명한 유일성에 의해 쌍 $(r, \theta)$를 $T$ 위에서
국소적으로 구성하면 충분하다. 따라서 $T = \Spec(R)$이 아핀이고
$\mathcal{L} = \mathcal{O}_T$이며, 어떤 동차원소 $f \in A_{+}$에 대해
$\psi(f)$가 $\mathcal{O}_T = \mathcal{O}_T^{\otimes \deg(f)}$를 생성한다고
가정해도 된다. 다시 말해 $\psi(f) = u \in R^*$는 가역원이다. 이 경우
사상 $\psi$는 등급 환의 사상
$$
A \longrightarrow R[x] = \Gamma_*(T, \mathcal{O}_T)
$$
이고, $f$를 $ux^{\deg(f)}$로 보낸다. 이는 $\mathbf{Z}$-등급 환의 사상
$\theta : A_f \to R[x, x^{-1}]$로 유일하게 확장되며, $1/f$를
$u^{-1}x^{-\deg(f)}$로 보낸다. 이 사상의 차수 영
부분은 환의 사상 $A_{(f)} \to R$이고, 이는 사상
$r : T = \Spec(R) \to \Spec(A_{(f)}) = D_{+}(f) \subset X$를 준다.
따라서 이 특수한 경우에 $(r, \theta)$를 구성하였다.

\medskip\noindent
보조정리의 마지막 명제를 보이자. Lemma
\ref{constructions-lemma-converse-construction}에 따르면, 거기서 구성한 사상은 그 명제에
표시된 도식이 가환하도록 하는 유일한 사상이다. 이 보조정리의 도식이 가환하면
$V$와 $A^{(d)}$로 제한한 도식도 가환한다. 따라서 결과가 따른다.
\end{proof}

\begin{remark}
\label{constructions-remark-not-in-invertible-locus}
위의 Lemma \ref{constructions-lemma-invertible-map-into-proj}와 같은 가정을 사용하자.
사상 $r_{\mathcal{L}, \psi}$의 상은 층 $\mathcal{O}_X(1)$이 가역인 자리에
포함될 필요가 없다. 다음은 그 예이다. $k$를 체라 하자.
$S = k[A, B, C]$에 $\deg(A) = 1$, $\deg(B) = 2$, $\deg(C) = 3$으로
등급을 주자. $X = \text{Proj}(S)$로 두자.
$T = \mathbf{P}^2_k = \text{Proj}(k[X_0, X_1, X_2])$로 두자.
$\mathcal{L} = \mathcal{O}_T(1)$은 가역이고
$\mathcal{O}_T(n) = \mathcal{L}^{\otimes n}$임을 상기하자. 사상들의 합성
$\psi$를 생각하자.
$$
S \to k[X_0, X_1, X_2] \to \Gamma_*(T, \mathcal{L}).
$$
여기서 첫 번째 사상은 $A \mapsto X_0$, $B \mapsto X_1^2$,
$C \mapsto X_2^3$이고, 두 번째 사상은 (\ref{constructions-equation-global-sections})이다.
보조정리에 의해 이는 사상
$r_{\mathcal{L}, \psi} : T \to X = \text{Proj}(S)$에 대응하며, 이 사상이
전사임은 쉽게 알 수 있다. 한편 Remark \ref{constructions-remark-not-isomorphism}에서 층
$\mathcal{O}_X(1)$이 $X$의 모든 점에서 가역인 것은 아님을 보였다.
\end{remark}

\section{붙이기를 통한 상대 Proj}
\label{constructions-section-relative-proj-via-glueing}

\begin{situation}
\label{constructions-situation-relative-proj}
여기서 $S$는 스킴이고 $\mathcal{A}$는 준연접 등급
$\mathcal{O}_S$-대수이다.
\end{situation}

\noindent
이 절에서는 동차 스펙트럼들을 붙여서 스킴의 사상
$$
\underline{\text{Proj}}_S(\mathcal{A}) \longrightarrow S
$$
을 구성하는 방법을 개략적으로 설명한다. 붙이는 대상은
$\text{Proj}(\Gamma(U, \mathcal{A}))$들이고, 여기서 $U$는 $S$의 아핀
열린집합들을 달린다. 먼저 Lemma \ref{constructions-lemma-relative-glueing}에서와 같이,
아핀 열린집합 위에서 취한 $\mathcal{A}$의 값들의 동차 스펙트럼이 적절한
스킴들의 모음을 이룸을 보이자.

\begin{lemma}
\label{constructions-lemma-proj-inclusion}
Situation \ref{constructions-situation-relative-proj}의 상황이라 하자. $U \subset U' \subset S$가
아핀 열린집합들이라고 가정하자. $A = \mathcal{A}(U)$ 및
$A' = \mathcal{A}(U')$로 두자. 등급 환의 사상 $A' \to A$는 사상
$r : \text{Proj}(A) \to \text{Proj}(A')$을 유도하고, 도식
$$
\xymatrix{
\text{Proj}(A) \ar[r] \ar[d] &
\text{Proj}(A') \ar[d] \\
U \ar[r] &
U'
}
$$
은 올곱이다. 또한 곱셈 사상들과 호환되는 표준 동형
$\theta : r^*\mathcal{O}_{\text{Proj}(A')}(n) \to
\mathcal{O}_{\text{Proj}(A)}(n)$이 존재한다.
\end{lemma}

\begin{proof}
$R = \mathcal{O}_S(U)$ 및 $R' = \mathcal{O}_S(U')$로 두자.
사상 $R \otimes_{R'} A' \to A$가 동형임에 유의하자. 이는
$\mathcal{A}$가 준연접이기 때문이다(예를 들어 Schemes, Lemma
\ref{schemes-lemma-widetilde-pullback}을 보라). 따라서 Lemma
\ref{constructions-lemma-base-change-map-proj}로부터 보조정리가 따른다.
\end{proof}

\noindent
특히 보조정리의 사상 $\text{Proj}(A) \to \text{Proj}(A')$은 열린
몰입이다.

\begin{lemma}
\label{constructions-lemma-transitive-proj}
Situation \ref{constructions-situation-relative-proj}의 상황이라 하자.
$U \subset U' \subset U'' \subset S$가 아핀 열린집합들이라고 가정하자.
$A = \mathcal{A}(U)$, $A' = \mathcal{A}(U')$ 및
$A'' = \mathcal{A}(U'')$로 두자. Lemma \ref{constructions-lemma-proj-inclusion}의 사상
$r : \text{Proj}(A) \to \text{Proj}(A')$ 및
$r' : \text{Proj}(A') \to \text{Proj}(A'')$의 합성은 Lemma
\ref{constructions-lemma-proj-inclusion}의 사상
$r'' : \text{Proj}(A) \to \text{Proj}(A'')$이다. 동형 $\theta$에 대해서도
비슷한 명제가 성립한다.
\end{lemma}

\begin{proof}
사상 $A'' \to A$가 $A'' \to A'$와 $A' \to A$의 합성이므로, 이는 Lemma
\ref{constructions-lemma-morphism-proj-transitive}에서 따른다.
\end{proof}

\begin{lemma}
\label{constructions-lemma-glue-relative-proj}
Situation \ref{constructions-situation-relative-proj}의 상황이라 하자. 다음 성질들을 갖는
스킴의 사상
$$
\pi : \underline{\text{Proj}}_S(\mathcal{A}) \longrightarrow S
$$
이 존재한다.
\begin{enumerate}
\item 모든 아핀 열린집합 $U \subset S$에 대해 동형
$i_U : \pi^{-1}(U) \to \text{Proj}(A)$가 존재하며, 여기서
$A = \mathcal{A}(U)$이다.
\item 아핀 열린집합 $U \subset U' \subset S$에 대해 합성
$$
\xymatrix{
\text{Proj}(A) \ar[r]^{i_U^{-1}} &
\pi^{-1}(U) \ar[rr]^{inclusion} & &
\pi^{-1}(U') \ar[r]^{i_{U'}} &
\text{Proj}(A')
}
$$
은 $A = \mathcal{A}(U)$, $A' = \mathcal{A}(U')$일 때 위의 Lemma
\ref{constructions-lemma-proj-inclusion}의 열린 몰입이다.
\end{enumerate}
\end{lemma}

\begin{proof}
Lemmas \ref{constructions-lemma-relative-glueing}, \ref{constructions-lemma-proj-inclusion} 및
\ref{constructions-lemma-transitive-proj}에서 곧바로 따른다.
\end{proof}

\begin{lemma}
\label{constructions-lemma-glue-relative-proj-twists}
Situation \ref{constructions-situation-relative-proj}의 상황이라 하자. Lemma
\ref{constructions-lemma-glue-relative-proj}의 사상
$\pi : \underline{\text{Proj}}_S(\mathcal{A}) \to S$에는 다음 추가 구조가
주어진다. 준연접 $\mathbf{Z}$-등급
$\mathcal{O}_{\underline{\text{Proj}}_S(\mathcal{A})}$-대수의 층
$\bigoplus\nolimits_{n \in \mathbf{Z}}
\mathcal{O}_{\underline{\text{Proj}}_S(\mathcal{A})}(n)$과 등급
$\mathcal{O}_S$-대수의 사상
$$
\psi :
\mathcal{A}
\longrightarrow
\bigoplus\nolimits_{n \geq 0}
\pi_*\left(\mathcal{O}_{\underline{\text{Proj}}_S(\mathcal{A})}(n)\right)
$$
이 존재하고, 다음 성질로 유일하게 결정된다. 모든 아핀 열린집합 $U \subset S$와
$A = \mathcal{A}(U)$에 대해 $\mathbf{Z}$-등급
$\mathcal{O}_{\pi^{-1}(U)}$-대수의 동형
$$
\theta_U :
i_U^*\left(
\bigoplus\nolimits_{n \in \mathbf{Z}} \mathcal{O}_{\text{Proj}(A)}(n)
\right)
\longrightarrow
\left(
\bigoplus\nolimits_{n \in \mathbf{Z}}
\mathcal{O}_{\underline{\text{Proj}}_S(\mathcal{A})}(n)
\right)|_{\pi^{-1}(U)}
$$
이 존재하여
$$
\xymatrix{
A_n
\ar[rr]_\psi
\ar[dr]_-{(\ref{constructions-equation-global-sections})}
& &
\Gamma(\pi^{-1}(U),
\mathcal{O}_{\underline{\text{Proj}}_S(\mathcal{A})}(n)) \\
&
\Gamma(\text{Proj}(A),
\mathcal{O}_{\text{Proj}(A)}(n))
\ar[ru]_-{\theta_U}
&
}
$$
이 가환한다.
\end{lemma}

\begin{proof}
Lemma \ref{constructions-lemma-relative-glueing-sheaves}를 사용하여 $\mathbf{Z}$-등급 대수의
층 $\bigoplus_{n \in \mathbf{Z}} \mathcal{O}_{\text{Proj}(A)}(n)$들을
붙이겠다. 여기서 $A = \mathcal{A}(U)$이고 $U \subset S$는 아핀 열린집합이며,
붙이기는 스킴 $\underline{\text{Proj}}_S(\mathcal{A})$ 위에서 수행한다.
이에 필요한 자료는 Lemma \ref{constructions-lemma-proj-inclusion}에서 구성했고,
Lemma \ref{constructions-lemma-relative-glueing-sheaves}의 조건 (d)는 Lemma
\ref{constructions-lemma-transitive-proj}에서 확인하였다. 따라서 $\mathbf{Z}$-등급
$\mathcal{O}_{\underline{\text{Proj}}_S(\mathcal{A})}$-대수의 층
$\bigoplus_{n \in \mathbf{Z}}
\mathcal{O}_{\underline{\text{Proj}}_S(\mathcal{A})}(n)$과 동형
$\theta_U$를 얻는다. 이는 모든 $U \subset S$ 아핀 열린집합 및 모든
$n \in \mathbf{Z}$에 대해 주어진다. 모든 아핀 열린집합 $U \subset S$와
$A = \mathcal{A}(U)$에 대해 사상
$A \to \Gamma(\text{Proj}(A),
\bigoplus_{n \geq 0} \mathcal{O}_{\text{Proj}(A)}(n))$이 있다. 따라서
상대 붙이기의 함자성에 의해 사상 $\psi$가 존재한다. Remark
\ref{constructions-remark-relative-glueing-functorial}을 보라. 보조정리의 도식은 구성상
가환한다. 이는 $\mathbf{Z}$-등급
$\mathcal{O}_{\underline{\text{Proj}}_S(\mathcal{A})}$-대수의 층
$\bigoplus \mathcal{O}_{\underline{\text{Proj}}_S(\mathcal{A})}(n)$을
특징짓는다. 실제로 Lemma \ref{constructions-lemma-morphism-proj}의 증명은 이 도식들이
가환한다는 것이 사상 $\theta_U$들을 유일하게 결정함을 보인다. 일부 세부사항은
생략한다.
\end{proof}

\section{함자로서의 상대 Proj}
\label{constructions-section-relative-proj}

\noindent
Situation \ref{constructions-situation-relative-proj}의 상황에 놓여 있다고 하자. 따라서 $S$는
스킴이고,
$\mathcal{A} = \bigoplus_{d \geq 0} \mathcal{A}_d$는 준연접 등급
$\mathcal{O}_S$-대수이다. 이 절에서는 상대 동차 스펙트럼이 표현할 함자를
구성함으로써 $\text{Proj}$의 구성을 상대화한다. 그 결과, $S$의 아핀
열린집합 위에서는 등급 환의 동차 스펙트럼처럼 보이는 스킴의 사상
$$
\underline{\text{Proj}}_S(\mathcal{A}) \longrightarrow S
$$
을 구성할 것이다. 논의는 Section \ref{constructions-section-spec}의 상대 스펙트럼에 관한
논의를 본뜬다. $\text{Proj}(A)$ 꼴의 스킴들을 붙이는 더 쉬운 방법은
$A = \Gamma(U, \mathcal{A})$이고 $U \subset S$가 아핀 열린집합인 경우에
Section \ref{constructions-section-relative-proj-via-glueing}에서 설명했으며, 이 절의
결과가 그 결과와 동형임을 보일 것이다.

\medskip\noindent
당분간 정수 $d \geq 1$을 고정하자. Algebra, Section
\ref{algebra-section-graded}의 표기와 비슷하게
$\mathcal{A}^{(d)} = \bigoplus_{n \geq 0} \mathcal{A}_{nd}$로 쓴다.
$T$를 스킴이라 하자.
{\it 사중항 $(d, f : T \to S, \mathcal{L}, \psi)$들을 $T$ 위에서} 생각하자.
여기서
\begin{enumerate}
\item $d$는 위에서 고정한 정수이고,
\item $f : T \to S$는 스킴의 사상이며,
\item $\mathcal{L}$은 가역 $\mathcal{O}_T$-가군이고,
\item
$\psi : f^*\mathcal{A}^{(d)} \to \bigoplus_{n \geq 0}\mathcal{L}^{\otimes n}$는 등급
$\mathcal{O}_T$-대수의 준동형이며
$f^*\mathcal{A}_d \to \mathcal{L}$은 전사이다.
\end{enumerate}
사상 $h : T' \to T$와 사중항 $(d, f, \mathcal{L}, \psi)$가 $T$ 위에
주어지면, 이를 사중항 $(d, f \circ h, h^*\mathcal{L}, h^*\psi)$로
$T'$ 위에 당겨올 수 있다. 두 사중항 $(d, f, \mathcal{L}, \psi)$ 및
$(d, f', \mathcal{L}', \psi')$가 $T$ 위에 있고 같은 정수 $d$를 갖는다고
하자. $f = f'$이고 동형 $\beta : \mathcal{L} \to \mathcal{L}'$이 존재하여
$\beta \circ \psi = \psi'$이면, 여기서 이 등식은 등급
$\mathcal{O}_T$-대수 사상
$f^*\mathcal{A}^{(d)} \to \bigoplus_{n \geq 0} (\mathcal{L}')^{\otimes n}$
사이의 등식이며, 이 두 사중항을 {\it 강하게 동치}라고 한다.

\medskip\noindent
각 정수 $d \geq 1$에 대해 위에서 정의한 당겨올림을 사용하여
\begin{eqnarray*}
F_d : \Sch^{opp} & \longrightarrow & \textit{Sets}, \\
T & \longmapsto &
\{\text{strict equivalence classes of }
(d, f : T \to S, \mathcal{L}, \psi)
\text{ as above}\}
\end{eqnarray*}
로 정의한다.

\begin{lemma}
\label{constructions-lemma-proj-base-change}
Situation \ref{constructions-situation-relative-proj}의 상황이라 하자. $d \geq 1$이라 하자.
$F_d$를 위의 $(S, \mathcal{A})$에 결부된 함자라 하자. $g : S' \to S$를
스킴의 사상이라 하고 $\mathcal{A}' = g^*\mathcal{A}$로 두자. $F_d'$를 위의
$(S', \mathcal{A}')$에 결부된 함자라 하자. 그러면 함자들의 표준 동형
$$
F'_d \cong h_{S'} \times_{h_S} F_d
$$
이 존재한다.
\end{lemma}

\begin{proof}
사중항
$(d, f' : T \to S', \mathcal{L}',
\psi' : (f')^*(\mathcal{A}')^{(d)} \to
\bigoplus_{n \geq 0} (\mathcal{L}')^{\otimes n})$은 사중항
$(d, f, \mathcal{L},
\psi : f^*\mathcal{A}^{(d)} \to
\bigoplus_{n \geq 0} \mathcal{L}^{\otimes n})$과 $f$의 분해
$f = g \circ f'$를 함께 준 것과 같다. 구체적으로 이 대응은
$f = g \circ f'$, $\mathcal{L} = \mathcal{L}'$ 및
$\psi = \psi'$이고, 여기서는 식별
$(f')^*(\mathcal{A}')^{(d)} = (f')^*g^*(\mathcal{A}^{(d)}) =
f^*\mathcal{A}^{(d)}$을 사용한다. 따라서 보조정리가 성립한다.
\end{proof}

\begin{lemma}
\label{constructions-lemma-relative-proj-affine}
Situation \ref{constructions-situation-relative-proj}의 상황이라 하자. $F_d$를 위의
$(d, S, \mathcal{A})$에 결부된 함자라 하자. $S$가 아핀이면, $F_d$는 열린
부분스킴 $U_d$ (\ref{constructions-equation-Ud}), 즉 스킴
$\text{Proj}(\Gamma(S, \mathcal{A}))$의 그 열린 부분스킴으로 표현된다.
\end{lemma}

\begin{proof}
$S = \Spec(R)$ 및 $A = \Gamma(S, \mathcal{A})$로 쓰자. 그러면 $A$는 등급
$R$-대수이고 $\mathcal{A} = \widetilde A$이다. 보조정리를 증명하려면 함자
$F_d$를 Section \ref{constructions-section-morphisms-proj}에서 정의한 삼중항의 함자
$F_d^{triples}$와 식별해야 한다.

\medskip\noindent
$(d, f : T \to S, \mathcal{L}, \psi)$를 사중항이라 하자. $\psi$를
$\mathcal{O}_S$-가군 사상
$\mathcal{A}^{(d)} \to \bigoplus_{n \geq 0} f_*\mathcal{L}^{\otimes n}$으로
생각할 수 있다.
$\mathcal{A}^{(d)}$가 준연접이므로, 이는 등급 환의 $R$-선형 준동형
$A^{(d)} \to \Gamma(S, \bigoplus_{n \geq 0} f_*\mathcal{L}^{\otimes n})$과
같다.
명백히
$\Gamma(S, \bigoplus_{n \geq 0} f_*\mathcal{L}^{\otimes n}) =
\Gamma_*(T, \mathcal{L})$이다. 따라서 이 사중항에 삼중항
$(d, \mathcal{L}, \psi)$를 결부시킬 수 있다.

\medskip\noindent
거꾸로 $(d, \mathcal{L}, \psi)$를 삼중항이라 하자. 합성
$R \to A_0 \to \Gamma(T, \mathcal{O}_T)$은 사상
$f : T \to S = \Spec(R)$를 결정한다. Schemes, Lemma
\ref{schemes-lemma-morphism-into-affine}을 보라. 이 $f$를 택하면 사상
$A^{(d)} \to \Gamma(S, \bigoplus_{n \geq 0} f_*\mathcal{L}^{\otimes n})$은
$R$-선형이고, 따라서 $\psi$에 대응하며 이를 사중항
$(d, f : T \to S, \mathcal{L}, \psi)$에 사용할 수 있다. 이것이 함자들의
동형 $F_d = F_d^{triples}$를 정한다는 확인은
생략한다.
\end{proof}

\begin{lemma}
\label{constructions-lemma-relative-proj-d}
Situation \ref{constructions-situation-relative-proj}의 상황이라 하자. 함자 $F_d$는
스킴으로 표현 가능하다.
\end{lemma}

\begin{proof}
Schemes, Lemma \ref{schemes-lemma-glue-functors}를 사용할 것이다.

\medskip\noindent
먼저 $F_d$가 자리스키 위상에 대한 층 성질을 만족함을 확인하자. 구체적으로
$T$가 스킴이고 $T = \bigcup_{i \in I} U_i$가 열린 덮개이며,
$(d, f_i, \mathcal{L}_i, \psi_i) \in F_d(U_i)$이고
$(d, f_i, \mathcal{L}_i, \psi_i)|_{U_i \cap U_j}$와
$(d, f_j, \mathcal{L}_j, \psi_j)|_{U_i \cap U_j}$가 강하게 동치라고
가정하자. 그러면 사상 $f_i : U_i \to S$들은 스킴의 사상
$f : T \to S$로 붙고 $f|_{I_i} = f_i$를 만족한다. Schemes, Section
\ref{schemes-section-glueing-schemes}를 보라. 따라서
$f_i^*\mathcal{A}^{(d)} = f^*\mathcal{A}^{(d)}|_{U_i}$이다. 또한 동형
$\beta_{ij} : \mathcal{L}_i|_{U_i \cap U_j} \to \mathcal{L}_j|_{U_i \cap U_j}$가
존재하여 $\beta_{ij} \circ \psi_i = \psi_j$를 $U_i \cap U_j$ 위에서
만족한다. 동형 $\beta_{ij}$는 사상
$f_i^*\mathcal{A}_d \to \mathcal{L}_i$가 전사이므로 이 조건으로 유일하게
결정된다는 점에 유의하자. 특히
$\beta_{jk} \circ \beta_{ij} = \beta_{ik}$가 $U_i \cap U_j \cap U_k$ 위에서
성립한다. 따라서 Sheaves,
Section \ref{sheaves-section-glueing-sheaves}에 의해 가역층
$\mathcal{L}_i$들은 가역 $\mathcal{O}_T$-가군 $\mathcal{L}$로 붙고,
사상 $\psi_i$들은 $\mathcal{O}_T$-대수의 사상
$\psi : f^*\mathcal{A}^{(d)} \to \bigoplus_{n \geq 0} \mathcal{L}^{\otimes n}$으로
붙는다. 이로써
$F_d$가 자리스키 위상에 대한 층 조건을 만족함을 증명하였다.

\medskip\noindent
$S = \bigcup_{i \in I} U_i$를 아핀 열린 덮개라 하자.
$F_{d, i} \subset F_d$를 쌍 $(f : T \to S, \varphi)$들 중
$f(T) \subset U_i$를 만족하는 것들로 이루어진 부분함자라 하자.

\medskip\noindent
각 $F_{d, i}$가 표현 가능함을 보여야 한다. Lemma
\ref{constructions-lemma-proj-base-change}에 의해 $F_{d, i}$는 $U_i$에 준연접 등급
$\mathcal{O}_{U_i}$-대수 $\mathcal{A}|_{U_i}$를 주어 얻는 함자와
식별되므로 그렇다. 따라서 Lemma
\ref{constructions-lemma-relative-proj-affine}에서 결과가 따른다.

\medskip\noindent
다음으로 $F_{d, i} \subset F_d$가 열린 몰입들로 표현됨을 보이자.
$(f : T \to S, \varphi) \in F_d(T)$라 하고 $V_i = f^{-1}(U_i)$를
생각하자. $F_{d, i}$의 정의에 의해 $a : T' \to T$가 주어졌을 때
$a^*(f, \varphi) \in F_{d, i}(T')$인 것과 $a(T') \subset V_i$인 것은
동치이다. 이것이 보일 명제였다.

\medskip\noindent
마지막으로 모음 $(F_{d, i})_{i \in I}$가 $F_d$를 덮음을 보여야 한다.
$(f : T \to S, \varphi) \in F_d(T)$라 하고
$V_i = f^{-1}(U_i)$를 생각하자. $S = \bigcup_{i \in I} U_i$가 $S$의
열린 덮개이므로 $T = \bigcup_{i \in I} V_i$는 $T$의 열린 덮개이다.
또한 $(f, \varphi)|_{V_i} \in F_{d, i}(V_i)$이다. 이로써 보조정리의
증명이 끝난다.
\end{proof}

\noindent
이제 Section \ref{constructions-section-morphisms-proj} 끝부분의 내용을 현재의 상대적 상황에서
되풀이하여 $\underline{\text{Proj}}_S(\mathcal{A})$로 표현되는 함자를 정의할
수 있다. 이를 위해 두 사중항
$(d, f : T \to S, \mathcal{L}, \psi)$ 및
$(d', f' : T \to S, \mathcal{L}', \psi')$ 사이의 동치 개념을 도입한다.
여기서 정수 $d, d'$의 값은 서로 다를 수도 있다.
구체적으로 $f = f'$이고 동형
$\beta : \mathcal{L}^{\otimes d'} \to (\mathcal{L}')^{\otimes d}$이
존재하여
$\beta \circ \psi|_{f^*\mathcal{A}^{(dd')}} = \psi'|_{f^*\mathcal{A}^{(dd')}}$이면
이들을 {\it 동치}라고 한다. 다음
보조정리는 이것이 동치관계를 정의함을 뜻한다. (이는 완전히 자명하지는 않다.)

\begin{lemma}
\label{constructions-lemma-equivalent-relative}
Situation \ref{constructions-situation-relative-proj}의 상황이라 하자. $T$를 스킴이라 하자.
$(d, f, \mathcal{L}, \psi)$, $(d', f', \mathcal{L}', \psi')$를 $T$ 위의 두
사중항이라 하자. 다음은 동치이다.
\begin{enumerate}
\item $m = \text{lcm}(d, d')$라 하고 $m = ad = a'd'$로 쓰자.
$f = f'$이고 동형
$\beta : \mathcal{L}^{\otimes a} \to (\mathcal{L}')^{\otimes a'}$이
존재하여
$\beta \circ \psi|_{f^*\mathcal{A}^{(m)}}$과
$\psi'|_{f^*\mathcal{A}^{(m)}}$이 등급 환 사상
$f^*\mathcal{A}^{(m)} \to \bigoplus_{n \geq 0} (\mathcal{L}')^{\otimes mn}$으로서
일치한다.
\item 사중항 $(d, f, \mathcal{L}, \psi)$와
$(d', f', \mathcal{L}', \psi')$는 동치이다.
\item $f = f'$이고, 어떤 양의 정수 $m = ad = a'd'$에 대해 동형
$\beta : \mathcal{L}^{\otimes a} \to (\mathcal{L}')^{\otimes a'}$이
존재하여
$\beta \circ \psi|_{f^*\mathcal{A}^{(m)}}$과
$\psi'|_{f^*\mathcal{A}^{(m)}}$이 등급 환 사상
$f^*\mathcal{A}^{(m)} \to \bigoplus_{n \geq 0} (\mathcal{L}')^{\otimes mn}$으로서
일치한다.
\end{enumerate}
\end{lemma}

\begin{proof}
더 나누어지는 차수들과 가역층들의 거듭제곱에 제한하면 되므로, (1)이 (2)를
함의하고 (2)가 (3)을 함의함은 명백하다. 어떤 정수 $m = ad = a'd'$에 대해
(3)을 가정하자. $m_0 = \text{lcm}(d, d')$라 하고
$m_0 = a_0d = a'_0d'$로 쓰자. (3)의 성질을 갖는 동형
$\beta : \mathcal{L}^{\otimes a} \to (\mathcal{L}')^{\otimes a'}$이
주어져 있다. 같은 성질을 갖는 동형
$\beta_0 : \mathcal{L}^{\otimes a_0} \to (\mathcal{L}')^{\otimes a'_0}$을
찾고자 한다. 가정에 의해 사상
$\psi : f^*\mathcal{A}_d \to \mathcal{L}$ 및
$\psi' : (f')^*\mathcal{A}_{d'} \to \mathcal{L}'$이 전사이므로, 사상
$\psi : f^*\mathcal{A}_{m_0} \to \mathcal{L}^{\otimes a_0}$ 및
$\psi' : (f')^*\mathcal{A}_{m_0} \to (\mathcal{L}')^{\otimes a_0}$도
전사이다. 따라서 $\beta_0$가 존재한다면
조건 $\beta_0 \circ \psi = \psi'$로 유일하게 결정된다. 이는 $T$ 위에서
국소적으로 작업해도 됨을 뜻한다. 그러므로
$f = f' : T \to S$가 어떤 아핀 열린집합으로 사상된다고 가정할 수 있고,
달리 말하면 $S$가 아핀이라고 가정할 수 있다. 이 경우 결과는 삼중항에 대한
대응하는 결과(Lemma \ref{constructions-lemma-equivalent})와 아핀 밑의 경우에 삼중항과
사중항이 대응한다는 사실(Lemma \ref{constructions-lemma-relative-proj-affine}의 증명 참조)
에서 따른다.
\end{proof}

\noindent
$d' = ad$라고 가정하자. 함자의 변환 $F_d \to F_{d'}$을 생각하자. 이
변환은 사중항 $(d, f, \mathcal{L}, \psi)$가 $T$ 위에 주어지면 사중항
$(d', f, \mathcal{L}^{\otimes a}, \psi|_{f^*\mathcal{A}^{(d')}})$을
결부시킨다. Lemma
\ref{constructions-lemma-equivalent-relative}의 함의 중 하나에 따르면 이 변환
$F_d \to F_{d'}$은 단사이다! 준콤팩트 스킴 $T$에 대해
$$
F(T) = \bigcup\nolimits_{d \in \mathbf{N}} F_d(T)
$$
로 정의하되, 전이 사상은 위에서 설명한 것을 사용한다. 이는 준콤팩트 스킴의
범주에서 집합의 범주로 가는 반변 함자를 명백히 정의한다. 일반 스킴 $T$에
대해서는
$$
F(T)
=
\lim_{V \subset T\text{ quasi-compact open}} F(V).
$$
로 정의한다. 다시 말해 원소 $\xi$가 $F(T)$에 속한다는 것은 원소
$\xi_V \in F(V)$들의 호환되는 선택계에 대응하며, 여기서 $V$는 $T$의
준콤팩트 열린집합들을 달린다. 당겨올림 사상 $F(T) \to F(T')$의 정의는
스킴의 사상 $T' \to T$에 대해 생략한다. 이로써 함자
\begin{equation}
\label{constructions-equation-proj}
F : \Sch^{opp} \longrightarrow \textit{Sets}
\end{equation}
를 정의하였다.

\begin{lemma}
\label{constructions-lemma-relative-proj}
Situation \ref{constructions-situation-relative-proj}의 상황이라 하자. 위의 함자 $F$는
스킴으로 표현 가능하다.
\end{lemma}

\begin{proof}
$U_d \to S$를 위에서 정의한 함자 $F_d$를 표현하는 스킴이라 하자. 보편
대상을
$\mathcal{L}_d$,
$\psi^d : \pi_d^*\mathcal{A}^{(d)} \to
\bigoplus_{n \geq 0} \mathcal{L}_d^{\otimes n}$로 쓰자. $d | d'$이면
사중항
$(d', \pi_d, \mathcal{L}_d^{\otimes d'/d}, \psi^d|_{\mathcal{A}^{(d')}})$을
생각할 수 있고, 이는 표준 사상 $U_d \to U_{d'}$을 $S$ 위에서 결정한다.
구성상 이 사상은 위에서 정의한 함자의 변환
$F_d \to F_{d'}$에 대응한다.

\medskip\noindent
모든 아핀 열린집합 $\Spec(R) = V \subset S$에 대해
$A = \Gamma(V, \mathcal{A})$로 두면, 밑변환 $U_{d, V}$와
$\text{Proj}(A)$의 대응하는 열린 부분스킴 사이의 표준 식별이 있다. Lemma
\ref{constructions-lemma-relative-proj-affine}을 보라. 또한 위에서 구성한 사상
$U_{d, V} \to U_{d', V}$는 $\text{Proj}(A)$의 열린집합들의 포함에
대응한다. 따라서 $U_d \to U_{d'}$는 열린 몰입이다.

\medskip\noindent
이를 통해 $X$를 구성할 수 있다. 즉 스킴 $U_d$들을 열린 몰입
$U_d \to U_{d'}$을 따라 붙인다. 기술적으로는 수열
$d_1 | d_2 | d_3 | \ldots$을 택하되 모든 양의 정수가 어느 $d_i$를
나누게 하고, 위의 열린 몰입을 사용하여 단순히
$X = \bigcup U_{d_i}$로 두는 것이 편리하다. 이제 $X$가 함자 $F$를
표현함을 증명하는 것은 간단하다.
\end{proof}

\begin{lemma}
\label{constructions-lemma-glueing-gives-functor-proj}
Situation \ref{constructions-situation-relative-proj}의 상황이라 하자. Lemma
\ref{constructions-lemma-glue-relative-proj}에서 구성한 스킴
$\pi : \underline{\text{Proj}}_S(\mathcal{A}) \to S$와 함자 $F$를
표현하는 스킴은 $S$ 위의 스킴으로서 표준적으로 동형이다.
\end{lemma}

\begin{proof}
$X$를 함자 $F$를 표현하는 스킴이라 하자. $X$는 $S$ 위의 스킴임에
유의하자. 실제로 함자 $F$에는 자연변환 $F \to h_S$가 주어져 있다.
$Y = \underline{\text{Proj}}_S(\mathcal{A})$로 쓰자.
$X \cong Y$가 $S$-스킴의 동형임을 보여야 한다. 두 가지 논증을 제시한다.

\medskip\noindent
첫 번째 논증은 Lemma \ref{constructions-lemma-relative-proj}의 증명에서 $X$를
스킴 $U_d$, 즉 함자 $F_d$를 표현하는 스킴들의 합집합으로 구성한 것을
사용한다.
$S$의 각 아핀 열린집합 위에서 $X$를 그 열린집합 위의 $\mathcal{A}$의
단면들의 동차 스펙트럼과 식별할 수 있다. 실제로 각 $U_d$의 열린집합에
대해 이 사실이 성립한다.
또한 이 식별들은 더 작은 아핀 열린집합으로 다시 제한하는 것과 호환된다.
한편 $Y$는 이 동차 스펙트럼들을 붙여 구성하였다. 따라서 이 동형들을 붙여
원하는 대로 $X$와 $\underline{\text{Proj}}_S(\mathcal{A})$ 사이의 동형을
얻는다. 세부사항은 생략한다.

\medskip\noindent
두 번째 논증은 다음과 같다. Lemma
\ref{constructions-lemma-glue-relative-proj-twists}에 따르면 $Y$ 위에서 등급 대수의 사상
$$
\psi : \pi^*\mathcal{A}
\longrightarrow
\bigoplus\nolimits_{n \geq 0} \mathcal{O}_Y(n)
$$
이 존재하고, 이는 $S$의 아핀 열린집합 위의 단면에서
(\ref{constructions-equation-global-sections})와 일치한다. 따라서 모든 $y \in Y$에 대해
열린 근방 $V \subset Y$가 $y$를 포함하고 정수 $d \geq 1$이 존재하여,
$d | n$일 때
층 $\mathcal{O}_Y(n)|_V$가 가역이고 곱셈 사상
$\mathcal{O}_Y(n)|_V \otimes_{\mathcal{O}_V} \mathcal{O}_Y(m)|_V
\to \mathcal{O}_Y(n + m)|_V$가 동형이다. 따라서 $\psi$를 층
$\pi^*\mathcal{A}^{(d)}|_V$에 제한하면 $F_d(V)$의 원소를 얻는다.
열린집합 $V$들이 $Y$를 덮으므로 “$\psi$”는 $F(Y)$의 원소를 정한다.
따라서 표준 사상 $Y \to X$를 $S$ 위에서 얻는다. 이 구성이 완전히
표준적이므로 그것이 동형임을 보이기 위해 $S$ 위에서 국소적으로 작업해도
된다. 따라서 $S$가 아핀인 경우로 환원되고, 이 경우 결과는 명백하다.
\end{proof}

\begin{definition}
\label{constructions-definition-relative-proj}
$S$를 스킴이라 하자. $\mathcal{A}$를 등급 $\mathcal{O}_S$-대수의 준연접
층이라 하자. $\mathcal{A}$의 $S$ 위 {\it 상대 동차 스펙트럼}, 또는
$\mathcal{A}$의 $S$ 위 {\it 동차 스펙트럼}, 또는
$\mathcal{A}$의 $S$ 위 {\it 상대 Proj}란 Lemma
\ref{constructions-lemma-glue-relative-proj}에서 구성되어 함자 $F$
(\ref{constructions-equation-proj})를 표현하는 스킴이다. Lemma
\ref{constructions-lemma-glueing-gives-functor-proj}를 보라. 이를
$\pi : \underline{\text{Proj}}_S(\mathcal{A}) \to S$로 나타낸다.
\end{definition}

\noindent
상대 Proj에는 $\mathbf{Z}$-등급 대수의 준연접층
$\bigoplus_{n \in \mathbf{Z}}
\mathcal{O}_{\underline{\text{Proj}}_S(\mathcal{A})}(n)$
(구조층의 꼬임들)과 등급 대수의 “보편” 준동형
$$
\psi_{univ} :
\mathcal{A}
\longrightarrow
\pi_*\left(
\bigoplus\nolimits_{n \geq 0}
\mathcal{O}_{\underline{\text{Proj}}_S(\mathcal{A})}(n)
\right)
$$
이 주어진다. Lemma \ref{constructions-lemma-glue-relative-proj-twists}를 보라. 원한다면
이를 준동형
$$
\psi_{univ} :
\pi^*\mathcal{A}
\longrightarrow
\bigoplus\nolimits_{n \geq 0}
\mathcal{O}_{\underline{\text{Proj}}_S(\mathcal{A})}(n)
$$
으로 생각할 수도 있다. 다음 보조정리는 이 대상의 보편성을 나타낸다.

\begin{lemma}
\label{constructions-lemma-tie-up-psi}
Situation \ref{constructions-situation-relative-proj}의 상황이라 하자.
$(f : T \to S, d, \mathcal{L}, \psi)$를 사중항이라 하자. 결부된 사상을
$r_{d, \mathcal{L}, \psi} : T \to \underline{\text{Proj}}_S(\mathcal{A})$로
쓰자. 이는 $S$-사상이다. $\mathbf{Z}$-등급
$\mathcal{O}_T$-대수의 동형
$$
\theta :
r_{d, \mathcal{L}, \psi}^*\left(
\bigoplus\nolimits_{n \in \mathbf{Z}}
\mathcal{O}_{\underline{\text{Proj}}_S(\mathcal{A})}(nd)
\right)
\longrightarrow
\bigoplus\nolimits_{n \in \mathbf{Z}} \mathcal{L}^{\otimes n}
$$
이 존재하여 다음 도식이 가환한다.
$$
\xymatrix{
\mathcal{A}^{(d)} \ar[rr]_-{\psi}
 \ar[rd]_-{\psi_{univ}} & &
f_*\left(
\bigoplus\nolimits_{n \in \mathbf{Z}}
\mathcal{L}^{\otimes n}
\right) \\
 &
\pi_*\left(
\bigoplus\nolimits_{n \geq 0}
\mathcal{O}_{\underline{\text{Proj}}_S(\mathcal{A})}(nd)
\right) \ar[ru]_\theta
}
$$
이 도식의 가환성은 $\theta$를 유일하게 결정한다.
\end{lemma}

\begin{proof}
사중항 $(f : T \to S, d, \mathcal{L}, \psi)$는 $F_d(T)$의 원소를
정한다는 점에 유의하자.
$U_d \subset \underline{\text{Proj}}_S(\mathcal{A})$를 층
$\mathcal{O}_{\underline{\text{Proj}}_S(\mathcal{A})}(d)$가 가역이고
$\psi_{univ} : \pi^*\mathcal{A}_d \to
\mathcal{O}_{\underline{\text{Proj}}_S(\mathcal{A})}(d)$의 상으로
생성되는 자취라 하자. $U_d$가 함자 $F_d$를 표현함을 상기하자. Lemma
\ref{constructions-lemma-relative-proj}의 증명을 보라. 따라서 사중항
$(U_d \to S, d, \mathcal{O}_{U_d}(d), \psi_{univ}|_{\mathcal{A}^{(d)}})$이
보편족, 즉 $F_d(U_d)$의 표현 대상임을
보이면 결과가 따른다. 다음 두 이유로 $S$의 아핀 열린집합에 제한한 뒤 이를
확인해도 된다. (a) 함자 $F_d$의 구성은 밑변환과 교환하고(Lemma
\ref{constructions-lemma-proj-base-change} 참조), (b) 쌍
$(\bigoplus_{n \in \mathbf{Z}}
\mathcal{O}_{\underline{\text{Proj}}_S(\mathcal{A})}(n),
\psi_{univ})$은 $S$의 아핀 열린집합들 위에서 붙이기로 구성된다(Lemma
\ref{constructions-lemma-glue-relative-proj-twists} 참조). 따라서 $S$가 아핀이라고
가정할 수 있다. 이 경우 사중항의 함자 $F_d$와 삼중항의 함자 $F_d$는
일치하고(Lemma \ref{constructions-lemma-relative-proj-affine}의 증명 참조), 또한 Lemma
\ref{constructions-lemma-proj-functor-strict}에 따르면
$(d, \mathcal{O}_{U_d}(d), \psi^d)$는 $U_d$ 위의 보편 삼중항이다. Lemma
\ref{constructions-lemma-relative-proj-affine}의 증명에 나온 식별들을 거꾸로 따라가면
$(U_d \to S, d, \mathcal{O}_{U_d}(d), \psi_{univ}|_{\mathcal{A}^{(d)}})$이
원하는 보편 사중항임을 알 수 있다.
\end{proof}

\begin{lemma}
\label{constructions-lemma-relative-proj-separated}
$S$를 스킴이라 하고 $\mathcal{A}$를 등급 $\mathcal{O}_S$-대수의 준연접층이라
하자. 사상 $\pi : \underline{\text{Proj}}_S(\mathcal{A}) \to S$는
분리사상이다.
\end{lemma}

\begin{proof}
사상이 분리사상임을 증명할 때 밑 위에서 국소적으로 작업해도 된다. Schemes,
Section \ref{schemes-section-separation-axioms}를 보라. 구성에 의해
$\underline{\text{Proj}}_S(\mathcal{A})$는 모든 아핀
$U \subset S$ 위에서 $\text{Proj}(A)$와 동형이며, 여기서
$A = \mathcal{A}(U)$이다. Lemma \ref{constructions-lemma-proj-separated}에 의해
$\text{Proj}(A)$가 분리임을 안다. 따라서 원하는 대로
$\text{Proj}(A) \to U$는 분리사상이다(Schemes, Lemma
\ref{schemes-lemma-compose-after-separated} 참조).
\end{proof}

\begin{lemma}
\label{constructions-lemma-relative-proj-base-change}
$S$를 스킴이라 하고 $\mathcal{A}$를 등급 $\mathcal{O}_S$-대수의 준연접층이라
하자. $g : S' \to S$를 임의의 스킴의 사상이라 하자. 그러면 표준 동형
$$
r :
\underline{\text{Proj}}_{S'}(g^*\mathcal{A})
\longrightarrow
S' \times_S \underline{\text{Proj}}_S(\mathcal{A})
$$
과 이에 대응하는 $\mathbf{Z}$-등급
$\mathcal{O}_{\underline{\text{Proj}}_{S'}(g^*\mathcal{A})}$-대수의 동형
$$
\theta :
r^*\text{pr}_2^*\left(\bigoplus\nolimits_{d \in \mathbf{Z}}
\mathcal{O}_{\underline{\text{Proj}}_S(\mathcal{A})}(d)\right)
\longrightarrow
\bigoplus\nolimits_{d \in \mathbf{Z}}
\mathcal{O}_{\underline{\text{Proj}}_{S'}(g^*\mathcal{A})}(d)
$$
이 존재한다.
\end{lemma}

\begin{proof}
이는 Lemma \ref{constructions-lemma-proj-base-change}와
$\underline{\text{Proj}}_S(\mathcal{A})$를 스킴 $U_d$, 즉 함자 $F_d$를
표현하는 스킴들의 합집합으로 구성한 Lemma
\ref{constructions-lemma-relative-proj}에서 따른다.
붙이기를 통한 상대 Proj의 구성으로 표현하면 이 동형은 Lemma
\ref{constructions-lemma-base-change-map-proj}에서 구성한 동형들로 주어지며, 이들이 동형
$\theta$를 준다. 일부 세부사항은 생략한다.
\end{proof}

\begin{lemma}
\label{constructions-lemma-apply-relative}
$S$를 스킴이라 하자. $\mathcal{A}$를 등급 $\mathcal{O}_S$-가군의
준연접층이라 하되, $\mathcal{A}_0$-대수로서 $\mathcal{A}_1$에 의해
생성된다고 하자. 이 경우 스킴
$X = \underline{\text{Proj}}_S(\mathcal{A})$는 함자 $F_1$을 표현한다.
이 함자는 스킴 $f : T \to S$가 $S$ 위에 주어지면 다음 쌍
$(\mathcal{L}, \psi)$들의 집합을 결부시킨다.
\begin{enumerate}
\item $\mathcal{L}$은 가역 $\mathcal{O}_T$-가군이고,
\item $\psi : f^*\mathcal{A} \to \bigoplus_{n \geq 0} \mathcal{L}^{\otimes n}$는 등급
$\mathcal{O}_T$-대수 준동형이며 $f^*\mathcal{A}_1 \to \mathcal{L}$은
전사이다.
\end{enumerate}
여기서 위와 같은 강한 동치까지 취한다. 또한 이 경우 모든 준연접층
$\mathcal{O}_{\underline{\text{Proj}}(\mathcal{A})}(n)$은 가역
$\mathcal{O}_{\underline{\text{Proj}}(\mathcal{A})}$-가군이고 곱셈 사상은
동형
$
\mathcal{O}_{\underline{\text{Proj}}(\mathcal{A})}(n)
\otimes_{\mathcal{O}_{\underline{\text{Proj}}(\mathcal{A})}}
\mathcal{O}_{\underline{\text{Proj}}(\mathcal{A})}(m) =
\mathcal{O}_{\underline{\text{Proj}}(\mathcal{A})}(n + m)$
을 유도한다.
\end{lemma}

\begin{proof}
보조정리의 가정 아래 층
$\mathcal{O}_{\underline{\text{Proj}}(\mathcal{A})}(n)$은 가역이고,
$S$의 아핀 열린집합들 위에서 Lemma \ref{constructions-lemma-relative-proj}와 Lemma
\ref{constructions-lemma-apply}에 의해 곱셈 사상들이 동형이다. 따라서 $X$는 실제로
함자 $F_1$을 표현한다. Lemma \ref{constructions-lemma-relative-proj}의 증명을 보라.
\end{proof}

\section{상대 Proj 위의 준연접층}
\label{constructions-section-quasi-coherent-relative-proj}

\noindent
상대적 상황에서 등급 가군을 다루는 방법을 간단히 논의한다.

\medskip\noindent
Situation \ref{constructions-situation-relative-proj}의 상황에 놓여 있다고 하자. 따라서
$S$는 스킴이고 $\mathcal{A}$는 준연접 등급 $\mathcal{O}_S$-대수이다.
$\mathcal{M} = \bigoplus_{n \in \mathbf{Z}} \mathcal{M}_n$을 등급
$\mathcal{A}$-가군이라 하되, $\mathcal{O}_S$-가군으로서 준연접이라고 하자.
$\underline{\text{Proj}}_S(\mathcal{A})$ 위의 이에 결부된 가군의 준연접층을
설명할 것이다. 먼저 이 층의 값을 상대 Proj로 사상되는 스킴 $T$ 위에서
설명한다.

\medskip\noindent
$T$를 스킴이라 하자. 사중항 $(d, f : T \to S, \mathcal{L}, \psi)$가
$T$ 위에 주어졌다고 하자. Section \ref{constructions-section-relative-proj}에서와 같다.
준연접층 $\widetilde{\mathcal{M}}_T$를 $\mathcal{O}_T$-가군의 층으로서
다음과 같이 정의한다.
\begin{equation}
\label{constructions-equation-widetilde-M}
\widetilde{\mathcal{M}}_T =
\left(
f^*\mathcal{M}^{(d)}
\otimes_{f^*\mathcal{A}^{(d)}}
\left(\bigoplus\nolimits_{n \in \mathbf{Z}} \mathcal{L}^{\otimes n}\right)
\right)_0
\end{equation}
즉 $\widetilde{\mathcal{M}}_T$는 차수 $0$인 부분이다. 구체적으로 등급
$f^*\mathcal{A}^{(d)}$-가군 $\mathcal{M}^{(d)}$와
$\bigoplus\nolimits_{n \in \mathbf{Z}} \mathcal{L}^{\otimes n}$의
텐서곱에서 취한 부분이다. 표기에서는 나타내지 않았지만 층
$\widetilde{\mathcal{M}}_T$는 사중항에 의존한다는 점에 유의하자. 이 구성은
임의의 사상 $g : T' \to T$가 주어졌을 때
$\widetilde{\mathcal{M}}_{T'} = g^*\widetilde{\mathcal{M}}_T$를 만족하는
좋은 성질을 갖는다. 여기서 $\widetilde{\mathcal{M}}_{T'}$는 당겨온 사중항
$(d, f \circ g, g^*\mathcal{L}, g^*\psi)$에 결부된 준연접층을 나타낸다.

\medskip\noindent
(\ref{constructions-equation-widetilde-M})의 모든 층이 준연접이므로, 아핀 열린집합
$\Spec(C) = V \subset T$가 아핀 열린집합
$\Spec(R) = U \subset S$로 사상되는 경우에 이 구성을 구체적으로 쓸 수 있다.
구체적으로 $\mathcal{A}|_U$가 등급 $R$-대수 $A$에 대응하고,
$\mathcal{M}|_U$가 등급 $A$-가군 $M$에 대응하며, $\mathcal{L}|_V$가 가역
$C$-가군 $L$에 대응한다고 하자. 사상 $\psi$는 등급 $R$-대수 사상
$\gamma : A^{(d)} \to \bigoplus_{n \geq 0} L^{\otimes n}$을 정한다.
($L$의 $C$ 위 텐서 거듭제곱이다.) 그러면
$(\widetilde{\mathcal{M}}_T)|_V$는 $C$-가군
$$
N_{R, C, A, M, \gamma} =
\left(
M^{(d)} \otimes_{A^{(d)}, \gamma}
\left(\bigoplus\nolimits_{n \in \mathbf{Z}} L^{\otimes n}\right)
\right)_0
$$
에 결부된 준연접층이다. 가정에 의해 $T$를 아핀 열린집합 $V$들로 덮되,
각각에 대해 어떤 $a \in A_d$가 존재하여 $\gamma(a) \in L$이
$C$-기저로서 가군 $L$을 생성하게 할 수도 있다. 그 경우
$N_{R, C, A, M, \gamma}$의 모든 원소는 순수 텐서들의 합
$\sum m_i \otimes \gamma(a)^{-n_i}$으로 쓸 수 있으며, 여기서
$m_i \in M_{n_id}$이다. 실제로 각 $m_i$에
$a$의 적당한 양의 거듭제곱을 곱하고 항들을 모으면
$N_{R, C, A, M, \gamma}$의 각 원소를
$m \otimes \gamma(a)^{-n}$ 꼴로 쓸 수 있다. 여기서 $m \in M_{nd}$이고
$n \gg 0$이다. 다시 말해 이 경우
$$
N_{R, C, A, M, \gamma} = M_{(a)} \otimes_{A_{(a)}} C
$$
이다. 여기서 사상 $A_{(a)} \to C$는
$x/a^n \mapsto \gamma(x)/\gamma(a)^n$인 사상이다. 달리 말하면 이는
$\widetilde{M}$의 $D_{+}(a) \subset \text{Proj}(A)$ 위 값을
$\Spec(C)$로 당겨온 것이다. 이 당겨올림은 사상
$\Spec(C) \to D_{+}(a)$을 통해 취하며, 이 사상은 $\gamma$에서 온다.

\begin{lemma}
\label{constructions-lemma-relative-proj-modules}
Situation \ref{constructions-situation-relative-proj}의 상황이라 하자. 준연접 등급
$\mathcal{A}$-가군층 $\mathcal{M}$이 $S$ 위에 주어질 때마다 이에
표준적으로 결부되고 다음 성질들을
갖는 $\mathcal{O}_{\underline{\text{Proj}}_S(\mathcal{A})}$-가군층
$\widetilde{\mathcal{M}}$이 존재한다.
\begin{enumerate}
\item 스킴 $T$와 사중항
$(T \to S, d, \mathcal{L}, \psi)$가 $T$ 위에 주어지고, 이것이 사상
$h : T \to \underline{\text{Proj}}_S(\mathcal{A})$에 대응하면 표준 동형
$\widetilde{\mathcal{M}}_T = h^*\widetilde{\mathcal{M}}$이 존재한다. 여기서
$\widetilde{\mathcal{M}}_T$는 (\ref{constructions-equation-widetilde-M})로 정의한다.
\item (1)의 동형들은 당겨올림과 호환된다.
\item 표준 사상
$$
\pi^*\mathcal{M}_0 \longrightarrow \widetilde{\mathcal{M}}.
$$
이 존재한다.
\item 구성 $\mathcal{M} \mapsto \widetilde{\mathcal{M}}$은
$\mathcal{M}$에 대해 함자적이다.
\item 구성 $\mathcal{M} \mapsto \widetilde{\mathcal{M}}$은 완전하다.
\item Lemma \ref{constructions-lemma-widetilde-tensor}에서와 같은 표준 사상
$$
\widetilde{\mathcal{M}}
\otimes_{\mathcal{O}_{\underline{\text{Proj}}_S(\mathcal{A})}}
\widetilde{\mathcal{N}}
\longrightarrow
\widetilde{\mathcal{M} \otimes_\mathcal{A} \mathcal{N}}
$$
이 존재한다.
\item (\ref{constructions-equation-global-sections-more-generally})를 일반화하는 표준 사상
$$
\pi^*\mathcal{M}
\longrightarrow
\bigoplus\nolimits_{n \in \mathbf{Z}}
\widetilde{\mathcal{M}(n)}
$$
이 존재한다.
\item $\widetilde{\mathcal{M}}$의 구성은 밑변환과 교환한다.
\end{enumerate}
\end{lemma}

\begin{proof}
생략한다. 이 보조정리는 여러 부분으로 나누고 각 부분을 따로 증명해야 한다.
\end{proof}

\section{상대 Proj의 함자성}
\label{constructions-section-functoriality-relative-proj}

\noindent
이 절은 상대 Proj에 대한 Section \ref{constructions-section-functoriality-proj}의
유사물이다. $S$를 스킴이라 하자. 등급 $\mathcal{O}_S$-대수의 사상
$\psi : \mathcal{A} \to \mathcal{B}$는 이에 결부된 상대 Proj 사이의 사상을
항상 주지는 않는다. 올바른 결과는 다음과 같다.

\begin{lemma}
\label{constructions-lemma-morphism-relative-proj}
$S$를 스킴이라 하자. $\mathcal{A}$, $\mathcal{B}$를 두 준연접 등급
$\mathcal{O}_S$-대수라 하자.
$p : X = \underline{\text{Proj}}_S(\mathcal{A}) \to S$ 및
$q : Y = \underline{\text{Proj}}_S(\mathcal{B}) \to S$로 두자.
$\psi : \mathcal{A} \to \mathcal{B}$를 등급 $\mathcal{O}_S$-대수의
준동형이라 하자. 표준 열린집합 $U(\psi) \subset Y$, $S$ 위의 표준 스킴 사상
$$
r_\psi :
U(\psi)
\longrightarrow
X
$$
및 $\mathbf{Z}$-등급 $\mathcal{O}_{U(\psi)}$-대수의 사상
$$
\theta = \theta_\psi :
r_\psi^*\left(
\bigoplus\nolimits_{d \in \mathbf{Z}} \mathcal{O}_X(d)
\right)
\longrightarrow
\bigoplus\nolimits_{d \in \mathbf{Z}} \mathcal{O}_{U(\psi)}(d).
$$
이 존재한다. 삼중항 $(U(\psi), r_\psi, \theta)$는 다음 성질로
특징지어진다. 모든 아핀 열린집합 $W \subset S$에 대해 삼중항
$$
(U(\psi) \cap p^{-1}W,\quad
r_\psi|_{U(\psi) \cap p^{-1}W} : U(\psi) \cap p^{-1}W \to q^{-1}W,\quad
\theta|_{U(\psi) \cap p^{-1}W})
$$
은 $\psi : \mathcal{A}(W) \to \mathcal{B}(W)$에 결부된 Lemma
\ref{constructions-lemma-morphism-proj}의 삼중항과 같다. 여기서는 식별
$p^{-1}W = \text{Proj}(\mathcal{A}(W))$ 및
$q^{-1}W = \text{Proj}(\mathcal{B}(W))$을 사용한다.
Section \ref{constructions-section-relative-proj-via-glueing}을 보라.
\end{lemma}

\begin{proof}
국소 삼중항들을 붙이면 보조정리가 저절로 증명된다.
\end{proof}

\begin{lemma}
\label{constructions-lemma-morphism-relative-proj-transitive}
$S$를 스킴이라 하자. $\mathcal{A}$, $\mathcal{B}$ 및 $\mathcal{C}$를
준연접 등급 $\mathcal{O}_S$-대수라 하자.
$X = \underline{\text{Proj}}_S(\mathcal{A})$,
$Y = \underline{\text{Proj}}_S(\mathcal{B})$ 및
$Z = \underline{\text{Proj}}_S(\mathcal{C})$로 두자.
$\varphi : \mathcal{A} \to \mathcal{B}$,
$\psi : \mathcal{B} \to \mathcal{C}$를 등급 $\mathcal{O}_S$-대수의
사상이라 하자. 그러면
$$
U(\psi \circ \varphi) = r_\varphi^{-1}(U(\psi))
\quad
\text{and}
\quad
r_{\psi \circ \varphi}
=
r_\varphi \circ r_\psi|_{U(\psi \circ \varphi)}.
$$
이다. 또한 자명한 표기 아래
$$
\theta_\psi \circ r_\psi^*\theta_\varphi
=
\theta_{\psi \circ \varphi}
$$
이다.
\end{lemma}

\begin{proof}
생략한다.
\end{proof}

\begin{lemma}
\label{constructions-lemma-surjective-graded-rings-map-relative-proj}
위의 Lemma \ref{constructions-lemma-morphism-relative-proj}와 같은 가정과 표기를 사용하자.
$\mathcal{A}_d \to \mathcal{B}_d$가 $d \gg 0$에 대해 전사라고 가정하자.
그러면
\begin{enumerate}
\item $U(\psi) = Y$이고,
\item $r_\psi : Y \to X$는 닫힌 몰입이며,
\item 사상 $\theta : r_\psi^*\mathcal{O}_X(n) \to \mathcal{O}_Y(n)$은
전사이지만 일반적으로 동형은 아니다($\mathcal{A} \to \mathcal{B}$가
전사이더라도 마찬가지이다).
\end{enumerate}
\end{lemma}

\begin{proof}
Lemma \ref{constructions-lemma-morphism-relative-proj}와 Lemma
\ref{constructions-lemma-surjective-graded-rings-map-proj}를 결합하면 따른다.
\end{proof}

\begin{lemma}
\label{constructions-lemma-eventual-iso-graded-rings-map-relative-proj}
위의 Lemma \ref{constructions-lemma-morphism-relative-proj}와 같은 가정과 표기를 사용하자.
$\mathcal{A}_d \to \mathcal{B}_d$가 모든 $d \gg 0$에 대해 동형이라고
가정하자. 그러면
\begin{enumerate}
\item $U(\psi) = Y$이고,
\item $r_\psi : Y \to X$는 동형이며,
\item 사상 $\theta : r_\psi^*\mathcal{O}_X(n) \to \mathcal{O}_Y(n)$은
동형이다.
\end{enumerate}
\end{lemma}

\begin{proof}
Lemma \ref{constructions-lemma-morphism-relative-proj}와 Lemma
\ref{constructions-lemma-eventual-iso-graded-rings-map-proj}를 결합하면 따른다.
\end{proof}

\begin{lemma}
\label{constructions-lemma-surjective-generated-degree-1-map-relative-proj}
위의 Lemma \ref{constructions-lemma-morphism-relative-proj}와 같은 가정과 표기를 사용하자.
$\mathcal{A}_d \to \mathcal{B}_d$가 $d \gg 0$에 대해 전사이고
$\mathcal{A}$가 $\mathcal{A}_1$에 의해 $\mathcal{A}_0$ 위에서 생성된다고
가정하자. 그러면
\begin{enumerate}
\item $U(\psi) = Y$이고,
\item $r_\psi : Y \to X$는 닫힌 몰입이며,
\item 사상 $\theta : r_\psi^*\mathcal{O}_X(n) \to \mathcal{O}_Y(n)$은
동형이다.
\end{enumerate}
\end{lemma}

\begin{proof}
Lemma \ref{constructions-lemma-morphism-relative-proj}와 Lemma
\ref{constructions-lemma-surjective-graded-rings-generated-degree-1-map-proj}를 결합하면
따른다.
\end{proof}

\section{가역층과 상대 Proj로 가는 사상}
\label{constructions-section-invertible-relative-proj}

\noindent
다음 보조정리가 어딘가에서 필요할 듯하다. 상황은 다음과 같다.
\begin{enumerate}
\item $S$를 스킴이라 하자.
\item $\mathcal{A}$를 준연접 등급 $\mathcal{O}_S$-대수라 하자.
\item $\pi : \underline{\text{Proj}}_S(\mathcal{A}) \to S$로 $S$ 위의 상대
동차 스펙트럼을 나타내자.
\item $f : X \to S$를 스킴의 사상이라 하자.
\item $\mathcal{L}$을 가역 $\mathcal{O}_X$-가군이라 하자.
\item $\psi : f^*\mathcal{A} \to
\bigoplus_{d \geq 0} \mathcal{L}^{\otimes d}$를 등급
$\mathcal{O}_X$-대수의 준동형이라 하자.
\end{enumerate}
이 자료가 주어지면
$$
U(\psi) = \bigcup\nolimits_{(U, V, a)} U_{\psi(a)}
$$
로 둔다. 여기서 $(U, V, a)$는 다음을 만족한다.
\begin{enumerate}
\item $V \subset S$는 아핀 열린집합이고,
\item $U = f^{-1}(V)$이며,
\item $a \in \mathcal{A}(V)_{+}$는 동차이다.
\end{enumerate}
실제로 그러면 $\psi(a) \in \Gamma(U, \mathcal{L}^{\otimes \deg(a)})$이고
$U_{\psi(a)}$는 이에 대응하는 열린집합이다(Modules, Lemma
\ref{modules-lemma-s-open} 참조).

\begin{lemma}
\label{constructions-lemma-invertible-map-into-relative-proj}
위와 같은 가정과 표기를 사용하자. 사상 $\psi$는 $S$ 위의 표준 스킴 사상
$$
r_{\mathcal{L}, \psi} :
U(\psi) \longrightarrow \underline{\text{Proj}}_S(\mathcal{A})
$$
과 등급 $\mathcal{O}_{U(\psi)}$-대수의 사상
$$
\theta :
r_{\mathcal{L}, \psi}^*\left(
\bigoplus\nolimits_{d \geq 0}
\mathcal{O}_{\underline{\text{Proj}}_S(\mathcal{A})}(d)
\right)
\longrightarrow
\bigoplus\nolimits_{d \geq 0} \mathcal{L}^{\otimes d}|_{U(\psi)}
$$
을 유도한다. 이들은 다음 성질들로 특징지어진다.
\begin{enumerate}
\item 모든 열린집합 $V \subset S$와 모든 $d \geq 0$에 대해 도식
$$
\xymatrix{
\mathcal{A}_d(V) \ar[d]_{\psi} \ar[r]_{\psi} &
\Gamma(f^{-1}(V), \mathcal{L}^{\otimes d}) \ar[d]^{restrict} \\
\Gamma(\pi^{-1}(V),
\mathcal{O}_{\underline{\text{Proj}}_S(\mathcal{A})}(d)) \ar[r]^{\theta} &
\Gamma(f^{-1}(V) \cap U(\psi), \mathcal{L}^{\otimes d})
}
$$
은 가환한다.
\item 임의의 $d \geq 1$과 임의의 열린 부분스킴 $W \subset X$에 대해
$\psi|_W : f^*\mathcal{A}_d|_W \to \mathcal{L}^{\otimes d}|_W$가 전사이면,
사상 $r_{\mathcal{L}, \psi}$의 제한은 상대 동차 스펙트럼의 구성으로부터
존재하는 사상
$W \to \underline{\text{Proj}}_S(\mathcal{A})$와 일치한다. Definition
\ref{constructions-definition-relative-proj}를 보라.
\item 임의의 아핀 열린집합 $V \subset S$에 대해 제한
$$
(U(\psi) \cap f^{-1}(V), r_{\mathcal{L}, \psi}|_{U(\psi) \cap f^{-1}(V)},
\theta|_{U(\psi) \cap f^{-1}(V)})
$$
은 $i_V$를 통한 식별 아래(Lemma \ref{constructions-lemma-glue-relative-proj} 참조)
Lemma \ref{constructions-lemma-invertible-map-into-proj}의 삼중항
$(U(\psi'), r_{\mathcal{L}, \psi'}, \theta')$과 일치한다. 이는 사상
$\psi' : A = \mathcal{A}(V) \to \Gamma_*(f^{-1}(V), \mathcal{L}|_{f^{-1}(V)})$에
결부된 삼중항이며, 이 사상은 $\psi$가 유도한다.
\end{enumerate}
\end{lemma}

\begin{proof}
특징화 (3)을 사용하여 사상 $r_{\mathcal{L}, \psi}$와 $\theta$를 $S$ 위에서
국소적으로 구성한다. Lemma \ref{constructions-lemma-invertible-map-into-proj}의 유일성을
사용하여 구성이 붙음을 보인다. 세부사항은 생략한다.
\end{proof}

\section{가역층에 의한 꼬임과 상대 Proj}
\label{constructions-section-twisting-and-proj}

\noindent
$S$를 스킴이라 하자.
$\mathcal{A} = \bigoplus_{d \geq 0} \mathcal{A}_d$를 준연접 등급
$\mathcal{O}_S$-대수라 하자. $\mathcal{L}$을 $S$ 위의 가역층이라 하자.
이 상황에서 또 다른 준연접 등급 $\mathcal{O}_S$-대수, 즉
$$
\mathcal{B}
=
\bigoplus\nolimits_{d \geq 0}
\mathcal{A}_d \otimes_{\mathcal{O}_S} \mathcal{L}^{\otimes d}
$$
를 얻는다. $\mathcal{A}$와 $\mathcal{B}$의 상대 동차 스펙트럼은
동형임이 밝혀진다.

\begin{lemma}
\label{constructions-lemma-twisting-and-proj}
위와 같은 표기 $S$, $\mathcal{A}$, $\mathcal{L}$ 및
$\mathcal{B}$를 사용하자. 다음 표준 동형이 존재한다.
$$
\xymatrix{
P = \underline{\text{Proj}}_S(\mathcal{A})
\ar[rr]_g \ar[rd]_\pi & &
\underline{\text{Proj}}_S(\mathcal{B}) = P'
\ar[ld]^{\pi'} \\
& S &
}
$$
이 동형은 다음 성질들을 갖는다.
\begin{enumerate}
\item 동형
$\theta_n : g^*\mathcal{O}_{P'}(n)
\to
\mathcal{O}_P(n) \otimes \pi^*\mathcal{L}^{\otimes n}$
들이 존재하고, 이들이 서로 맞물려 $\mathbf{Z}$-등급 대수의 동형
$$
\theta :
g^*\left(
\bigoplus\nolimits_{n \in \mathbf{Z}} \mathcal{O}_{P'}(n)
\right)
\longrightarrow
\bigoplus\nolimits_{n \in \mathbf{Z}} \mathcal{O}_P(n)
\otimes \pi^*\mathcal{L}^{\otimes n}
$$
을 준다.
\item 모든 열린집합 $V \subset S$에 대해 도식
$$
\xymatrix{
\mathcal{A}_n(V) \otimes \mathcal{L}^{\otimes n}(V)
\ar[r]_{multiply} \ar[d]^{\psi \otimes \pi^*}
&
\mathcal{B}_n(V) \ar[dd]^\psi \\
\Gamma(\pi^{-1}V, \mathcal{O}_P(n)) \otimes
\Gamma(\pi^{-1}V, \pi^*\mathcal{L}^{\otimes n})
\ar[d]^{multiply} \\
\Gamma(\pi^{-1}V, \mathcal{O}_P(n) \otimes \pi^*\mathcal{L}^{\otimes n})
&
\Gamma(\pi'^{-1}V, \mathcal{O}_{P'}(n)) \ar[l]_-{\theta_n}
}
$$
은 가환한다.
\item 필요에 따라 여기에 더 추가한다.
\end{enumerate}
\end{lemma}

\begin{proof}
$\mathcal{L} \cong \mathcal{O}_S$일 때 이는 항등사상이다. 일반적으로
$S$의 열린 덮개를 택하여 각 조각 위에서 $\mathcal{L}$을 자명화하고,
이에 대응하는 사상들을 붙인다. 세부사항은 생략한다.
\end{proof}

\section{사영 다발}
\label{constructions-section-projective-bundle}

\noindent
$S$를 스킴이라 하자. $\mathcal{E}$를 $\mathcal{O}_S$-가군의 준연접층이라
하자. Modules, Lemma \ref{modules-lemma-whole-tensor-algebra-permanence}에
의해 대칭대수 $\text{Sym}(\mathcal{E})$, 즉 $\mathcal{E}$의
$\mathcal{O}_S$ 위 대칭대수는 $\mathcal{O}_S$-대수의 준연접층이다.
차수 $1$의 원소들로 $\mathcal{O}_S$ 위에서 생성된다는 점에 유의하자.
따라서 앞 절의 구성, 구체적으로 Lemmas \ref{constructions-lemma-relative-proj}와
\ref{constructions-lemma-apply-relative}를 여기에 적용하는 것이 의미가 있다.

\begin{definition}
\label{constructions-definition-projective-bundle}
$S$를 스킴이라 하자. $\mathcal{E}$를 준연접
$\mathcal{O}_S$-가군이라 하자.\footnote{독자는 여기서
$\mathcal{E}$가 유한 국소 자유라는 조건을 기대할 수도 있다. 우리는
\cite[II, Definition 4.1.1]{EGA}와 일관되도록 이 조건을 부과하지 않는다.}
다음과 같이 쓰고
$$
\pi :
\mathbf{P}(\mathcal{E}) = \underline{\text{Proj}}_S(\text{Sym}(\mathcal{E}))
\longrightarrow
S
$$
이를 {\it $\mathcal{E}$에 결부된 사영 다발}이라고 한다. 기호
$\mathcal{O}_{\mathbf{P}(\mathcal{E})}(n)$은 Lemma
\ref{constructions-lemma-apply-relative}의 가역
$\mathcal{O}_{\mathbf{P}(\mathcal{E})}$-가군을 나타내고, $n$번째
{\it 구조층의 꼬임}이라고 한다.
\end{definition}

\noindent
Lemma \ref{constructions-lemma-glue-relative-proj-twists}에 따르면 표준
$\mathcal{O}_S$-가군 준동형
$$
\text{Sym}^n(\mathcal{E})
\longrightarrow
\pi_*\mathcal{O}_{\mathbf{P}(\mathcal{E})}(n)
\quad\text{equivalently}\quad
\pi^*\text{Sym}^n(\mathcal{E})
\longrightarrow
\mathcal{O}_{\mathbf{P}(\mathcal{E})}(n)
$$
이 모든 $n \geq 0$에 대해 존재한다. 특히 $n = 1$일 때
$$
\mathcal{E}
\longrightarrow
\pi_*\mathcal{O}_{\mathbf{P}(\mathcal{E})}(1)
\quad\text{equivalently}\quad
\pi^*\mathcal{E}
\longrightarrow
\mathcal{O}_{\mathbf{P}(\mathcal{E})}(1)
$$
이고, 사상
$\pi^*\mathcal{E} \to \mathcal{O}_{\mathbf{P}(\mathcal{E})}(1)$은 Lemma
\ref{constructions-lemma-apply-relative}에 의해 전사이다. 이는
$\mathbf{P}(\mathcal{E})$의 구성을 어떻게 정규화했는지를 기억하는 좋은
방법이다.

\medskip\noindent
주의: 어떤 문헌에서는 스킴 $\mathbf{P}(\mathcal{E})$를 $\mathcal{E}$가
$S$ 위에서 유한 국소 자유일 때에만 정의한다. 또한 어떤 경우에는
$\mathbf{P}(\mathcal{E})$를 실제로 우리의
$\mathbf{P}(\mathcal{E}^\vee)$로 정의하는데, 여기서
$\mathcal{E}^\vee$는 $\mathcal{E}$의 쌍대이다(그리고 이 정의는
$\mathcal{E}$가 유한 국소 자유일 때에만 사용한다).

\medskip\noindent
$S$, $\mathcal{E}$, $\mathbf{P}(\mathcal{E}) \to S$가 Definition
\ref{constructions-definition-projective-bundle}에서와 같다고 하자.
$f : T \to S$를 $S$ 위의 스킴이라 하자.
$\psi : f^*\mathcal{E} \to \mathcal{L}$을 전사라 하되,
$\mathcal{L}$은 가역 $\mathcal{O}_T$-가군이라고 하자. 유도된 등급
$\mathcal{O}_T$-대수 사상
$$
f^*\text{Sym}(\mathcal{E}) = \text{Sym}(f^*\mathcal{E}) \to
\text{Sym}(\mathcal{L}) = \bigoplus\nolimits_{n \geq 0} \mathcal{L}^{\otimes n}
$$
은 사상
$$
\varphi_{\mathcal{L}, \psi} : T \longrightarrow \mathbf{P}(\mathcal{E})
$$
에 대응한다. 이는 $S$ 위의 사상이며, Section
\ref{constructions-section-relative-proj}에서 상대 Proj를 함자 $F$를 표현하는 스킴으로
구성했기 때문에 생긴다. 한편 사상
$\varphi : T \to \mathbf{P}(\mathcal{E})$가 $S$ 위에 주어지면
$\mathcal{L} = \varphi^*\mathcal{O}_{\mathbf{P}(\mathcal{E})}(1)$로 두고,
$\psi : f^*\mathcal{E} \to \mathcal{L}$을 $\varphi$에 의해 표준 전사
$\pi^*\mathcal{E} \to \mathcal{O}_{\mathbf{P}(\mathcal{E})}(1)$를 당겨온
사상으로 둘 수 있다. Lemma \ref{constructions-lemma-apply-relative}에 의해 이 두 구성은
쌍 $(\mathcal{L}, \psi)$의 동형류들의 집합과 사상
$\varphi : T \to \mathbf{P}(\mathcal{E})$들의 집합 사이의 서로 역인
전단사이며, 뒤의 사상들은 $S$ 위에서 취한다. 따라서
$\mathbf{P}(\mathcal{E})$는 $f : T \to S$에
$\mathcal{O}_T$-가군 몫 중 $f^*\mathcal{E}$의 몫이고 국소 자유 계수가
$1$인 것들의 집합을 결부시키는 함자를 표현한다.

\begin{example}[벡터공간의 사영공간]
\label{constructions-example-projective-space}
$k$를 체라 하자. $V$를 $k$-벡터공간이라 하자. 이에 대응하는
{\it 사영공간}은 $k$-스킴
$$
\mathbf{P}(V) = \text{Proj}(\text{Sym}(V))
$$
이다. 여기서 $\text{Sym}(V)$는 $V$의 $k$ 위 대칭대수이다. 물론
$\mathbf{P}(V) \cong \mathbf{P}^n_k$이고, 이는 $\dim(V) = n + 1$일 때
성립한다. 실제로
이때 $V$ 위의 대칭대수는 $n + 1$개 변수를 갖는 다항식환과 동형이다.
$V$를 $\Spec(k)$ 위의 준연접 가군으로 생각하면 $\mathbf{P}(V)$는
$\Spec(k)$ 위의 이에 대응하는 사영공간 다발이다. 위의 논의에 의해
$k$-값점 $p$가 $\mathbf{P}(V)$에 속한다는 것은 $k$-벡터공간의 전사
$V \to L_p$가 주어지고 $\dim(L_p) = 1$인 것과 같다. 더 일반적으로
$X$를 $k$ 위의 스킴,
$\mathcal{L}$을 가역 $\mathcal{O}_X$-가군이라 하고,
$\psi : V \to \Gamma(X, \mathcal{L})$를 $k$-선형 사상이라 하자.
$\mathcal{L}$이 $\mathcal{O}_X$-가군으로서 $\psi$의 상에 있는 단면들에
의해 생성된다고 가정한다. 그러면 위의 논의는 $k$ 위의
표준 스킴 사상
$$
\varphi_{\mathcal{L}, \psi} : X \longrightarrow \mathbf{P}(V)
$$
을 주며, 동형
$\theta : \varphi_{\mathcal{L}, \psi}^*\mathcal{O}_{\mathbf{P}(V)}(1)
\to \mathcal{L}$이 존재하고 $\psi$는 합성
$$
V \to
\Gamma(\mathbf{P}(V), \mathcal{O}_{\mathbf{P}(V)}(1))
\to
\Gamma(X, \varphi_{\mathcal{L}, \psi}^*\mathcal{O}_{\mathbf{P}(V)}(1))
\to
\Gamma(X, \mathcal{L})
$$
과 일치한다. Lemma \ref{constructions-lemma-invertible-map-into-proj}를 보라.
$V \subset \Gamma(X, \mathcal{L})$가 부분공간이면 위에서 구성한 사상을
단순히 $\varphi_{\mathcal{L}, V}$로 나타낼 것이다.
$\dim(V) = n + 1$이고 기저 $v_0, \ldots, v_n$을 $V$에서 택하면 도식
$$
\xymatrix{
X \ar@{=}[d] \ar[rr]_{\varphi_{\mathcal{L}, \psi}} & &
\mathbf{P}(V) \ar[d]^{\cong} \\
X \ar[rr]^{\varphi_{(\mathcal{L}, (s_0, \ldots, s_n))}} & &
\mathbf{P}^n_k
}
$$
은 가환한다. 여기서
$s_i = \psi(v_i) \in \Gamma(X, \mathcal{L})$이고,
$\varphi_{(\mathcal{L}, (s_0, \ldots, s_n))}$는 Section
\ref{constructions-section-projective-space}에서와 같으며, 오른쪽 세로 화살표는 동형
$k[T_0, \ldots, T_n] \to \text{Sym}(V)$에 대응한다. 이 동형은
$T_i$를 $v_i$로 보낸다.
\end{example}

\begin{example}
\label{constructions-example-projective-bundle}
사상
$\text{Sym}^n(\mathcal{E}) \to
\pi_*(\mathcal{O}_{\mathbf{P}(\mathcal{E})}(n))$은 $\mathcal{E}$가
국소 자유이면 동형이지만, 일반적으로는 동형일 필요가 없다. 실제로 이 사상이
$n = 1$일 때 단사가 아닌 예를 제시하겠다. $S = \Spec(A)$로 두되
$$
A = k[u, v, s_1, s_2, t_1, t_2]/I
$$
라 하자. 여기서 $k$는 체이고
$$
I = (-us_1 + vt_1 + ut_2, vs_1 + us_2 - vt_2, vs_2, ut_1).
$$
$\overline{u}$로 $u$의 $A$에서의 류를 나타내고 다른 변수들도 같은 방식으로
나타내자.
$M = (Ax \oplus Ay)/A(\overline{u}x + \overline{v}y)$로 두면
$$
\text{Sym}(M) = A[x, y]/(\overline{u}x + \overline{v}y)
= k[x, y, u, v, s_1, s_2, t_1, t_2]/J
$$
이고, 여기서
$$
J = (-us_1 + vt_1 + ut_2, vs_1 + us_2 - vt_2, vs_2, ut_1, ux + vy).
$$
이다. 이 경우 준연접층 $\mathcal{E} = \widetilde{M}$이
$S = \Spec(A)$ 위에 주어질 때 이에 결부된 사영 다발은 스킴
$$
P =
\text{Proj}(\text{Sym}(M)).
$$
이다. 이 스킴은 아핀 열린 덮개 $P = D_{+}(x) \cup D_{+}(y)$를 갖는다.
원소 $m \in M$을 생각하되, 이를 원소 $us_1x + vt_2y$의 상으로 택하자.
다음에 유의하자.
$$
x(us_1x + vt_2y) = (s_1x + s_2y)(ux + vy) \bmod I
$$
그리고
$$
y(us_1x + vt_2y) = (t_1x + t_2y)(ux + vy) \bmod I.
$$
첫 번째 등식에 의해 $m$은 $\mathcal{O}_P(1)$의 단면으로서
$D_{+}(x)$ 위에서 영으로 사상되고, 두 번째 등식에 의해
$\mathcal{O}_P(1)$의 단면으로서 $D_{+}(y)$ 위에서 영으로 사상된다. 따라서
$m$은 $\Gamma(P, \mathcal{O}_P(1))$에서 영으로 사상된다. 한편
$m \not = 0$임을 주장한다. 그러면 $m$은 $\mathcal{E}$의 영이 아닌
대역단면이 $\Gamma(P, \mathcal{O}_P(1))$에서 영으로 사상되는 예를 준다.
모순을 얻기 위해 $m = 0$이라고 가정하자. 이 경우 어떤 원소
$f \in k[u, v, s_1, s_2, t_1, t_2]$가 존재하여
$$
us_1x + vt_2y = f(ux + vy) \bmod I
$$
이다. $I$가 차수 $2$인 동차 다항식들로 생성되므로 $f$를 동차 성분들로
분해하여 차수 1인 성분을 취할 수 있다. 다시 말해 어떤
$a, b, \alpha_1, \alpha_2, \beta_1, \beta_2 \in k$에 대해
$$
f = au + bv + \alpha_1s_1 + \alpha_2s_2 + \beta_1t_1 + \beta_2t_2
$$
라고 가정할 수 있다. 그 결과 얻는 조건들은
$$
\begin{matrix}
us_1 - u(au + bv + \alpha_1s_1 + \alpha_2s_2 + \beta_1t_1 + \beta_2t_2)
\in I \\
vt_2 - v(au + bv + \alpha_1s_1 + \alpha_2s_2 + \beta_1t_1 + \beta_2t_2)
\in I
\end{matrix}
$$
이다. 항 $u^2, uv, v^2$는 $I$의 생성자에 없으므로
$a = b = 0$임을 알 수 있다. 따라서 관계식
$$
\begin{matrix}
us_1 - u(\alpha_1s_1 + \alpha_2s_2 + \beta_1t_1 + \beta_2t_2)
\in I \\
vt_2 - v(\alpha_1s_1 + \alpha_2s_2 + \beta_1t_1 + \beta_2t_2)
\in I
\end{matrix}
$$
을 얻는다. $I$의 첫째 생성자를 사용하여 $us_1$이 나타날 때마다
$vt_1 + ut_2$로 바꾸고, $I$의 둘째 생성자를 사용하여 $vs_1$이 나타날
때마다 $-us_2 + vt_2$로 바꾸며, 셋째 생성자를 사용하여 $vs_2$가 나타나는
항을 제거하고 셋째 생성자를 사용하여 $ut_1$이 나타나는 항을 제거할 수
있다. 그러면 관계식
$$
\begin{matrix}
(1 - \alpha_1)vt_1 + (1 - \alpha_1)ut_2 - \alpha_2us_2 - \beta_2ut_2 = 0 \\
(1 - \alpha_1)vt_2 + \alpha_1us_2 - \beta_1vt_1 - \beta_2vt_2 = 0
\end{matrix}
$$
을 얻는다. 이는 $\alpha_1$이 $0$과 $1$이어야 함을 동시에 뜻하므로,
원하는 모순을 얻는다.
\end{example}


\begin{lemma}
\label{constructions-lemma-projective-bundle-separated}
$S$를 스킴이라 하자. 사영 다발의 구조사상
$\mathbf{P}(\mathcal{E}) \to S$는 $S$ 위에서 분리사상이다.
\end{lemma}

\begin{proof}
Lemma \ref{constructions-lemma-relative-proj-separated}에서 곧바로 따른다.
\end{proof}

\begin{lemma}
\label{constructions-lemma-projective-space-bundle}
$S$를 스킴이라 하자. $n \geq 0$이라 하자. 그러면
$\mathbf{P}^n_S$는 $S$ 위의 사영 다발이다.
\end{lemma}

\begin{proof}
다음에 유의하자.
$$
\mathbf{P}^n_{\mathbf{Z}} =
\text{Proj}(\mathbf{Z}[T_0, \ldots, T_n]) =
\underline{\text{Proj}}_{\Spec(\mathbf{Z})}
\left(\widetilde{\mathbf{Z}[T_0, \ldots, T_n]}\right)
$$
여기서 환 $\mathbf{Z}[T_0, \ldots, T_n]$의 등급은
$\deg(T_i) = 1$로 주어지고 $\mathbf{Z}$의 원소들은 차수 $0$에 있다.
$\mathbf{P}^n_S$가
$\mathbf{P}^n_{\mathbf{Z}} \times_{\Spec(\mathbf{Z})} S$로 정의됨을
상기하자. 또한 상대 동차 스펙트럼의 구성은 밑변환과 교환한다. Lemma
\ref{constructions-lemma-relative-proj-base-change}를 보라. 임의의 스킴
$g : S \to \Spec(\mathbf{Z})$에 대해
$g^*\mathcal{O}_{\Spec(\mathbf{Z})}[T_0, \ldots, T_n]
= \mathcal{O}_S[T_0, \ldots, T_n]$이다. 위의 사실들을 결합하면
$$
\mathbf{P}^n_S = \underline{\text{Proj}}_S(\mathcal{O}_S[T_0, \ldots, T_n]).
$$
를 얻는다. 마지막으로
$\mathcal{O}_S[T_0, \ldots, T_n] = \text{Sym}(\mathcal{O}_S^{\oplus n + 1})$임에
유의하자. 따라서
$\mathbf{P}^n_S$는 $S$ 위의 사영 다발이다.
\end{proof}

\section{그라스만 다양체}
\label{constructions-section-grassmannian}

\noindent
이 절에서는 표준 그라스만 다양체 함자들을 도입하고 이들이 스킴으로
표현됨을 보인다. 정수 $k$, $n$을 $0 < k < n$이 되도록 택하자. 대략적으로
$k$-차원 부분공간들을 $n$-공간 안에서 매개화하는 함자
\begin{equation}
\label{constructions-equation-gkn}
G(k, n) : \Sch \longrightarrow \textit{Sets}
\end{equation}
를 구성할 것이다. 하지만 기술적인 이유로 $(n - k)$-차원 몫들을
매개화하는 편이 더 편리하므로 그렇게 하겠다.

\medskip\noindent
더 정확히 말해 $G(k, n)$은 스킴 $S$에 전사들의 동형류들의 집합
$G(k, n)(S)$를 결부시킨다. 여기서 전사는
$$
q : \mathcal{O}_S^{\oplus n} \longrightarrow \mathcal{Q}
$$
꼴이고 $\mathcal{Q}$는 유한 국소 자유 $\mathcal{O}_S$-가군이며 그 계수는
$n - k$이다. 이것은 실제로 집합임에 유의하자. 예를 들어
Modules, Lemma \ref{modules-lemma-set-isomorphism-classes-finite-type-modules}를
사용하거나, 전사 $q$의 동형류가 $q$의 핵으로 결정된다는 사실을 사용할 수
있다(그리고 층 하나가 주어지면 그 부분층들은 집합을 이룬다). 스킴의 사상
$f : T \to S$가 주어지면
$G(k, n)(f) : G(k, n)(S) \to G(k, n)(T)$는
$q : \mathcal{O}_S^{\oplus n} \longrightarrow \mathcal{Q}$의 동형류를
$f^*q : \mathcal{O}_T^{\oplus n} \longrightarrow f^*\mathcal{Q}$의 동형류로
보낸다. 이는 다음 이유들로 의미가 있다. (1)
$f^*\mathcal{O}_S = \mathcal{O}_T$이고, (2) $f^*$는 가법적이며, (3)
$f^*$는 국소 자유 가군을 보존하고(Modules, Lemma
\ref{modules-lemma-pullback-locally-free}), (4) $f^*$는 오른쪽 완전하다
(Modules, Lemma \ref{modules-lemma-exactness-pushforward-pullback}).

\begin{lemma}
\label{constructions-lemma-gkn-representable}
$0 < k < n$이라 하자. (\ref{constructions-equation-gkn})의 함자 $G(k, n)$은 스킴으로
표현 가능하다.
\end{lemma}

\begin{proof}
$F = G(k, n)$로 두자. 보조정리를 증명하기 위해 Schemes, Lemma
\ref{schemes-lemma-glue-functors}의 판정법을 사용하겠다. $F$가 자리스키
위상에 대한 층 성질을 만족하는 이유는 층들을 붙일 수 있기 때문이다.
Sheaves, Section \ref{sheaves-section-glueing-sheaves}를 보라(일부
세부사항은 생략한다).

\medskip\noindent
부분함자들의 족 $F_i$. $I$를 $\{1, \ldots, n\}$의 원소 수 $n - k$인
부분집합들의 집합이라 하자. 스킴 $S$와
$j \in \{1, \ldots, n\}$가 주어지면 $e_j$로 대역단면
$$
e_j = (0, \ldots, 0, 1, 0, \ldots, 0)\quad(1\text{ in }j\text{th spot})
$$
을 나타낸다. 이는 $\mathcal{O}_S^{\oplus n}$의 대역단면이다. 물론 이
단면들은 $\mathcal{O}_S^{\oplus n}$을 자유롭게 생성한다. 마찬가지로
$j \in \{1, \ldots, n - k\}$에 대해 $f_j$로
$\mathcal{O}_S^{\oplus n - k}$의 대역단면 중 $j$번째 성분만 $1$이고
나머지 성분은 모두 영인 것을 나타낸다. $i \in I$에 대해
$$
s_i : \mathcal{O}_S^{\oplus n - k} \longrightarrow \mathcal{O}_S^{\oplus n}
$$
로 두되, 이는 쌍대사영
$\mathcal{O}_S \to \mathcal{O}_S^{\oplus n}$들 중 $I$의 원소들에
대응하는 것들의 직합이다. 더 정확히
$i = \{i_1, \ldots, i_{n - k}\}$이고
$i_1 < i_2 < \ldots < i_{n - k}$이면 $s_i$는 $f_j$를 $e_{i_j}$로
보내며, 여기서 $j \in \{1, \ldots, n - k\}$이다. 이 표기를 사용하여
$$
F_i(S) = \{q : \mathcal{O}_S^{\oplus n} \to \mathcal{Q} \in F(S) \mid
q \circ s_i \text{ is surjective}\}
\subset F(S)
$$
로 둘 수 있다. 스킴의 사상 $f : T \to S$가 주어지면 당겨올림
$f^*s_i$는 $T$ 위의 대응하는 사상이다. $f^*$가 오른쪽 완전하므로
(Modules, Lemma \ref{modules-lemma-exactness-pushforward-pullback}),
$F_i$가 $F$의 부분함자라고 결론짓는다.

\medskip\noindent
$F_i$의 표현 가능성. 이를 증명할 때 (번호를 다시 매겨)
$i = \{1, \ldots, n - k\}$라고 가정해도 된다. 이는 $s_i$가 첫
$n - k$개 성분의 포함임을 뜻한다. $q \circ s_i$가 전사이면 같은 계수의
유한 국소 자유 가군 사이의 전사이므로 $q \circ s_i$는 동형이다(Modules,
Lemma \ref{modules-lemma-map-finite-locally-free}). 따라서
$q : \mathcal{O}_S^{\oplus n} \to \mathcal{Q}$가 $F_i(S)$의 원소이면
$q \circ s_i$를 사용하여 $\mathcal{Q}$를
$\mathcal{O}_S^{\oplus n - k}$와 식별할 수 있다. 그렇게 하면
$$
q : \mathcal{O}_S^{\oplus n} \longrightarrow \mathcal{O}_S^{\oplus n - k}
$$
를 얻고, 이는 $e_j$를 $f_j$로 보내며 $j = 1, \ldots, n - k$이다
(표기는 위와 같다). $q$를 완전히 결정하려면
$q(e_{n - k + 1}), \ldots, q(e_n)$의 상들을
$\Gamma(S, \mathcal{O}_S^{\oplus n - k})$에서 정해야 한다. 따라서
$F_i$는 함자
$$
S \longmapsto
\prod\nolimits_{j = n - k + 1, \ldots, n}
\Gamma(S,  \mathcal{O}_S^{\oplus n - k})
$$
와 동형이다. 이 함자는 $k(n - k)$중 자기곱과 동형이며, 여기서 거듭제곱하는
함자는 $S \mapsto \Gamma(S, \mathcal{O}_S)$이다. Schemes, Example
\ref{schemes-example-global-sections}에 의해 뒤의 함자는
$\mathbf{A}^1_\mathbf{Z}$로 표현된다. 따라서 $F_i$는
$\mathbf{A}^{k(n - k)}_\mathbf{Z}$로 표현된다. 실제로
$\Spec(\mathbf{Z})$ 위의 올곱은 스킴의 범주에서의 곱이다.

\medskip\noindent
포함 $F_i \subset F$는 열린 몰입들로 표현된다. $S$를 스킴이라 하고,
$q : \mathcal{O}_S^{\oplus n} \to \mathcal{Q}$를 $F(S)$의 원소라 하자.
Modules, Lemma \ref{modules-lemma-finite-type-surjective-on-stalk}에 의해
집합
$U_i = \{s \in S \mid (q \circ s_i)_s\text{ surjective}\}$는 $S$에서
열려 있다. $\mathcal{O}_{S, s}$는 국소환이고 $\mathcal{Q}_s$는 유한
$\mathcal{O}_{S, s}$-가군이므로, 나카야마 보조정리(Algebra, Lemma
\ref{algebra-lemma-NAK})에 의해
$$
s \in U_i \Leftrightarrow
\left(
\text{the map }
\kappa(s)^{\oplus n - k} \to \mathcal{Q}_s/\mathfrak m_s\mathcal{Q}_s
\text{ induced by }
(q \circ s_i)_s
\text{ is surjective}
\right)
$$
이다. $f : T \to S$를 스킴의 사상이라 하고 $t \in T$를
$s \in S$로 사상되는 점이라 하자. 다음이 성립한다.
$(f^*\mathcal{Q})_t =
\mathcal{Q}_s \otimes_{\mathcal{O}_{S, s}} \mathcal{O}_{T, t}$
(Sheaves, Lemma \ref{sheaves-lemma-stalk-pullback-modules}), 이하도
마찬가지이다. 따라서 사상
$$
\kappa(t)^{\oplus n - k} \to (f^*\mathcal{Q})_t/\mathfrak m_t(f^*\mathcal{Q})_t
$$
은 $(f^*q \circ f^*s_i)_t$에 의해 유도되며, 위의 사상
$\kappa(s)^{\oplus n - k} \to \mathcal{Q}_s/\mathfrak m_s\mathcal{Q}_s$를
체의 확대
$\kappa(t)/\kappa(s)$로 밑변환한 것이다. 따라서 $s \in U_i$인 것과
$t$가 $f^*q$에 대한 대응하는 열린집합에 속하는 것은 동치이다. 특히
$T \to S$가 $U_i$를 통해 분해되는 것과
$f^*q \in F_i(T)$인 것은 동치이며, 이것이 원한 결과이다.

\medskip\noindent
모음 $F_i$, $i \in I$는 $F$를 덮는다.
$q : \mathcal{O}_S^{\oplus n} \to \mathcal{Q}$를 $F(S)$의 원소라 하자.
모든 점 $s$에 대해, 그 점이 놓인 $S$에서 $i \in I$가 존재하여
$s_i$가 $s$의 한 근방에서
전사임을 보여야 한다. 따라서 합성들
$$
\kappa(s)^{\oplus n - k} \xrightarrow{s_i}
\kappa(s)^{\oplus n} \rightarrow
\mathcal{Q}_s/\mathfrak m_s\mathcal{Q}_s
$$
중 하나가 전사임을 보이면 된다(앞 문단을 보라).
$\mathcal{Q}_s/\mathfrak m_s\mathcal{Q}_s$는 차원 $n - k$인
벡터공간이므로 이는 벡터공간 이론에서 따른다.
\end{proof}

\begin{definition}
\label{constructions-definition-grassmannian}
$0 < k < n$이라 하자. 스킴 $\mathbf{G}(k, n)$, 즉 함자
$G(k, n)$을 표현하는 스킴을 {\it $\mathbf{Z}$ 위의 그라스만 다양체}라고
한다. 그 밑변환 $\mathbf{G}(k, n)_S$를 스킴 $S$에 대해
{\it $S$ 위의 그라스만 다양체}라고 한다. $R$가 환이면
$\Spec(R)$로의 밑변환을 $\mathbf{G}(k, n)_R$로 나타내고
{\it $R$ 위의 그라스만 다양체}라고 한다.
\end{definition}

\noindent
Lemma \ref{constructions-lemma-gkn-representable}에서 이 함자들이 실제로 표현 가능함을
보였으므로 정의는 의미가 있다.

\begin{lemma}
\label{constructions-lemma-projective-space-grassmannian}
$n \geq 1$이라 하자. 표준 동형
$\mathbf{G}(n, n + 1) = \mathbf{P}^n_\mathbf{Z}$이 존재한다.
\end{lemma}

\begin{proof}
Lemma \ref{constructions-lemma-projective-space}에 따르면 스킴
$\mathbf{P}^n_\mathbf{Z}$는 스킴 $S$에 다음 쌍들의 동형류들의 집합을
결부시키는 함자를 표현한다. 이 쌍은
$(\mathcal{L}, (s_0, \ldots, s_n))$ 꼴이고, 가역 가군
$\mathcal{L}$과 대역단면들의 $(n + 1)$-튜플로 이루어지며, 이 단면들은
$\mathcal{L}$을 생성한다.
이러한 쌍이 주어지면 몫
$$
\mathcal{O}_S^{\oplus n + 1} \longrightarrow \mathcal{L},\quad
(h_0, \ldots, h_n) \longmapsto \sum h_i s_i.
$$
을 얻는다. 거꾸로
$q : \mathcal{O}_S^{\oplus n + 1} \to \mathcal{Q}$가
$G(n, n + 1)(S)$의 원소로 주어지면 이러한 쌍, 즉
$(\mathcal{Q}, (q(e_1), \ldots, q(e_{n + 1})))$을 얻는다. 여기서
$e_i$, $i = 1, \ldots, n + 1$은 자유 가군
$\mathcal{O}_S^{\oplus n + 1}$의 표준 생성 단면들이다. 이 구성들이
함자들의 서로 역인 변환들을 정의한다는 확인은 생략한다.
\end{proof}
